\documentclass[10pt]{article}
\usepackage{amsmath,amssymb,amsthm,mathtools}
\usepackage{graphicx}
\usepackage[margin=1in]{geometry}
\usepackage{enumitem}
\usepackage[colorlinks=true,linkcolor=blue,citecolor=blue,urlcolor=blue]{hyperref}

\newtheorem{theorem}{Theorem}[section]
\newtheorem{lemma}[theorem]{Lemma}
\newtheorem{proposition}[theorem]{Proposition}
\newtheorem{corollary}[theorem]{Corollary}
\theoremstyle{definition}
\newtheorem{definition}[theorem]{Definition}
\newtheorem{remark}[theorem]{Remark}
\newtheorem{algorithm}[theorem]{Algorithm}

\DeclareMathOperator{\SO}{SO}
\DeclareMathOperator{\Ort}{O}
\DeclareMathOperator{\tr}{tr}
\DeclareMathOperator{\dist}{dist}
\DeclareMathOperator{\diag}{diag}
\DeclareMathOperator{\dbsup}{db-sup}
\newcommand{\R}{\mathbb{R}}
\newcommand{\Z}{\mathbb{Z}}
\newcommand{\Q}{\mathbb{Q}}
\newcommand{\Sph}{\mathbb{S}}
\newcommand{\rr}{\rho}
\newcommand{\Dom}{\mathcal{D}}
\newcommand{\ip}[2]{\langle #1,#2\rangle}

\begin{document}

\title{A computer-assisted proof of Kuperberg's six-cylinder conjecture}
\author{Ivan Mati\'c\thanks{Department of Mathematics, Baruch College,
City University of New York.}
\and Rado\v{s} Radoi\v{c}i\'c\footnotemark[1]}
 
  \section{Term $\ell=1$}
\subsection{Matrix $D_1$}
We will write the matrix $D^{-1}$ instead of $D$. The matrices $D$ and $D^{-1}$ are diagonal, so one is easy to construct from the other. The matrix $D^{-1}$ is more convenient becauce the denominators of diagonal terms are powers of $2$. Therefore, their binary expressions are finite and there is no rounding in computer storage.\[D_1^{-1}=\left[\begin{array}{ccc}
3 & 0 & 0\\
0 & 3 & 0\\
0 & 0 & 3
\end{array}\right].\]
\subsection{Matrix $C_1$}
\[C_1=\left[\begin{array}{ccc}
\frac{789965}{2^{20}} & 0 & 0\\
0 & \frac{175}{2^{19}} & 0\\
0 & 0 & \frac{175}{2^{19}}
\end{array}\right].\]
\subsection{Matrix $T_{h}^1(q_0,q_1,q_2,q_3)$}
\begin{align*}(T_{h}^1)_{1,1}&=1\cdot q_0^{2}-1\cdot q_1^{2}-1\cdot q_2^{2}+1\cdot q_3^{2}\end{align*}
\begin{align*}(T_{h}^1)_{1,2}&=2\cdot q_0^{1}\cdot q_1^{1}+2\cdot q_2^{1}\cdot q_3^{1}\end{align*}
\begin{align*}(T_{h}^1)_{1,3}&=-2\cdot q_0^{1}\cdot q_2^{1}+2\cdot q_1^{1}\cdot q_3^{1}\end{align*}
\begin{align*}(T_{h}^1)_{2,1}&=-2\cdot q_0^{1}\cdot q_1^{1}+2\cdot q_2^{1}\cdot q_3^{1}\end{align*}
\begin{align*}(T_{h}^1)_{2,2}&=1\cdot q_0^{2}-1\cdot q_1^{2}+1\cdot q_2^{2}-1\cdot q_3^{2}\end{align*}
\begin{align*}(T_{h}^1)_{2,3}&=2\cdot q_0^{1}\cdot q_3^{1}+2\cdot q_1^{1}\cdot q_2^{1}\end{align*}
\begin{align*}(T_{h}^1)_{3,1}&=2\cdot q_0^{1}\cdot q_2^{1}+2\cdot q_1^{1}\cdot q_3^{1}\end{align*}
\begin{align*}(T_{h}^1)_{3,2}&=-2\cdot q_0^{1}\cdot q_3^{1}+2\cdot q_1^{1}\cdot q_2^{1}\end{align*}
\begin{align*}(T_{h}^1)_{3,3}&=1\cdot q_0^{2}+1\cdot q_1^{2}-1\cdot q_2^{2}-1\cdot q_3^{2}\end{align*}
\section{Term $\ell=2$}
\subsection{Matrix $D_2$}
We will write the matrix $D^{-1}$ instead of $D$. The matrices $D$ and $D^{-1}$ are diagonal, so one is easy to construct from the other. The matrix $D^{-1}$ is more convenient becauce the denominators of diagonal terms are powers of $2$. Therefore, their binary expressions are finite and there is no rounding in computer storage.\[D_2^{-1}=\left[\begin{array}{ccccc}
15 & 0 & 0 & 0 & 0\\
0 & \frac{15}{2^{2}} & 0 & 0 & 0\\
0 & 0 & 15 & 0 & 0\\
0 & 0 & 0 & 15 & 0\\
0 & 0 & 0 & 0 & \frac{5}{2^{2}}
\end{array}\right].\]
\subsection{Matrix $C_2$}
\[C_2=\left[\begin{array}{ccccc}
\frac{35}{2^{19}} & 0 & 0 & 0 & 0\\
0 & \frac{261701}{2^{19}} & 0 & 0 & \frac{62205}{2^{19}}\\
0 & 0 & \frac{35}{2^{19}} & 0 & 0\\
0 & 0 & 0 & \frac{35}{2^{19}} & 0\\
0 & \frac{62205}{2^{19}} & 0 & 0 & \frac{15213}{2^{19}}
\end{array}\right].\]
\subsection{Matrix $T_{h}^2(q_0,q_1,q_2,q_3)$}
\begin{align*}(T_{h}^2)_{1,1}&=1\cdot q_0^{4}+1\cdot q_1^{4}-1\cdot q_2^{4}-1\cdot q_3^{4}-6\cdot q_0^{2}\cdot q_1^{2}+6\cdot q_2^{2}\cdot q_3^{2}\end{align*}
\begin{align*}(T_{h}^2)_{1,2}&=-8\cdot q_0^{1}\cdot q_1^{3}+8\cdot q_0^{3}\cdot q_1^{1}-8\cdot q_2^{1}\cdot q_3^{3}+8\cdot q_2^{3}\cdot q_3^{1}\end{align*}
\begin{align*}(T_{h}^2)_{1,3}&=2\cdot q_0^{1}\cdot q_3^{3}+2\cdot q_0^{3}\cdot q_3^{1}-2\cdot q_1^{1}\cdot q_2^{3}-2\cdot q_1^{3}\cdot q_2^{1}\\
&\quad -6\cdot q_0^{1}\cdot q_1^{2}\cdot q_3^{1}-6\cdot q_0^{1}\cdot q_2^{2}\cdot q_3^{1}+6\cdot q_0^{2}\cdot q_1^{1}\cdot q_2^{1}\\
&\quad +6\cdot q_1^{1}\cdot q_2^{1}\cdot q_3^{2}\end{align*}
\begin{align*}(T_{h}^2)_{1,4}&=-2\cdot q_0^{1}\cdot q_2^{3}-2\cdot q_0^{3}\cdot q_2^{1}-2\cdot q_1^{1}\cdot q_3^{3}-2\cdot q_1^{3}\cdot q_3^{1}\\
&\quad +6\cdot q_0^{1}\cdot q_1^{2}\cdot q_2^{1}+6\cdot q_0^{1}\cdot q_2^{1}\cdot q_3^{2}+6\cdot q_0^{2}\cdot q_1^{1}\cdot q_3^{1}\\
&\quad +6\cdot q_1^{1}\cdot q_2^{2}\cdot q_3^{1}\end{align*}
\begin{align*}(T_{h}^2)_{1,5}&=-24\cdot q_0^{1}\cdot q_1^{1}\cdot q_2^{2}+24\cdot q_0^{1}\cdot q_1^{1}\cdot q_3^{2}-24\cdot q_0^{2}\cdot q_2^{1}\cdot q_3^{1}\\
&\quad +24\cdot q_1^{2}\cdot q_2^{1}\cdot q_3^{1}\end{align*}
\begin{align*}(T_{h}^2)_{2,1}&=2\cdot q_0^{1}\cdot q_1^{3}-2\cdot q_0^{3}\cdot q_1^{1}-2\cdot q_2^{1}\cdot q_3^{3}+2\cdot q_2^{3}\cdot q_3^{1}\end{align*}
\begin{align*}(T_{h}^2)_{2,2}&=1\cdot q_0^{4}+1\cdot q_1^{4}+1\cdot q_2^{4}+1\cdot q_3^{4}-6\cdot q_0^{2}\cdot q_1^{2}-6\cdot q_2^{2}\cdot q_3^{2}\end{align*}
\begin{align*}(T_{h}^2)_{2,3}&=-1\cdot q_0^{1}\cdot q_2^{3}+1\cdot q_0^{3}\cdot q_2^{1}-1\cdot q_1^{1}\cdot q_3^{3}+1\cdot q_1^{3}\cdot q_3^{1}\\
&\quad -3\cdot q_0^{1}\cdot q_1^{2}\cdot q_2^{1}+3\cdot q_0^{1}\cdot q_2^{1}\cdot q_3^{2}-3\cdot q_0^{2}\cdot q_1^{1}\cdot q_3^{1}\\
&\quad +3\cdot q_1^{1}\cdot q_2^{2}\cdot q_3^{1}\end{align*}
\begin{align*}(T_{h}^2)_{2,4}&=-1\cdot q_0^{1}\cdot q_3^{3}+1\cdot q_0^{3}\cdot q_3^{1}+1\cdot q_1^{1}\cdot q_2^{3}-1\cdot q_1^{3}\cdot q_2^{1}\\
&\quad -3\cdot q_0^{1}\cdot q_1^{2}\cdot q_3^{1}+3\cdot q_0^{1}\cdot q_2^{2}\cdot q_3^{1}+3\cdot q_0^{2}\cdot q_1^{1}\cdot q_2^{1}\\
&\quad -3\cdot q_1^{1}\cdot q_2^{1}\cdot q_3^{2}\end{align*}
\begin{align*}(T_{h}^2)_{2,5}&=-6\cdot q_0^{2}\cdot q_2^{2}+6\cdot q_0^{2}\cdot q_3^{2}+6\cdot q_1^{2}\cdot q_2^{2}-6\cdot q_1^{2}\cdot q_3^{2}\\
&\quad +24\cdot q_0^{1}\cdot q_1^{1}\cdot q_2^{1}\cdot q_3^{1}\end{align*}
\begin{align*}(T_{h}^2)_{3,1}&=-2\cdot q_0^{1}\cdot q_3^{3}-2\cdot q_0^{3}\cdot q_3^{1}-2\cdot q_1^{1}\cdot q_2^{3}-2\cdot q_1^{3}\cdot q_2^{1}\\
&\quad +6\cdot q_0^{1}\cdot q_1^{2}\cdot q_3^{1}+6\cdot q_0^{1}\cdot q_2^{2}\cdot q_3^{1}+6\cdot q_0^{2}\cdot q_1^{1}\cdot q_2^{1}\\
&\quad +6\cdot q_1^{1}\cdot q_2^{1}\cdot q_3^{2}\end{align*}
\begin{align*}(T_{h}^2)_{3,2}&=4\cdot q_0^{1}\cdot q_2^{3}-4\cdot q_0^{3}\cdot q_2^{1}-4\cdot q_1^{1}\cdot q_3^{3}+4\cdot q_1^{3}\cdot q_3^{1}\\
&\quad +12\cdot q_0^{1}\cdot q_1^{2}\cdot q_2^{1}-12\cdot q_0^{1}\cdot q_2^{1}\cdot q_3^{2}-12\cdot q_0^{2}\cdot q_1^{1}\cdot q_3^{1}\\
&\quad +12\cdot q_1^{1}\cdot q_2^{2}\cdot q_3^{1}\end{align*}
\begin{align*}(T_{h}^2)_{3,3}&=1\cdot q_0^{4}-1\cdot q_1^{4}+1\cdot q_2^{4}-1\cdot q_3^{4}-6\cdot q_0^{2}\cdot q_2^{2}+6\cdot q_1^{2}\cdot q_3^{2}\end{align*}
\begin{align*}(T_{h}^2)_{3,4}&=2\cdot q_0^{1}\cdot q_1^{3}+2\cdot q_0^{3}\cdot q_1^{1}-2\cdot q_2^{1}\cdot q_3^{3}-2\cdot q_2^{3}\cdot q_3^{1}\\
&\quad -6\cdot q_0^{1}\cdot q_1^{1}\cdot q_2^{2}-6\cdot q_0^{1}\cdot q_1^{1}\cdot q_3^{2}+6\cdot q_0^{2}\cdot q_2^{1}\cdot q_3^{1}\\
&\quad +6\cdot q_1^{2}\cdot q_2^{1}\cdot q_3^{1}\end{align*}
\begin{align*}(T_{h}^2)_{3,5}&=12\cdot q_0^{1}\cdot q_2^{3}-12\cdot q_0^{3}\cdot q_2^{1}-12\cdot q_1^{1}\cdot q_3^{3}+12\cdot q_1^{3}\cdot q_3^{1}\\
&\quad -12\cdot q_0^{1}\cdot q_1^{2}\cdot q_2^{1}+12\cdot q_0^{1}\cdot q_2^{1}\cdot q_3^{2}+12\cdot q_0^{2}\cdot q_1^{1}\cdot q_3^{1}\\
&\quad -12\cdot q_1^{1}\cdot q_2^{2}\cdot q_3^{1}\end{align*}
\begin{align*}(T_{h}^2)_{4,1}&=2\cdot q_0^{1}\cdot q_2^{3}+2\cdot q_0^{3}\cdot q_2^{1}-2\cdot q_1^{1}\cdot q_3^{3}-2\cdot q_1^{3}\cdot q_3^{1}\\
&\quad -6\cdot q_0^{1}\cdot q_1^{2}\cdot q_2^{1}-6\cdot q_0^{1}\cdot q_2^{1}\cdot q_3^{2}+6\cdot q_0^{2}\cdot q_1^{1}\cdot q_3^{1}\\
&\quad +6\cdot q_1^{1}\cdot q_2^{2}\cdot q_3^{1}\end{align*}
\begin{align*}(T_{h}^2)_{4,2}&=4\cdot q_0^{1}\cdot q_3^{3}-4\cdot q_0^{3}\cdot q_3^{1}+4\cdot q_1^{1}\cdot q_2^{3}-4\cdot q_1^{3}\cdot q_2^{1}\\
&\quad +12\cdot q_0^{1}\cdot q_1^{2}\cdot q_3^{1}-12\cdot q_0^{1}\cdot q_2^{2}\cdot q_3^{1}+12\cdot q_0^{2}\cdot q_1^{1}\cdot q_2^{1}\\
&\quad -12\cdot q_1^{1}\cdot q_2^{1}\cdot q_3^{2}\end{align*}
\begin{align*}(T_{h}^2)_{4,3}&=-2\cdot q_0^{1}\cdot q_1^{3}-2\cdot q_0^{3}\cdot q_1^{1}-2\cdot q_2^{1}\cdot q_3^{3}-2\cdot q_2^{3}\cdot q_3^{1}\\
&\quad +6\cdot q_0^{1}\cdot q_1^{1}\cdot q_2^{2}+6\cdot q_0^{1}\cdot q_1^{1}\cdot q_3^{2}+6\cdot q_0^{2}\cdot q_2^{1}\cdot q_3^{1}\\
&\quad +6\cdot q_1^{2}\cdot q_2^{1}\cdot q_3^{1}\end{align*}
\begin{align*}(T_{h}^2)_{4,4}&=1\cdot q_0^{4}-1\cdot q_1^{4}-1\cdot q_2^{4}+1\cdot q_3^{4}-6\cdot q_0^{2}\cdot q_3^{2}+6\cdot q_1^{2}\cdot q_2^{2}\end{align*}
\begin{align*}(T_{h}^2)_{4,5}&=-12\cdot q_0^{1}\cdot q_3^{3}+12\cdot q_0^{3}\cdot q_3^{1}-12\cdot q_1^{1}\cdot q_2^{3}+12\cdot q_1^{3}\cdot q_2^{1}\\
&\quad +12\cdot q_0^{1}\cdot q_1^{2}\cdot q_3^{1}-12\cdot q_0^{1}\cdot q_2^{2}\cdot q_3^{1}+12\cdot q_0^{2}\cdot q_1^{1}\cdot q_2^{1}\\
&\quad -12\cdot q_1^{1}\cdot q_2^{1}\cdot q_3^{2}\end{align*}
\begin{align*}(T_{h}^2)_{5,1}&=2\cdot q_0^{1}\cdot q_1^{1}\cdot q_2^{2}-2\cdot q_0^{1}\cdot q_1^{1}\cdot q_3^{2}-2\cdot q_0^{2}\cdot q_2^{1}\cdot q_3^{1}\\
&\quad +2\cdot q_1^{2}\cdot q_2^{1}\cdot q_3^{1}\end{align*}
\begin{align*}(T_{h}^2)_{5,2}&=-2\cdot q_0^{2}\cdot q_2^{2}+2\cdot q_0^{2}\cdot q_3^{2}+2\cdot q_1^{2}\cdot q_2^{2}-2\cdot q_1^{2}\cdot q_3^{2}\\
&\quad -8\cdot q_0^{1}\cdot q_1^{1}\cdot q_2^{1}\cdot q_3^{1}\end{align*}
\begin{align*}(T_{h}^2)_{5,3}&=-1\cdot q_0^{1}\cdot q_2^{3}+1\cdot q_0^{3}\cdot q_2^{1}-1\cdot q_1^{1}\cdot q_3^{3}+1\cdot q_1^{3}\cdot q_3^{1}\\
&\quad +1\cdot q_0^{1}\cdot q_1^{2}\cdot q_2^{1}-1\cdot q_0^{1}\cdot q_2^{1}\cdot q_3^{2}+1\cdot q_0^{2}\cdot q_1^{1}\cdot q_3^{1}\\
&\quad -1\cdot q_1^{1}\cdot q_2^{2}\cdot q_3^{1}\end{align*}
\begin{align*}(T_{h}^2)_{5,4}&=1\cdot q_0^{1}\cdot q_3^{3}-1\cdot q_0^{3}\cdot q_3^{1}-1\cdot q_1^{1}\cdot q_2^{3}+1\cdot q_1^{3}\cdot q_2^{1}\\
&\quad -1\cdot q_0^{1}\cdot q_1^{2}\cdot q_3^{1}+1\cdot q_0^{1}\cdot q_2^{2}\cdot q_3^{1}+1\cdot q_0^{2}\cdot q_1^{1}\cdot q_2^{1}\\
&\quad -1\cdot q_1^{1}\cdot q_2^{1}\cdot q_3^{2}\end{align*}
\begin{align*}(T_{h}^2)_{5,5}&=1\cdot q_0^{4}+1\cdot q_1^{4}+1\cdot q_2^{4}+1\cdot q_3^{4}+2\cdot q_0^{2}\cdot q_1^{2}-4\cdot q_0^{2}\cdot q_2^{2}\\
&\quad -4\cdot q_0^{2}\cdot q_3^{2}-4\cdot q_1^{2}\cdot q_2^{2}-4\cdot q_1^{2}\cdot q_3^{2}+2\cdot q_2^{2}\cdot q_3^{2}\end{align*}
\section{Term $\ell=3$}
\subsection{Matrix $D_3$}
We will write the matrix $D^{-1}$ instead of $D$. The matrices $D$ and $D^{-1}$ are diagonal, so one is easy to construct from the other. The matrix $D^{-1}$ is more convenient becauce the denominators of diagonal terms are powers of $2$. Therefore, their binary expressions are finite and there is no rounding in computer storage.\[D_3^{-1}=\left[\begin{array}{ccccccc}
\frac{35}{2^{3}} & 0 & 0 & 0 & 0 & 0 & 0\\
0 & \frac{35}{2^{3}} & 0 & 0 & 0 & 0 & 0\\
0 & 0 & 105 & 0 & 0 & 0 & 0\\
0 & 0 & 0 & \frac{105}{2^{2}} & 0 & 0 & 0\\
0 & 0 & 0 & 0 & \frac{21}{2^{3}} & 0 & 0\\
0 & 0 & 0 & 0 & 0 & \frac{21}{2^{3}} & 0\\
0 & 0 & 0 & 0 & 0 & 0 & \frac{7}{2^{2}}
\end{array}\right].\]
\subsection{Matrix $C_3$}
\[C_3=\left[\begin{array}{ccccccc}
\frac{252053}{2^{20}} & 0 & 0 & 0 & -\frac{49961}{2^{20}} & 0 & 0\\
0 & \frac{15}{2^{16}} & 0 & 0 & 0 & 0 & 0\\
0 & 0 & \frac{5}{2^{19}} & 0 & 0 & 0 & 0\\
0 & 0 & 0 & \frac{5}{2^{17}} & 0 & 0 & 0\\
-\frac{49961}{2^{20}} & 0 & 0 & 0 & \frac{1289}{2^{17}} & 0 & 0\\
0 & 0 & 0 & 0 & 0 & \frac{399}{2^{20}} & 0\\
0 & 0 & 0 & 0 & 0 & 0 & \frac{599}{2^{20}}
\end{array}\right].\]
\subsection{Matrix $T_{h}^3(q_0,q_1,q_2,q_3)$}
\begin{align*}(T_{h}^3)_{1,1}&=1\cdot q_0^{6}-1\cdot q_1^{6}-1\cdot q_2^{6}+1\cdot q_3^{6}+15\cdot q_0^{2}\cdot q_1^{4}-15\cdot q_0^{4}\cdot q_1^{2}\\
&\quad -15\cdot q_2^{2}\cdot q_3^{4}+15\cdot q_2^{4}\cdot q_3^{2}\end{align*}
\begin{align*}(T_{h}^3)_{1,2}&=6\cdot q_0^{1}\cdot q_1^{5}-20\cdot q_0^{3}\cdot q_1^{3}+6\cdot q_0^{5}\cdot q_1^{1}+6\cdot q_2^{1}\cdot q_3^{5}\\
&\quad -20\cdot q_2^{3}\cdot q_3^{3}+6\cdot q_2^{5}\cdot q_3^{1}\end{align*}
\begin{align*}(T_{h}^3)_{1,3}&=-\frac{1}{2^{1}}\cdot q_0^{1}\cdot q_3^{5}+\frac{1}{2^{1}}\cdot q_0^{5}\cdot q_3^{1}-\frac{1}{2^{1}}\cdot q_1^{1}\cdot q_2^{5}\\
&\quad +\frac{1}{2^{1}}\cdot q_1^{5}\cdot q_2^{1}+\frac{5}{2^{1}}\cdot q_0^{1}\cdot q_1^{4}\cdot q_3^{1}+5\cdot q_0^{1}\cdot q_2^{2}\cdot q_3^{3}\\
&\quad -\frac{5}{2^{1}}\cdot q_0^{1}\cdot q_2^{4}\cdot q_3^{1}-5\cdot q_0^{2}\cdot q_1^{3}\cdot q_2^{1}-5\cdot q_0^{3}\cdot q_1^{2}\cdot q_3^{1}\\
&\quad +\frac{5}{2^{1}}\cdot q_0^{4}\cdot q_1^{1}\cdot q_2^{1}-\frac{5}{2^{1}}\cdot q_1^{1}\cdot q_2^{1}\cdot q_3^{4}\\
&\quad +5\cdot q_1^{1}\cdot q_2^{3}\cdot q_3^{2}\end{align*}
\begin{align*}(T_{h}^3)_{1,4}&=-1\cdot q_0^{1}\cdot q_2^{5}-1\cdot q_0^{5}\cdot q_2^{1}+1\cdot q_1^{1}\cdot q_3^{5}+1\cdot q_1^{5}\cdot q_3^{1}\\
&\quad -5\cdot q_0^{1}\cdot q_1^{4}\cdot q_2^{1}-5\cdot q_0^{1}\cdot q_2^{1}\cdot q_3^{4}+10\cdot q_0^{1}\cdot q_2^{3}\cdot q_3^{2}\\
&\quad -10\cdot q_0^{2}\cdot q_1^{3}\cdot q_3^{1}+10\cdot q_0^{3}\cdot q_1^{2}\cdot q_2^{1}+5\cdot q_0^{4}\cdot q_1^{1}\cdot q_3^{1}\\
&\quad -10\cdot q_1^{1}\cdot q_2^{2}\cdot q_3^{3}+5\cdot q_1^{1}\cdot q_2^{4}\cdot q_3^{1}\end{align*}
\begin{align*}(T_{h}^3)_{1,5}&=5\cdot q_0^{2}\cdot q_2^{4}+5\cdot q_0^{2}\cdot q_3^{4}-5\cdot q_0^{4}\cdot q_2^{2}+5\cdot q_0^{4}\cdot q_3^{2}\\
&\quad -5\cdot q_1^{2}\cdot q_2^{4}-5\cdot q_1^{2}\cdot q_3^{4}-5\cdot q_1^{4}\cdot q_2^{2}+5\cdot q_1^{4}\cdot q_3^{2}\\
&\quad +30\cdot q_0^{2}\cdot q_1^{2}\cdot q_2^{2}-30\cdot q_0^{2}\cdot q_1^{2}\cdot q_3^{2}-30\cdot q_0^{2}\cdot q_2^{2}\cdot q_3^{2}\\
&\quad +30\cdot q_1^{2}\cdot q_2^{2}\cdot q_3^{2}+40\cdot q_0^{1}\cdot q_1^{1}\cdot q_2^{1}\cdot q_3^{3}-40\cdot q_0^{1}\cdot q_1^{1}\cdot q_2^{3}\cdot q_3^{1}\\
&\quad -40\cdot q_0^{1}\cdot q_1^{3}\cdot q_2^{1}\cdot q_3^{1}+40\cdot q_0^{3}\cdot q_1^{1}\cdot q_2^{1}\cdot q_3^{1}\end{align*}
\begin{align*}(T_{h}^3)_{1,6}&=-10\cdot q_0^{1}\cdot q_1^{1}\cdot q_2^{4}-10\cdot q_0^{1}\cdot q_1^{1}\cdot q_3^{4}+20\cdot q_0^{1}\cdot q_1^{3}\cdot q_2^{2}\\
&\quad -20\cdot q_0^{1}\cdot q_1^{3}\cdot q_3^{2}+20\cdot q_0^{2}\cdot q_2^{1}\cdot q_3^{3}-20\cdot q_0^{2}\cdot q_2^{3}\cdot q_3^{1}\\
&\quad -20\cdot q_0^{3}\cdot q_1^{1}\cdot q_2^{2}+20\cdot q_0^{3}\cdot q_1^{1}\cdot q_3^{2}-10\cdot q_0^{4}\cdot q_2^{1}\cdot q_3^{1}\\
&\quad -20\cdot q_1^{2}\cdot q_2^{1}\cdot q_3^{3}+20\cdot q_1^{2}\cdot q_2^{3}\cdot q_3^{1}-10\cdot q_1^{4}\cdot q_2^{1}\cdot q_3^{1}\\
&\quad +60\cdot q_0^{1}\cdot q_1^{1}\cdot q_2^{2}\cdot q_3^{2}+60\cdot q_0^{2}\cdot q_1^{2}\cdot q_2^{1}\cdot q_3^{1}\end{align*}
\begin{align*}(T_{h}^3)_{1,7}&=10\cdot q_0^{3}\cdot q_2^{3}-10\cdot q_1^{3}\cdot q_3^{3}-30\cdot q_0^{1}\cdot q_1^{2}\cdot q_2^{3}+30\cdot q_0^{2}\cdot q_1^{1}\cdot q_3^{3}\\
&\quad -30\cdot q_0^{3}\cdot q_2^{1}\cdot q_3^{2}+30\cdot q_1^{3}\cdot q_2^{2}\cdot q_3^{1}+90\cdot q_0^{1}\cdot q_1^{2}\cdot q_2^{1}\cdot q_3^{2}\\
&\quad -90\cdot q_0^{2}\cdot q_1^{1}\cdot q_2^{2}\cdot q_3^{1}\end{align*}
\begin{align*}(T_{h}^3)_{2,1}&=-6\cdot q_0^{1}\cdot q_1^{5}+20\cdot q_0^{3}\cdot q_1^{3}-6\cdot q_0^{5}\cdot q_1^{1}+6\cdot q_2^{1}\cdot q_3^{5}\\
&\quad -20\cdot q_2^{3}\cdot q_3^{3}+6\cdot q_2^{5}\cdot q_3^{1}\end{align*}
\begin{align*}(T_{h}^3)_{2,2}&=1\cdot q_0^{6}-1\cdot q_1^{6}+1\cdot q_2^{6}-1\cdot q_3^{6}+15\cdot q_0^{2}\cdot q_1^{4}-15\cdot q_0^{4}\cdot q_1^{2}\\
&\quad +15\cdot q_2^{2}\cdot q_3^{4}-15\cdot q_2^{4}\cdot q_3^{2}\end{align*}
\begin{align*}(T_{h}^3)_{2,3}&=-\frac{1}{2^{1}}\cdot q_0^{1}\cdot q_2^{5}+\frac{1}{2^{1}}\cdot q_0^{5}\cdot q_2^{1}+\frac{1}{2^{1}}\cdot q_1^{1}\cdot q_3^{5}\\
&\quad -\frac{1}{2^{1}}\cdot q_1^{5}\cdot q_3^{1}+\frac{5}{2^{1}}\cdot q_0^{1}\cdot q_1^{4}\cdot q_2^{1}-\frac{5}{2^{1}}\cdot q_0^{1}\cdot q_2^{1}\cdot q_3^{4}\\
&\quad +5\cdot q_0^{1}\cdot q_2^{3}\cdot q_3^{2}+5\cdot q_0^{2}\cdot q_1^{3}\cdot q_3^{1}-5\cdot q_0^{3}\cdot q_1^{2}\cdot q_2^{1}\\
&\quad -\frac{5}{2^{1}}\cdot q_0^{4}\cdot q_1^{1}\cdot q_3^{1}-5\cdot q_1^{1}\cdot q_2^{2}\cdot q_3^{3}+\frac{5}{2^{1}}\cdot q_1^{1}\cdot q_2^{4}\cdot q_3^{1}\end{align*}
\begin{align*}(T_{h}^3)_{2,4}&=1\cdot q_0^{1}\cdot q_3^{5}+1\cdot q_0^{5}\cdot q_3^{1}+1\cdot q_1^{1}\cdot q_2^{5}+1\cdot q_1^{5}\cdot q_2^{1}\\
&\quad +5\cdot q_0^{1}\cdot q_1^{4}\cdot q_3^{1}-10\cdot q_0^{1}\cdot q_2^{2}\cdot q_3^{3}+5\cdot q_0^{1}\cdot q_2^{4}\cdot q_3^{1}\\
&\quad -10\cdot q_0^{2}\cdot q_1^{3}\cdot q_2^{1}-10\cdot q_0^{3}\cdot q_1^{2}\cdot q_3^{1}+5\cdot q_0^{4}\cdot q_1^{1}\cdot q_2^{1}\\
&\quad +5\cdot q_1^{1}\cdot q_2^{1}\cdot q_3^{4}-10\cdot q_1^{1}\cdot q_2^{3}\cdot q_3^{2}\end{align*}
\begin{align*}(T_{h}^3)_{2,5}&=-10\cdot q_0^{1}\cdot q_1^{1}\cdot q_2^{4}-10\cdot q_0^{1}\cdot q_1^{1}\cdot q_3^{4}-20\cdot q_0^{1}\cdot q_1^{3}\cdot q_2^{2}\\
&\quad +20\cdot q_0^{1}\cdot q_1^{3}\cdot q_3^{2}+20\cdot q_0^{2}\cdot q_2^{1}\cdot q_3^{3}-20\cdot q_0^{2}\cdot q_2^{3}\cdot q_3^{1}\\
&\quad +20\cdot q_0^{3}\cdot q_1^{1}\cdot q_2^{2}-20\cdot q_0^{3}\cdot q_1^{1}\cdot q_3^{2}+10\cdot q_0^{4}\cdot q_2^{1}\cdot q_3^{1}\\
&\quad -20\cdot q_1^{2}\cdot q_2^{1}\cdot q_3^{3}+20\cdot q_1^{2}\cdot q_2^{3}\cdot q_3^{1}+10\cdot q_1^{4}\cdot q_2^{1}\cdot q_3^{1}\\
&\quad +60\cdot q_0^{1}\cdot q_1^{1}\cdot q_2^{2}\cdot q_3^{2}-60\cdot q_0^{2}\cdot q_1^{2}\cdot q_2^{1}\cdot q_3^{1}\end{align*}
\begin{align*}(T_{h}^3)_{2,6}&=-5\cdot q_0^{2}\cdot q_2^{4}-5\cdot q_0^{2}\cdot q_3^{4}-5\cdot q_0^{4}\cdot q_2^{2}+5\cdot q_0^{4}\cdot q_3^{2}\\
&\quad +5\cdot q_1^{2}\cdot q_2^{4}+5\cdot q_1^{2}\cdot q_3^{4}-5\cdot q_1^{4}\cdot q_2^{2}+5\cdot q_1^{4}\cdot q_3^{2}\\
&\quad +30\cdot q_0^{2}\cdot q_1^{2}\cdot q_2^{2}-30\cdot q_0^{2}\cdot q_1^{2}\cdot q_3^{2}+30\cdot q_0^{2}\cdot q_2^{2}\cdot q_3^{2}\\
&\quad -30\cdot q_1^{2}\cdot q_2^{2}\cdot q_3^{2}-40\cdot q_0^{1}\cdot q_1^{1}\cdot q_2^{1}\cdot q_3^{3}+40\cdot q_0^{1}\cdot q_1^{1}\cdot q_2^{3}\cdot q_3^{1}\\
&\quad -40\cdot q_0^{1}\cdot q_1^{3}\cdot q_2^{1}\cdot q_3^{1}+40\cdot q_0^{3}\cdot q_1^{1}\cdot q_2^{1}\cdot q_3^{1}\end{align*}
\begin{align*}(T_{h}^3)_{2,7}&=10\cdot q_0^{3}\cdot q_3^{3}+10\cdot q_1^{3}\cdot q_2^{3}-30\cdot q_0^{1}\cdot q_1^{2}\cdot q_3^{3}-30\cdot q_0^{2}\cdot q_1^{1}\cdot q_2^{3}\\
&\quad -30\cdot q_0^{3}\cdot q_2^{2}\cdot q_3^{1}-30\cdot q_1^{3}\cdot q_2^{1}\cdot q_3^{2}+90\cdot q_0^{1}\cdot q_1^{2}\cdot q_2^{2}\cdot q_3^{1}\\
&\quad +90\cdot q_0^{2}\cdot q_1^{1}\cdot q_2^{1}\cdot q_3^{2}\end{align*}
\begin{align*}(T_{h}^3)_{3,1}&=12\cdot q_0^{1}\cdot q_3^{5}-12\cdot q_0^{5}\cdot q_3^{1}-12\cdot q_1^{1}\cdot q_2^{5}+12\cdot q_1^{5}\cdot q_2^{1}\\
&\quad -60\cdot q_0^{1}\cdot q_1^{4}\cdot q_3^{1}-120\cdot q_0^{1}\cdot q_2^{2}\cdot q_3^{3}+60\cdot q_0^{1}\cdot q_2^{4}\cdot q_3^{1}\\
&\quad -120\cdot q_0^{2}\cdot q_1^{3}\cdot q_2^{1}+120\cdot q_0^{3}\cdot q_1^{2}\cdot q_3^{1}+60\cdot q_0^{4}\cdot q_1^{1}\cdot q_2^{1}\\
&\quad -60\cdot q_1^{1}\cdot q_2^{1}\cdot q_3^{4}+120\cdot q_1^{1}\cdot q_2^{3}\cdot q_3^{2}\end{align*}
\begin{align*}(T_{h}^3)_{3,2}&=12\cdot q_0^{1}\cdot q_2^{5}-12\cdot q_0^{5}\cdot q_2^{1}+12\cdot q_1^{1}\cdot q_3^{5}-12\cdot q_1^{5}\cdot q_3^{1}\\
&\quad -60\cdot q_0^{1}\cdot q_1^{4}\cdot q_2^{1}+60\cdot q_0^{1}\cdot q_2^{1}\cdot q_3^{4}-120\cdot q_0^{1}\cdot q_2^{3}\cdot q_3^{2}\\
&\quad +120\cdot q_0^{2}\cdot q_1^{3}\cdot q_3^{1}+120\cdot q_0^{3}\cdot q_1^{2}\cdot q_2^{1}-60\cdot q_0^{4}\cdot q_1^{1}\cdot q_3^{1}\\
&\quad -120\cdot q_1^{1}\cdot q_2^{2}\cdot q_3^{3}+60\cdot q_1^{1}\cdot q_2^{4}\cdot q_3^{1}\end{align*}
\begin{align*}(T_{h}^3)_{3,3}&=1\cdot q_0^{6}+1\cdot q_1^{6}+1\cdot q_2^{6}+1\cdot q_3^{6}-5\cdot q_0^{2}\cdot q_1^{4}-5\cdot q_0^{2}\cdot q_2^{4}\\
&\quad -5\cdot q_0^{2}\cdot q_3^{4}-5\cdot q_0^{4}\cdot q_1^{2}-5\cdot q_0^{4}\cdot q_2^{2}-5\cdot q_0^{4}\cdot q_3^{2}\\
&\quad -5\cdot q_1^{2}\cdot q_2^{4}-5\cdot q_1^{2}\cdot q_3^{4}-5\cdot q_1^{4}\cdot q_2^{2}-5\cdot q_1^{4}\cdot q_3^{2}\\
&\quad -5\cdot q_2^{2}\cdot q_3^{4}-5\cdot q_2^{4}\cdot q_3^{2}+30\cdot q_0^{2}\cdot q_1^{2}\cdot q_2^{2}+30\cdot q_0^{2}\cdot q_1^{2}\cdot q_3^{2}\\
&\quad +30\cdot q_0^{2}\cdot q_2^{2}\cdot q_3^{2}+30\cdot q_1^{2}\cdot q_2^{2}\cdot q_3^{2}\end{align*}
\begin{align*}(T_{h}^3)_{3,4}&=-8\cdot q_0^{1}\cdot q_1^{5}+8\cdot q_0^{5}\cdot q_1^{1}+8\cdot q_2^{1}\cdot q_3^{5}-8\cdot q_2^{5}\cdot q_3^{1}\\
&\quad +40\cdot q_0^{1}\cdot q_1^{3}\cdot q_2^{2}+40\cdot q_0^{1}\cdot q_1^{3}\cdot q_3^{2}-40\cdot q_0^{2}\cdot q_2^{1}\cdot q_3^{3}\\
&\quad +40\cdot q_0^{2}\cdot q_2^{3}\cdot q_3^{1}-40\cdot q_0^{3}\cdot q_1^{1}\cdot q_2^{2}-40\cdot q_0^{3}\cdot q_1^{1}\cdot q_3^{2}\\
&\quad -40\cdot q_1^{2}\cdot q_2^{1}\cdot q_3^{3}+40\cdot q_1^{2}\cdot q_2^{3}\cdot q_3^{1}\end{align*}
\begin{align*}(T_{h}^3)_{3,5}&=-20\cdot q_0^{1}\cdot q_3^{5}+20\cdot q_0^{5}\cdot q_3^{1}+20\cdot q_1^{1}\cdot q_2^{5}-20\cdot q_1^{5}\cdot q_2^{1}\\
&\quad +160\cdot q_0^{1}\cdot q_1^{2}\cdot q_3^{3}-60\cdot q_0^{1}\cdot q_1^{4}\cdot q_3^{1}+40\cdot q_0^{1}\cdot q_2^{2}\cdot q_3^{3}\\
&\quad +60\cdot q_0^{1}\cdot q_2^{4}\cdot q_3^{1}-160\cdot q_0^{2}\cdot q_1^{1}\cdot q_2^{3}+40\cdot q_0^{2}\cdot q_1^{3}\cdot q_2^{1}\\
&\quad -40\cdot q_0^{3}\cdot q_1^{2}\cdot q_3^{1}-160\cdot q_0^{3}\cdot q_2^{2}\cdot q_3^{1}+60\cdot q_0^{4}\cdot q_1^{1}\cdot q_2^{1}\\
&\quad -60\cdot q_1^{1}\cdot q_2^{1}\cdot q_3^{4}-40\cdot q_1^{1}\cdot q_2^{3}\cdot q_3^{2}+160\cdot q_1^{3}\cdot q_2^{1}\cdot q_3^{2}\end{align*}
\begin{align*}(T_{h}^3)_{3,6}&=20\cdot q_0^{1}\cdot q_2^{5}-20\cdot q_0^{5}\cdot q_2^{1}+20\cdot q_1^{1}\cdot q_3^{5}-20\cdot q_1^{5}\cdot q_3^{1}\\
&\quad -160\cdot q_0^{1}\cdot q_1^{2}\cdot q_2^{3}+60\cdot q_0^{1}\cdot q_1^{4}\cdot q_2^{1}-60\cdot q_0^{1}\cdot q_2^{1}\cdot q_3^{4}\\
&\quad -40\cdot q_0^{1}\cdot q_2^{3}\cdot q_3^{2}-160\cdot q_0^{2}\cdot q_1^{1}\cdot q_3^{3}+40\cdot q_0^{2}\cdot q_1^{3}\cdot q_3^{1}\\
&\quad +40\cdot q_0^{3}\cdot q_1^{2}\cdot q_2^{1}+160\cdot q_0^{3}\cdot q_2^{1}\cdot q_3^{2}+60\cdot q_0^{4}\cdot q_1^{1}\cdot q_3^{1}\\
&\quad -40\cdot q_1^{1}\cdot q_2^{2}\cdot q_3^{3}-60\cdot q_1^{1}\cdot q_2^{4}\cdot q_3^{1}+160\cdot q_1^{3}\cdot q_2^{2}\cdot q_3^{1}\end{align*}
\begin{align*}(T_{h}^3)_{3,7}&=120\cdot q_0^{1}\cdot q_1^{1}\cdot q_2^{4}-120\cdot q_0^{1}\cdot q_1^{1}\cdot q_3^{4}-120\cdot q_0^{1}\cdot q_1^{3}\cdot q_2^{2}\\
&\quad +120\cdot q_0^{1}\cdot q_1^{3}\cdot q_3^{2}+120\cdot q_0^{2}\cdot q_2^{1}\cdot q_3^{3}+120\cdot q_0^{2}\cdot q_2^{3}\cdot q_3^{1}\\
&\quad -120\cdot q_0^{3}\cdot q_1^{1}\cdot q_2^{2}+120\cdot q_0^{3}\cdot q_1^{1}\cdot q_3^{2}-120\cdot q_0^{4}\cdot q_2^{1}\cdot q_3^{1}\\
&\quad -120\cdot q_1^{2}\cdot q_2^{1}\cdot q_3^{3}-120\cdot q_1^{2}\cdot q_2^{3}\cdot q_3^{1}+120\cdot q_1^{4}\cdot q_2^{1}\cdot q_3^{1}\end{align*}
\begin{align*}(T_{h}^3)_{4,1}&=6\cdot q_0^{1}\cdot q_2^{5}+6\cdot q_0^{5}\cdot q_2^{1}+6\cdot q_1^{1}\cdot q_3^{5}+6\cdot q_1^{5}\cdot q_3^{1}\\
&\quad +30\cdot q_0^{1}\cdot q_1^{4}\cdot q_2^{1}+30\cdot q_0^{1}\cdot q_2^{1}\cdot q_3^{4}-60\cdot q_0^{1}\cdot q_2^{3}\cdot q_3^{2}\\
&\quad -60\cdot q_0^{2}\cdot q_1^{3}\cdot q_3^{1}-60\cdot q_0^{3}\cdot q_1^{2}\cdot q_2^{1}+30\cdot q_0^{4}\cdot q_1^{1}\cdot q_3^{1}\\
&\quad -60\cdot q_1^{1}\cdot q_2^{2}\cdot q_3^{3}+30\cdot q_1^{1}\cdot q_2^{4}\cdot q_3^{1}\end{align*}
\begin{align*}(T_{h}^3)_{4,2}&=-6\cdot q_0^{1}\cdot q_3^{5}-6\cdot q_0^{5}\cdot q_3^{1}+6\cdot q_1^{1}\cdot q_2^{5}+6\cdot q_1^{5}\cdot q_2^{1}\\
&\quad -30\cdot q_0^{1}\cdot q_1^{4}\cdot q_3^{1}+60\cdot q_0^{1}\cdot q_2^{2}\cdot q_3^{3}-30\cdot q_0^{1}\cdot q_2^{4}\cdot q_3^{1}\\
&\quad -60\cdot q_0^{2}\cdot q_1^{3}\cdot q_2^{1}+60\cdot q_0^{3}\cdot q_1^{2}\cdot q_3^{1}+30\cdot q_0^{4}\cdot q_1^{1}\cdot q_2^{1}\\
&\quad +30\cdot q_1^{1}\cdot q_2^{1}\cdot q_3^{4}-60\cdot q_1^{1}\cdot q_2^{3}\cdot q_3^{2}\end{align*}
\begin{align*}(T_{h}^3)_{4,3}&=2\cdot q_0^{1}\cdot q_1^{5}-2\cdot q_0^{5}\cdot q_1^{1}+2\cdot q_2^{1}\cdot q_3^{5}-2\cdot q_2^{5}\cdot q_3^{1}\\
&\quad -10\cdot q_0^{1}\cdot q_1^{3}\cdot q_2^{2}-10\cdot q_0^{1}\cdot q_1^{3}\cdot q_3^{2}-10\cdot q_0^{2}\cdot q_2^{1}\cdot q_3^{3}\\
&\quad +10\cdot q_0^{2}\cdot q_2^{3}\cdot q_3^{1}+10\cdot q_0^{3}\cdot q_1^{1}\cdot q_2^{2}+10\cdot q_0^{3}\cdot q_1^{1}\cdot q_3^{2}\\
&\quad -10\cdot q_1^{2}\cdot q_2^{1}\cdot q_3^{3}+10\cdot q_1^{2}\cdot q_2^{3}\cdot q_3^{1}\end{align*}
\begin{align*}(T_{h}^3)_{4,4}&=1\cdot q_0^{6}+1\cdot q_1^{6}-1\cdot q_2^{6}-1\cdot q_3^{6}-5\cdot q_0^{2}\cdot q_1^{4}+5\cdot q_0^{2}\cdot q_2^{4}\\
&\quad +5\cdot q_0^{2}\cdot q_3^{4}-5\cdot q_0^{4}\cdot q_1^{2}-5\cdot q_0^{4}\cdot q_2^{2}-5\cdot q_0^{4}\cdot q_3^{2}\\
&\quad +5\cdot q_1^{2}\cdot q_2^{4}+5\cdot q_1^{2}\cdot q_3^{4}-5\cdot q_1^{4}\cdot q_2^{2}-5\cdot q_1^{4}\cdot q_3^{2}\\
&\quad +5\cdot q_2^{2}\cdot q_3^{4}+5\cdot q_2^{4}\cdot q_3^{2}+30\cdot q_0^{2}\cdot q_1^{2}\cdot q_2^{2}+30\cdot q_0^{2}\cdot q_1^{2}\cdot q_3^{2}\\
&\quad -30\cdot q_0^{2}\cdot q_2^{2}\cdot q_3^{2}-30\cdot q_1^{2}\cdot q_2^{2}\cdot q_3^{2}\end{align*}
\begin{align*}(T_{h}^3)_{4,5}&=10\cdot q_0^{1}\cdot q_2^{5}-40\cdot q_0^{3}\cdot q_2^{3}+10\cdot q_0^{5}\cdot q_2^{1}+10\cdot q_1^{1}\cdot q_3^{5}\\
&\quad -40\cdot q_1^{3}\cdot q_3^{3}+10\cdot q_1^{5}\cdot q_3^{1}+40\cdot q_0^{1}\cdot q_1^{2}\cdot q_2^{3}-30\cdot q_0^{1}\cdot q_1^{4}\cdot q_2^{1}\\
&\quad -30\cdot q_0^{1}\cdot q_2^{1}\cdot q_3^{4}-20\cdot q_0^{1}\cdot q_2^{3}\cdot q_3^{2}+40\cdot q_0^{2}\cdot q_1^{1}\cdot q_3^{3}\\
&\quad -20\cdot q_0^{2}\cdot q_1^{3}\cdot q_3^{1}-20\cdot q_0^{3}\cdot q_1^{2}\cdot q_2^{1}+40\cdot q_0^{3}\cdot q_2^{1}\cdot q_3^{2}\\
&\quad -30\cdot q_0^{4}\cdot q_1^{1}\cdot q_3^{1}-20\cdot q_1^{1}\cdot q_2^{2}\cdot q_3^{3}-30\cdot q_1^{1}\cdot q_2^{4}\cdot q_3^{1}\\
&\quad +40\cdot q_1^{3}\cdot q_2^{2}\cdot q_3^{1}+120\cdot q_0^{1}\cdot q_1^{2}\cdot q_2^{1}\cdot q_3^{2}+120\cdot q_0^{2}\cdot q_1^{1}\cdot q_2^{2}\cdot q_3^{1}\end{align*}
\begin{align*}(T_{h}^3)_{4,6}&=10\cdot q_0^{1}\cdot q_3^{5}-40\cdot q_0^{3}\cdot q_3^{3}+10\cdot q_0^{5}\cdot q_3^{1}-10\cdot q_1^{1}\cdot q_2^{5}\\
&\quad +40\cdot q_1^{3}\cdot q_2^{3}-10\cdot q_1^{5}\cdot q_2^{1}+40\cdot q_0^{1}\cdot q_1^{2}\cdot q_3^{3}-30\cdot q_0^{1}\cdot q_1^{4}\cdot q_3^{1}\\
&\quad -20\cdot q_0^{1}\cdot q_2^{2}\cdot q_3^{3}-30\cdot q_0^{1}\cdot q_2^{4}\cdot q_3^{1}-40\cdot q_0^{2}\cdot q_1^{1}\cdot q_2^{3}\\
&\quad +20\cdot q_0^{2}\cdot q_1^{3}\cdot q_2^{1}-20\cdot q_0^{3}\cdot q_1^{2}\cdot q_3^{1}+40\cdot q_0^{3}\cdot q_2^{2}\cdot q_3^{1}\\
&\quad +30\cdot q_0^{4}\cdot q_1^{1}\cdot q_2^{1}+30\cdot q_1^{1}\cdot q_2^{1}\cdot q_3^{4}+20\cdot q_1^{1}\cdot q_2^{3}\cdot q_3^{2}\\
&\quad -40\cdot q_1^{3}\cdot q_2^{1}\cdot q_3^{2}+120\cdot q_0^{1}\cdot q_1^{2}\cdot q_2^{2}\cdot q_3^{1}-120\cdot q_0^{2}\cdot q_1^{1}\cdot q_2^{1}\cdot q_3^{2}\end{align*}
\begin{align*}(T_{h}^3)_{4,7}&=30\cdot q_0^{2}\cdot q_2^{4}-30\cdot q_0^{2}\cdot q_3^{4}-30\cdot q_0^{4}\cdot q_2^{2}+30\cdot q_0^{4}\cdot q_3^{2}\\
&\quad -30\cdot q_1^{2}\cdot q_2^{4}+30\cdot q_1^{2}\cdot q_3^{4}+30\cdot q_1^{4}\cdot q_2^{2}-30\cdot q_1^{4}\cdot q_3^{2}\\
&\quad -120\cdot q_0^{1}\cdot q_1^{1}\cdot q_2^{1}\cdot q_3^{3}-120\cdot q_0^{1}\cdot q_1^{1}\cdot q_2^{3}\cdot q_3^{1}\\
&\quad +120\cdot q_0^{1}\cdot q_1^{3}\cdot q_2^{1}\cdot q_3^{1}+120\cdot q_0^{3}\cdot q_1^{1}\cdot q_2^{1}\cdot q_3^{1}\end{align*}
\begin{align*}(T_{h}^3)_{5,1}&=3\cdot q_0^{2}\cdot q_2^{4}+3\cdot q_0^{2}\cdot q_3^{4}-3\cdot q_0^{4}\cdot q_2^{2}+3\cdot q_0^{4}\cdot q_3^{2}\\
&\quad -3\cdot q_1^{2}\cdot q_2^{4}-3\cdot q_1^{2}\cdot q_3^{4}-3\cdot q_1^{4}\cdot q_2^{2}+3\cdot q_1^{4}\cdot q_3^{2}\\
&\quad +18\cdot q_0^{2}\cdot q_1^{2}\cdot q_2^{2}-18\cdot q_0^{2}\cdot q_1^{2}\cdot q_3^{2}-18\cdot q_0^{2}\cdot q_2^{2}\cdot q_3^{2}\\
&\quad +18\cdot q_1^{2}\cdot q_2^{2}\cdot q_3^{2}-24\cdot q_0^{1}\cdot q_1^{1}\cdot q_2^{1}\cdot q_3^{3}+24\cdot q_0^{1}\cdot q_1^{1}\cdot q_2^{3}\cdot q_3^{1}\\
&\quad +24\cdot q_0^{1}\cdot q_1^{3}\cdot q_2^{1}\cdot q_3^{1}-24\cdot q_0^{3}\cdot q_1^{1}\cdot q_2^{1}\cdot q_3^{1}\end{align*}
\begin{align*}(T_{h}^3)_{5,2}&=6\cdot q_0^{1}\cdot q_1^{1}\cdot q_2^{4}+6\cdot q_0^{1}\cdot q_1^{1}\cdot q_3^{4}+12\cdot q_0^{1}\cdot q_1^{3}\cdot q_2^{2}\\
&\quad -12\cdot q_0^{1}\cdot q_1^{3}\cdot q_3^{2}+12\cdot q_0^{2}\cdot q_2^{1}\cdot q_3^{3}-12\cdot q_0^{2}\cdot q_2^{3}\cdot q_3^{1}\\
&\quad -12\cdot q_0^{3}\cdot q_1^{1}\cdot q_2^{2}+12\cdot q_0^{3}\cdot q_1^{1}\cdot q_3^{2}+6\cdot q_0^{4}\cdot q_2^{1}\cdot q_3^{1}\\
&\quad -12\cdot q_1^{2}\cdot q_2^{1}\cdot q_3^{3}+12\cdot q_1^{2}\cdot q_2^{3}\cdot q_3^{1}+6\cdot q_1^{4}\cdot q_2^{1}\cdot q_3^{1}\\
&\quad -36\cdot q_0^{1}\cdot q_1^{1}\cdot q_2^{2}\cdot q_3^{2}-36\cdot q_0^{2}\cdot q_1^{2}\cdot q_2^{1}\cdot q_3^{1}\end{align*}
\begin{align*}(T_{h}^3)_{5,3}&=\frac{1}{2^{1}}\cdot q_0^{1}\cdot q_3^{5}-\frac{1}{2^{1}}\cdot q_0^{5}\cdot q_3^{1}+\frac{1}{2^{1}}\cdot q_1^{1}\cdot q_2^{5}\\
&\quad -\frac{1}{2^{1}}\cdot q_1^{5}\cdot q_2^{1}-4\cdot q_0^{1}\cdot q_1^{2}\cdot q_3^{3}+\frac{3}{2^{1}}\cdot q_0^{1}\cdot q_1^{4}\cdot q_3^{1}\\
&\quad -1\cdot q_0^{1}\cdot q_2^{2}\cdot q_3^{3}-\frac{3}{2^{1}}\cdot q_0^{1}\cdot q_2^{4}\cdot q_3^{1}-4\cdot q_0^{2}\cdot q_1^{1}\cdot q_2^{3}\\
&\quad +1\cdot q_0^{2}\cdot q_1^{3}\cdot q_2^{1}+1\cdot q_0^{3}\cdot q_1^{2}\cdot q_3^{1}+4\cdot q_0^{3}\cdot q_2^{2}\cdot q_3^{1}\\
&\quad +\frac{3}{2^{1}}\cdot q_0^{4}\cdot q_1^{1}\cdot q_2^{1}-\frac{3}{2^{1}}\cdot q_1^{1}\cdot q_2^{1}\cdot q_3^{4}\\
&\quad -1\cdot q_1^{1}\cdot q_2^{3}\cdot q_3^{2}+4\cdot q_1^{3}\cdot q_2^{1}\cdot q_3^{2}\end{align*}
\begin{align*}(T_{h}^3)_{5,4}&=-1\cdot q_0^{1}\cdot q_2^{5}+4\cdot q_0^{3}\cdot q_2^{3}-1\cdot q_0^{5}\cdot q_2^{1}+1\cdot q_1^{1}\cdot q_3^{5}\\
&\quad -4\cdot q_1^{3}\cdot q_3^{3}+1\cdot q_1^{5}\cdot q_3^{1}-4\cdot q_0^{1}\cdot q_1^{2}\cdot q_2^{3}+3\cdot q_0^{1}\cdot q_1^{4}\cdot q_2^{1}\\
&\quad +3\cdot q_0^{1}\cdot q_2^{1}\cdot q_3^{4}+2\cdot q_0^{1}\cdot q_2^{3}\cdot q_3^{2}+4\cdot q_0^{2}\cdot q_1^{1}\cdot q_3^{3}\\
&\quad -2\cdot q_0^{2}\cdot q_1^{3}\cdot q_3^{1}+2\cdot q_0^{3}\cdot q_1^{2}\cdot q_2^{1}-4\cdot q_0^{3}\cdot q_2^{1}\cdot q_3^{2}\\
&\quad -3\cdot q_0^{4}\cdot q_1^{1}\cdot q_3^{1}-2\cdot q_1^{1}\cdot q_2^{2}\cdot q_3^{3}-3\cdot q_1^{1}\cdot q_2^{4}\cdot q_3^{1}\\
&\quad +4\cdot q_1^{3}\cdot q_2^{2}\cdot q_3^{1}-12\cdot q_0^{1}\cdot q_1^{2}\cdot q_2^{1}\cdot q_3^{2}+12\cdot q_0^{2}\cdot q_1^{1}\cdot q_2^{2}\cdot q_3^{1}\end{align*}
\begin{align*}(T_{h}^3)_{5,5}&=1\cdot q_0^{6}-1\cdot q_1^{6}-1\cdot q_2^{6}+1\cdot q_3^{6}-1\cdot q_0^{2}\cdot q_1^{4}+14\cdot q_0^{2}\cdot q_2^{4}\\
&\quad -2\cdot q_0^{2}\cdot q_3^{4}+1\cdot q_0^{4}\cdot q_1^{2}-14\cdot q_0^{4}\cdot q_2^{2}-2\cdot q_0^{4}\cdot q_3^{2}\\
&\quad +2\cdot q_1^{2}\cdot q_2^{4}-14\cdot q_1^{2}\cdot q_3^{4}+2\cdot q_1^{4}\cdot q_2^{2}+14\cdot q_1^{4}\cdot q_3^{2}\\
&\quad +1\cdot q_2^{2}\cdot q_3^{4}-1\cdot q_2^{4}\cdot q_3^{2}-12\cdot q_0^{2}\cdot q_1^{2}\cdot q_2^{2}+12\cdot q_0^{2}\cdot q_1^{2}\cdot q_3^{2}\\
&\quad +12\cdot q_0^{2}\cdot q_2^{2}\cdot q_3^{2}-12\cdot q_1^{2}\cdot q_2^{2}\cdot q_3^{2}\end{align*}
\begin{align*}(T_{h}^3)_{5,6}&=2\cdot q_0^{1}\cdot q_1^{5}+4\cdot q_0^{3}\cdot q_1^{3}+2\cdot q_0^{5}\cdot q_1^{1}+2\cdot q_2^{1}\cdot q_3^{5}\\
&\quad +4\cdot q_2^{3}\cdot q_3^{3}+2\cdot q_2^{5}\cdot q_3^{1}+12\cdot q_0^{1}\cdot q_1^{1}\cdot q_2^{4}+12\cdot q_0^{1}\cdot q_1^{1}\cdot q_3^{4}\\
&\quad -16\cdot q_0^{1}\cdot q_1^{3}\cdot q_2^{2}-16\cdot q_0^{1}\cdot q_1^{3}\cdot q_3^{2}-16\cdot q_0^{2}\cdot q_2^{1}\cdot q_3^{3}\\
&\quad -16\cdot q_0^{2}\cdot q_2^{3}\cdot q_3^{1}-16\cdot q_0^{3}\cdot q_1^{1}\cdot q_2^{2}-16\cdot q_0^{3}\cdot q_1^{1}\cdot q_3^{2}\\
&\quad +12\cdot q_0^{4}\cdot q_2^{1}\cdot q_3^{1}-16\cdot q_1^{2}\cdot q_2^{1}\cdot q_3^{3}-16\cdot q_1^{2}\cdot q_2^{3}\cdot q_3^{1}\\
&\quad +12\cdot q_1^{4}\cdot q_2^{1}\cdot q_3^{1}+24\cdot q_0^{1}\cdot q_1^{1}\cdot q_2^{2}\cdot q_3^{2}+24\cdot q_0^{2}\cdot q_1^{2}\cdot q_2^{1}\cdot q_3^{1}\end{align*}
\begin{align*}(T_{h}^3)_{5,7}&=-6\cdot q_0^{1}\cdot q_2^{5}+18\cdot q_0^{3}\cdot q_2^{3}-6\cdot q_0^{5}\cdot q_2^{1}+6\cdot q_1^{1}\cdot q_3^{5}\\
&\quad -18\cdot q_1^{3}\cdot q_3^{3}+6\cdot q_1^{5}\cdot q_3^{1}+18\cdot q_0^{1}\cdot q_1^{2}\cdot q_2^{3}-6\cdot q_0^{1}\cdot q_1^{4}\cdot q_2^{1}\\
&\quad -6\cdot q_0^{1}\cdot q_2^{1}\cdot q_3^{4}-12\cdot q_0^{1}\cdot q_2^{3}\cdot q_3^{2}-18\cdot q_0^{2}\cdot q_1^{1}\cdot q_3^{3}\\
&\quad +12\cdot q_0^{2}\cdot q_1^{3}\cdot q_3^{1}-12\cdot q_0^{3}\cdot q_1^{2}\cdot q_2^{1}+18\cdot q_0^{3}\cdot q_2^{1}\cdot q_3^{2}\\
&\quad +6\cdot q_0^{4}\cdot q_1^{1}\cdot q_3^{1}+12\cdot q_1^{1}\cdot q_2^{2}\cdot q_3^{3}+6\cdot q_1^{1}\cdot q_2^{4}\cdot q_3^{1}\\
&\quad -18\cdot q_1^{3}\cdot q_2^{2}\cdot q_3^{1}+18\cdot q_0^{1}\cdot q_1^{2}\cdot q_2^{1}\cdot q_3^{2}-18\cdot q_0^{2}\cdot q_1^{1}\cdot q_2^{2}\cdot q_3^{1}\end{align*}
\begin{align*}(T_{h}^3)_{6,1}&=6\cdot q_0^{1}\cdot q_1^{1}\cdot q_2^{4}+6\cdot q_0^{1}\cdot q_1^{1}\cdot q_3^{4}-12\cdot q_0^{1}\cdot q_1^{3}\cdot q_2^{2}\\
&\quad +12\cdot q_0^{1}\cdot q_1^{3}\cdot q_3^{2}+12\cdot q_0^{2}\cdot q_2^{1}\cdot q_3^{3}-12\cdot q_0^{2}\cdot q_2^{3}\cdot q_3^{1}\\
&\quad +12\cdot q_0^{3}\cdot q_1^{1}\cdot q_2^{2}-12\cdot q_0^{3}\cdot q_1^{1}\cdot q_3^{2}-6\cdot q_0^{4}\cdot q_2^{1}\cdot q_3^{1}\\
&\quad -12\cdot q_1^{2}\cdot q_2^{1}\cdot q_3^{3}+12\cdot q_1^{2}\cdot q_2^{3}\cdot q_3^{1}-6\cdot q_1^{4}\cdot q_2^{1}\cdot q_3^{1}\\
&\quad -36\cdot q_0^{1}\cdot q_1^{1}\cdot q_2^{2}\cdot q_3^{2}+36\cdot q_0^{2}\cdot q_1^{2}\cdot q_2^{1}\cdot q_3^{1}\end{align*}
\begin{align*}(T_{h}^3)_{6,2}&=-3\cdot q_0^{2}\cdot q_2^{4}-3\cdot q_0^{2}\cdot q_3^{4}-3\cdot q_0^{4}\cdot q_2^{2}+3\cdot q_0^{4}\cdot q_3^{2}\\
&\quad +3\cdot q_1^{2}\cdot q_2^{4}+3\cdot q_1^{2}\cdot q_3^{4}-3\cdot q_1^{4}\cdot q_2^{2}+3\cdot q_1^{4}\cdot q_3^{2}\\
&\quad +18\cdot q_0^{2}\cdot q_1^{2}\cdot q_2^{2}-18\cdot q_0^{2}\cdot q_1^{2}\cdot q_3^{2}+18\cdot q_0^{2}\cdot q_2^{2}\cdot q_3^{2}\\
&\quad -18\cdot q_1^{2}\cdot q_2^{2}\cdot q_3^{2}+24\cdot q_0^{1}\cdot q_1^{1}\cdot q_2^{1}\cdot q_3^{3}-24\cdot q_0^{1}\cdot q_1^{1}\cdot q_2^{3}\cdot q_3^{1}\\
&\quad +24\cdot q_0^{1}\cdot q_1^{3}\cdot q_2^{1}\cdot q_3^{1}-24\cdot q_0^{3}\cdot q_1^{1}\cdot q_2^{1}\cdot q_3^{1}\end{align*}
\begin{align*}(T_{h}^3)_{6,3}&=-\frac{1}{2^{1}}\cdot q_0^{1}\cdot q_2^{5}+\frac{1}{2^{1}}\cdot q_0^{5}\cdot q_2^{1}+\frac{1}{2^{1}}\cdot q_1^{1}\cdot q_3^{5}\\
&\quad -\frac{1}{2^{1}}\cdot q_1^{5}\cdot q_3^{1}+4\cdot q_0^{1}\cdot q_1^{2}\cdot q_2^{3}-\frac{3}{2^{1}}\cdot q_0^{1}\cdot q_1^{4}\cdot q_2^{1}\\
&\quad +\frac{3}{2^{1}}\cdot q_0^{1}\cdot q_2^{1}\cdot q_3^{4}+1\cdot q_0^{1}\cdot q_2^{3}\cdot q_3^{2}-4\cdot q_0^{2}\cdot q_1^{1}\cdot q_3^{3}\\
&\quad +1\cdot q_0^{2}\cdot q_1^{3}\cdot q_3^{1}-1\cdot q_0^{3}\cdot q_1^{2}\cdot q_2^{1}-4\cdot q_0^{3}\cdot q_2^{1}\cdot q_3^{2}\\
&\quad +\frac{3}{2^{1}}\cdot q_0^{4}\cdot q_1^{1}\cdot q_3^{1}-1\cdot q_1^{1}\cdot q_2^{2}\cdot q_3^{3}-\frac{3}{2^{1}}\cdot q_1^{1}\cdot q_2^{4}\cdot q_3^{1}\\
&\quad +4\cdot q_1^{3}\cdot q_2^{2}\cdot q_3^{1}\end{align*}
\begin{align*}(T_{h}^3)_{6,4}&=-1\cdot q_0^{1}\cdot q_3^{5}+4\cdot q_0^{3}\cdot q_3^{3}-1\cdot q_0^{5}\cdot q_3^{1}-1\cdot q_1^{1}\cdot q_2^{5}\\
&\quad +4\cdot q_1^{3}\cdot q_2^{3}-1\cdot q_1^{5}\cdot q_2^{1}-4\cdot q_0^{1}\cdot q_1^{2}\cdot q_3^{3}+3\cdot q_0^{1}\cdot q_1^{4}\cdot q_3^{1}\\
&\quad +2\cdot q_0^{1}\cdot q_2^{2}\cdot q_3^{3}+3\cdot q_0^{1}\cdot q_2^{4}\cdot q_3^{1}-4\cdot q_0^{2}\cdot q_1^{1}\cdot q_2^{3}\\
&\quad +2\cdot q_0^{2}\cdot q_1^{3}\cdot q_2^{1}+2\cdot q_0^{3}\cdot q_1^{2}\cdot q_3^{1}-4\cdot q_0^{3}\cdot q_2^{2}\cdot q_3^{1}\\
&\quad +3\cdot q_0^{4}\cdot q_1^{1}\cdot q_2^{1}+3\cdot q_1^{1}\cdot q_2^{1}\cdot q_3^{4}+2\cdot q_1^{1}\cdot q_2^{3}\cdot q_3^{2}\\
&\quad -4\cdot q_1^{3}\cdot q_2^{1}\cdot q_3^{2}-12\cdot q_0^{1}\cdot q_1^{2}\cdot q_2^{2}\cdot q_3^{1}-12\cdot q_0^{2}\cdot q_1^{1}\cdot q_2^{1}\cdot q_3^{2}\end{align*}
\begin{align*}(T_{h}^3)_{6,5}&=-2\cdot q_0^{1}\cdot q_1^{5}-4\cdot q_0^{3}\cdot q_1^{3}-2\cdot q_0^{5}\cdot q_1^{1}+2\cdot q_2^{1}\cdot q_3^{5}\\
&\quad +4\cdot q_2^{3}\cdot q_3^{3}+2\cdot q_2^{5}\cdot q_3^{1}-12\cdot q_0^{1}\cdot q_1^{1}\cdot q_2^{4}-12\cdot q_0^{1}\cdot q_1^{1}\cdot q_3^{4}\\
&\quad +16\cdot q_0^{1}\cdot q_1^{3}\cdot q_2^{2}+16\cdot q_0^{1}\cdot q_1^{3}\cdot q_3^{2}-16\cdot q_0^{2}\cdot q_2^{1}\cdot q_3^{3}\\
&\quad -16\cdot q_0^{2}\cdot q_2^{3}\cdot q_3^{1}+16\cdot q_0^{3}\cdot q_1^{1}\cdot q_2^{2}+16\cdot q_0^{3}\cdot q_1^{1}\cdot q_3^{2}\\
&\quad +12\cdot q_0^{4}\cdot q_2^{1}\cdot q_3^{1}-16\cdot q_1^{2}\cdot q_2^{1}\cdot q_3^{3}-16\cdot q_1^{2}\cdot q_2^{3}\cdot q_3^{1}\\
&\quad +12\cdot q_1^{4}\cdot q_2^{1}\cdot q_3^{1}-24\cdot q_0^{1}\cdot q_1^{1}\cdot q_2^{2}\cdot q_3^{2}+24\cdot q_0^{2}\cdot q_1^{2}\cdot q_2^{1}\cdot q_3^{1}\end{align*}
\begin{align*}(T_{h}^3)_{6,6}&=1\cdot q_0^{6}-1\cdot q_1^{6}+1\cdot q_2^{6}-1\cdot q_3^{6}-1\cdot q_0^{2}\cdot q_1^{4}-2\cdot q_0^{2}\cdot q_2^{4}\\
&\quad +14\cdot q_0^{2}\cdot q_3^{4}+1\cdot q_0^{4}\cdot q_1^{2}-2\cdot q_0^{4}\cdot q_2^{2}-14\cdot q_0^{4}\cdot q_3^{2}\\
&\quad -14\cdot q_1^{2}\cdot q_2^{4}+2\cdot q_1^{2}\cdot q_3^{4}+14\cdot q_1^{4}\cdot q_2^{2}+2\cdot q_1^{4}\cdot q_3^{2}\\
&\quad -1\cdot q_2^{2}\cdot q_3^{4}+1\cdot q_2^{4}\cdot q_3^{2}+12\cdot q_0^{2}\cdot q_1^{2}\cdot q_2^{2}-12\cdot q_0^{2}\cdot q_1^{2}\cdot q_3^{2}\\
&\quad +12\cdot q_0^{2}\cdot q_2^{2}\cdot q_3^{2}-12\cdot q_1^{2}\cdot q_2^{2}\cdot q_3^{2}\end{align*}
\begin{align*}(T_{h}^3)_{6,7}&=6\cdot q_0^{1}\cdot q_3^{5}-18\cdot q_0^{3}\cdot q_3^{3}+6\cdot q_0^{5}\cdot q_3^{1}+6\cdot q_1^{1}\cdot q_2^{5}\\
&\quad -18\cdot q_1^{3}\cdot q_2^{3}+6\cdot q_1^{5}\cdot q_2^{1}-18\cdot q_0^{1}\cdot q_1^{2}\cdot q_3^{3}+6\cdot q_0^{1}\cdot q_1^{4}\cdot q_3^{1}\\
&\quad +12\cdot q_0^{1}\cdot q_2^{2}\cdot q_3^{3}+6\cdot q_0^{1}\cdot q_2^{4}\cdot q_3^{1}-18\cdot q_0^{2}\cdot q_1^{1}\cdot q_2^{3}\\
&\quad +12\cdot q_0^{2}\cdot q_1^{3}\cdot q_2^{1}+12\cdot q_0^{3}\cdot q_1^{2}\cdot q_3^{1}-18\cdot q_0^{3}\cdot q_2^{2}\cdot q_3^{1}\\
&\quad +6\cdot q_0^{4}\cdot q_1^{1}\cdot q_2^{1}+6\cdot q_1^{1}\cdot q_2^{1}\cdot q_3^{4}+12\cdot q_1^{1}\cdot q_2^{3}\cdot q_3^{2}\\
&\quad -18\cdot q_1^{3}\cdot q_2^{1}\cdot q_3^{2}-18\cdot q_0^{1}\cdot q_1^{2}\cdot q_2^{2}\cdot q_3^{1}-18\cdot q_0^{2}\cdot q_1^{1}\cdot q_2^{1}\cdot q_3^{2}\end{align*}
\begin{align*}(T_{h}^3)_{7,1}&=-4\cdot q_0^{3}\cdot q_2^{3}-4\cdot q_1^{3}\cdot q_3^{3}+12\cdot q_0^{1}\cdot q_1^{2}\cdot q_2^{3}+12\cdot q_0^{2}\cdot q_1^{1}\cdot q_3^{3}\\
&\quad +12\cdot q_0^{3}\cdot q_2^{1}\cdot q_3^{2}+12\cdot q_1^{3}\cdot q_2^{2}\cdot q_3^{1}-36\cdot q_0^{1}\cdot q_1^{2}\cdot q_2^{1}\cdot q_3^{2}\\
&\quad -36\cdot q_0^{2}\cdot q_1^{1}\cdot q_2^{2}\cdot q_3^{1}\end{align*}
\begin{align*}(T_{h}^3)_{7,2}&=-4\cdot q_0^{3}\cdot q_3^{3}+4\cdot q_1^{3}\cdot q_2^{3}+12\cdot q_0^{1}\cdot q_1^{2}\cdot q_3^{3}-12\cdot q_0^{2}\cdot q_1^{1}\cdot q_2^{3}\\
&\quad +12\cdot q_0^{3}\cdot q_2^{2}\cdot q_3^{1}-12\cdot q_1^{3}\cdot q_2^{1}\cdot q_3^{2}-36\cdot q_0^{1}\cdot q_1^{2}\cdot q_2^{2}\cdot q_3^{1}\\
&\quad +36\cdot q_0^{2}\cdot q_1^{1}\cdot q_2^{1}\cdot q_3^{2}\end{align*}
\begin{align*}(T_{h}^3)_{7,3}&=-2\cdot q_0^{1}\cdot q_1^{1}\cdot q_2^{4}+2\cdot q_0^{1}\cdot q_1^{1}\cdot q_3^{4}+2\cdot q_0^{1}\cdot q_1^{3}\cdot q_2^{2}\\
&\quad -2\cdot q_0^{1}\cdot q_1^{3}\cdot q_3^{2}+2\cdot q_0^{2}\cdot q_2^{1}\cdot q_3^{3}+2\cdot q_0^{2}\cdot q_2^{3}\cdot q_3^{1}\\
&\quad +2\cdot q_0^{3}\cdot q_1^{1}\cdot q_2^{2}-2\cdot q_0^{3}\cdot q_1^{1}\cdot q_3^{2}-2\cdot q_0^{4}\cdot q_2^{1}\cdot q_3^{1}\\
&\quad -2\cdot q_1^{2}\cdot q_2^{1}\cdot q_3^{3}-2\cdot q_1^{2}\cdot q_2^{3}\cdot q_3^{1}+2\cdot q_1^{4}\cdot q_2^{1}\cdot q_3^{1}\end{align*}
\begin{align*}(T_{h}^3)_{7,4}&=2\cdot q_0^{2}\cdot q_2^{4}-2\cdot q_0^{2}\cdot q_3^{4}-2\cdot q_0^{4}\cdot q_2^{2}+2\cdot q_0^{4}\cdot q_3^{2}\\
&\quad -2\cdot q_1^{2}\cdot q_2^{4}+2\cdot q_1^{2}\cdot q_3^{4}+2\cdot q_1^{4}\cdot q_2^{2}-2\cdot q_1^{4}\cdot q_3^{2}\\
&\quad +8\cdot q_0^{1}\cdot q_1^{1}\cdot q_2^{1}\cdot q_3^{3}+8\cdot q_0^{1}\cdot q_1^{1}\cdot q_2^{3}\cdot q_3^{1}\\
&\quad -8\cdot q_0^{1}\cdot q_1^{3}\cdot q_2^{1}\cdot q_3^{1}-8\cdot q_0^{3}\cdot q_1^{1}\cdot q_2^{1}\cdot q_3^{1}\end{align*}
\begin{align*}(T_{h}^3)_{7,5}&=4\cdot q_0^{1}\cdot q_2^{5}-12\cdot q_0^{3}\cdot q_2^{3}+4\cdot q_0^{5}\cdot q_2^{1}+4\cdot q_1^{1}\cdot q_3^{5}\\
&\quad -12\cdot q_1^{3}\cdot q_3^{3}+4\cdot q_1^{5}\cdot q_3^{1}-12\cdot q_0^{1}\cdot q_1^{2}\cdot q_2^{3}+4\cdot q_0^{1}\cdot q_1^{4}\cdot q_2^{1}\\
&\quad +4\cdot q_0^{1}\cdot q_2^{1}\cdot q_3^{4}+8\cdot q_0^{1}\cdot q_2^{3}\cdot q_3^{2}-12\cdot q_0^{2}\cdot q_1^{1}\cdot q_3^{3}\\
&\quad +8\cdot q_0^{2}\cdot q_1^{3}\cdot q_3^{1}+8\cdot q_0^{3}\cdot q_1^{2}\cdot q_2^{1}-12\cdot q_0^{3}\cdot q_2^{1}\cdot q_3^{2}\\
&\quad +4\cdot q_0^{4}\cdot q_1^{1}\cdot q_3^{1}+8\cdot q_1^{1}\cdot q_2^{2}\cdot q_3^{3}+4\cdot q_1^{1}\cdot q_2^{4}\cdot q_3^{1}\\
&\quad -12\cdot q_1^{3}\cdot q_2^{2}\cdot q_3^{1}-12\cdot q_0^{1}\cdot q_1^{2}\cdot q_2^{1}\cdot q_3^{2}-12\cdot q_0^{2}\cdot q_1^{1}\cdot q_2^{2}\cdot q_3^{1}\end{align*}
\begin{align*}(T_{h}^3)_{7,6}&=-4\cdot q_0^{1}\cdot q_3^{5}+12\cdot q_0^{3}\cdot q_3^{3}-4\cdot q_0^{5}\cdot q_3^{1}+4\cdot q_1^{1}\cdot q_2^{5}\\
&\quad -12\cdot q_1^{3}\cdot q_2^{3}+4\cdot q_1^{5}\cdot q_2^{1}+12\cdot q_0^{1}\cdot q_1^{2}\cdot q_3^{3}-4\cdot q_0^{1}\cdot q_1^{4}\cdot q_3^{1}\\
&\quad -8\cdot q_0^{1}\cdot q_2^{2}\cdot q_3^{3}-4\cdot q_0^{1}\cdot q_2^{4}\cdot q_3^{1}-12\cdot q_0^{2}\cdot q_1^{1}\cdot q_2^{3}\\
&\quad +8\cdot q_0^{2}\cdot q_1^{3}\cdot q_2^{1}-8\cdot q_0^{3}\cdot q_1^{2}\cdot q_3^{1}+12\cdot q_0^{3}\cdot q_2^{2}\cdot q_3^{1}\\
&\quad +4\cdot q_0^{4}\cdot q_1^{1}\cdot q_2^{1}+4\cdot q_1^{1}\cdot q_2^{1}\cdot q_3^{4}+8\cdot q_1^{1}\cdot q_2^{3}\cdot q_3^{2}\\
&\quad -12\cdot q_1^{3}\cdot q_2^{1}\cdot q_3^{2}+12\cdot q_0^{1}\cdot q_1^{2}\cdot q_2^{2}\cdot q_3^{1}-12\cdot q_0^{2}\cdot q_1^{1}\cdot q_2^{1}\cdot q_3^{2}\end{align*}
\begin{align*}(T_{h}^3)_{7,7}&=1\cdot q_0^{6}+1\cdot q_1^{6}-1\cdot q_2^{6}-1\cdot q_3^{6}+3\cdot q_0^{2}\cdot q_1^{4}+9\cdot q_0^{2}\cdot q_2^{4}\\
&\quad +9\cdot q_0^{2}\cdot q_3^{4}+3\cdot q_0^{4}\cdot q_1^{2}-9\cdot q_0^{4}\cdot q_2^{2}-9\cdot q_0^{4}\cdot q_3^{2}\\
&\quad +9\cdot q_1^{2}\cdot q_2^{4}+9\cdot q_1^{2}\cdot q_3^{4}-9\cdot q_1^{4}\cdot q_2^{2}-9\cdot q_1^{4}\cdot q_3^{2}\\
&\quad -3\cdot q_2^{2}\cdot q_3^{4}-3\cdot q_2^{4}\cdot q_3^{2}-18\cdot q_0^{2}\cdot q_1^{2}\cdot q_2^{2}-18\cdot q_0^{2}\cdot q_1^{2}\cdot q_3^{2}\\
&\quad +18\cdot q_0^{2}\cdot q_2^{2}\cdot q_3^{2}+18\cdot q_1^{2}\cdot q_2^{2}\cdot q_3^{2}\end{align*}
\section{Term $\ell=4$}
\subsection{Matrix $D_4$}
We will write the matrix $D^{-1}$ instead of $D$. The matrices $D$ and $D^{-1}$ are diagonal, so one is easy to construct from the other. The matrix $D^{-1}$ is more convenient becauce the denominators of diagonal terms are powers of $2$. Therefore, their binary expressions are finite and there is no rounding in computer storage.\[D_4^{-1}=\left[\begin{array}{ccccccccc}
\frac{315}{2^{2}} & 0 & 0 & 0 & 0 & 0 & 0 & 0 & 0\\
0 & \frac{315}{2^{6}} & 0 & 0 & 0 & 0 & 0 & 0 & 0\\
0 & 0 & \frac{315}{2^{3}} & 0 & 0 & 0 & 0 & 0 & 0\\
0 & 0 & 0 & \frac{315}{2^{3}} & 0 & 0 & 0 & 0 & 0\\
0 & 0 & 0 & 0 & \frac{45}{2^{2}} & 0 & 0 & 0 & 0\\
0 & 0 & 0 & 0 & 0 & \frac{45}{2^{4}} & 0 & 0 & 0\\
0 & 0 & 0 & 0 & 0 & 0 & \frac{45}{2^{3}} & 0 & 0\\
0 & 0 & 0 & 0 & 0 & 0 & 0 & \frac{45}{2^{3}} & 0\\
0 & 0 & 0 & 0 & 0 & 0 & 0 & 0 & \frac{9}{2^{6}}
\end{array}\right].\]
\subsection{Matrix $C_4$}
\[C_4=\left[\begin{array}{ccccccccc}
\frac{13}{2^{20}} & 0 & 0 & 0 & 0 & 0 & 0 & 0 & 0\\
0 & \frac{13569}{2^{18}} & 0 & 0 & 0 & -\frac{5593}{2^{17}} & 0 & 0 & -\frac{797}{2^{15}}\\
0 & 0 & \frac{27}{2^{20}} & 0 & 0 & 0 & 0 & 0 & 0\\
0 & 0 & 0 & \frac{27}{2^{20}} & 0 & 0 & 0 & 0 & 0\\
0 & 0 & 0 & 0 & \frac{93}{2^{20}} & 0 & 0 & 0 & 0\\
0 & -\frac{5593}{2^{17}} & 0 & 0 & 0 & \frac{37405}{2^{20}} & 0 & 0 & \frac{5277}{2^{18}}\\
0 & 0 & 0 & 0 & 0 & 0 & \frac{93}{2^{19}} & 0 & 0\\
0 & 0 & 0 & 0 & 0 & 0 & 0 & \frac{93}{2^{19}} & 0\\
0 & -\frac{797}{2^{15}} & 0 & 0 & 0 & \frac{5277}{2^{18}} & 0 & 0 & \frac{609}{2^{15}}
\end{array}\right].\]
\subsection{Matrix $T_{h}^4(q_0,q_1,q_2,q_3)$}
\begin{align*}(T_{h}^4)_{1,1}&=1\cdot q_0^{8}+1\cdot q_1^{8}-1\cdot q_2^{8}-1\cdot q_3^{8}-28\cdot q_0^{2}\cdot q_1^{6}+70\cdot q_0^{4}\cdot q_1^{4}\\
&\quad -28\cdot q_0^{6}\cdot q_1^{2}+28\cdot q_2^{2}\cdot q_3^{6}-70\cdot q_2^{4}\cdot q_3^{4}+28\cdot q_2^{6}\cdot q_3^{2}\end{align*}
\begin{align*}(T_{h}^4)_{1,2}&=-32\cdot q_0^{1}\cdot q_1^{7}+224\cdot q_0^{3}\cdot q_1^{5}-224\cdot q_0^{5}\cdot q_1^{3}+32\cdot q_0^{7}\cdot q_1^{1}\\
&\quad -32\cdot q_2^{1}\cdot q_3^{7}+224\cdot q_2^{3}\cdot q_3^{5}-224\cdot q_2^{5}\cdot q_3^{3}+32\cdot q_2^{7}\cdot q_3^{1}\end{align*}
\begin{align*}(T_{h}^4)_{1,3}&=4\cdot q_0^{1}\cdot q_3^{7}+4\cdot q_0^{7}\cdot q_3^{1}-4\cdot q_1^{1}\cdot q_2^{7}-4\cdot q_1^{7}\cdot q_2^{1}\\
&\quad -28\cdot q_0^{1}\cdot q_1^{6}\cdot q_3^{1}-84\cdot q_0^{1}\cdot q_2^{2}\cdot q_3^{5}+140\cdot q_0^{1}\cdot q_2^{4}\cdot q_3^{3}\\
&\quad -28\cdot q_0^{1}\cdot q_2^{6}\cdot q_3^{1}+84\cdot q_0^{2}\cdot q_1^{5}\cdot q_2^{1}+140\cdot q_0^{3}\cdot q_1^{4}\cdot q_3^{1}\\
&\quad -140\cdot q_0^{4}\cdot q_1^{3}\cdot q_2^{1}-84\cdot q_0^{5}\cdot q_1^{2}\cdot q_3^{1}+28\cdot q_0^{6}\cdot q_1^{1}\cdot q_2^{1}\\
&\quad +28\cdot q_1^{1}\cdot q_2^{1}\cdot q_3^{6}-140\cdot q_1^{1}\cdot q_2^{3}\cdot q_3^{4}+84\cdot q_1^{1}\cdot q_2^{5}\cdot q_3^{2}\end{align*}
\begin{align*}(T_{h}^4)_{1,4}&=-4\cdot q_0^{1}\cdot q_2^{7}-4\cdot q_0^{7}\cdot q_2^{1}-4\cdot q_1^{1}\cdot q_3^{7}-4\cdot q_1^{7}\cdot q_3^{1}\\
&\quad +28\cdot q_0^{1}\cdot q_1^{6}\cdot q_2^{1}+28\cdot q_0^{1}\cdot q_2^{1}\cdot q_3^{6}-140\cdot q_0^{1}\cdot q_2^{3}\cdot q_3^{4}\\
&\quad +84\cdot q_0^{1}\cdot q_2^{5}\cdot q_3^{2}+84\cdot q_0^{2}\cdot q_1^{5}\cdot q_3^{1}-140\cdot q_0^{3}\cdot q_1^{4}\cdot q_2^{1}\\
&\quad -140\cdot q_0^{4}\cdot q_1^{3}\cdot q_3^{1}+84\cdot q_0^{5}\cdot q_1^{2}\cdot q_2^{1}+28\cdot q_0^{6}\cdot q_1^{1}\cdot q_3^{1}\\
&\quad +84\cdot q_1^{1}\cdot q_2^{2}\cdot q_3^{5}-140\cdot q_1^{1}\cdot q_2^{4}\cdot q_3^{3}+28\cdot q_1^{1}\cdot q_2^{6}\cdot q_3^{1}\end{align*}
\begin{align*}(T_{h}^4)_{1,5}&=14\cdot q_0^{2}\cdot q_2^{6}-14\cdot q_0^{2}\cdot q_3^{6}-14\cdot q_0^{6}\cdot q_2^{2}+14\cdot q_0^{6}\cdot q_3^{2}\\
&\quad -14\cdot q_1^{2}\cdot q_2^{6}+14\cdot q_1^{2}\cdot q_3^{6}+14\cdot q_1^{6}\cdot q_2^{2}-14\cdot q_1^{6}\cdot q_3^{2}\\
&\quad -210\cdot q_0^{2}\cdot q_1^{4}\cdot q_2^{2}+210\cdot q_0^{2}\cdot q_1^{4}\cdot q_3^{2}+210\cdot q_0^{2}\cdot q_2^{2}\cdot q_3^{4}\\
&\quad -210\cdot q_0^{2}\cdot q_2^{4}\cdot q_3^{2}+210\cdot q_0^{4}\cdot q_1^{2}\cdot q_2^{2}-210\cdot q_0^{4}\cdot q_1^{2}\cdot q_3^{2}\\
&\quad -210\cdot q_1^{2}\cdot q_2^{2}\cdot q_3^{4}+210\cdot q_1^{2}\cdot q_2^{4}\cdot q_3^{2}-168\cdot q_0^{1}\cdot q_1^{1}\cdot q_2^{1}\cdot q_3^{5}\\
&\quad +560\cdot q_0^{1}\cdot q_1^{1}\cdot q_2^{3}\cdot q_3^{3}-168\cdot q_0^{1}\cdot q_1^{1}\cdot q_2^{5}\cdot q_3^{1}\\
&\quad +168\cdot q_0^{1}\cdot q_1^{5}\cdot q_2^{1}\cdot q_3^{1}-560\cdot q_0^{3}\cdot q_1^{3}\cdot q_2^{1}\cdot q_3^{1}\\
&\quad +168\cdot q_0^{5}\cdot q_1^{1}\cdot q_2^{1}\cdot q_3^{1}\end{align*}
\begin{align*}(T_{h}^4)_{1,6}&=-56\cdot q_0^{1}\cdot q_1^{1}\cdot q_2^{6}+56\cdot q_0^{1}\cdot q_1^{1}\cdot q_3^{6}-168\cdot q_0^{1}\cdot q_1^{5}\cdot q_2^{2}\\
&\quad +168\cdot q_0^{1}\cdot q_1^{5}\cdot q_3^{2}-168\cdot q_0^{2}\cdot q_2^{1}\cdot q_3^{5}+560\cdot q_0^{2}\cdot q_2^{3}\cdot q_3^{3}\\
&\quad -168\cdot q_0^{2}\cdot q_2^{5}\cdot q_3^{1}+560\cdot q_0^{3}\cdot q_1^{3}\cdot q_2^{2}-560\cdot q_0^{3}\cdot q_1^{3}\cdot q_3^{2}\\
&\quad -168\cdot q_0^{5}\cdot q_1^{1}\cdot q_2^{2}+168\cdot q_0^{5}\cdot q_1^{1}\cdot q_3^{2}-56\cdot q_0^{6}\cdot q_2^{1}\cdot q_3^{1}\\
&\quad +168\cdot q_1^{2}\cdot q_2^{1}\cdot q_3^{5}-560\cdot q_1^{2}\cdot q_2^{3}\cdot q_3^{3}+168\cdot q_1^{2}\cdot q_2^{5}\cdot q_3^{1}\\
&\quad +56\cdot q_1^{6}\cdot q_2^{1}\cdot q_3^{1}-840\cdot q_0^{1}\cdot q_1^{1}\cdot q_2^{2}\cdot q_3^{4}+840\cdot q_0^{1}\cdot q_1^{1}\cdot q_2^{4}\cdot q_3^{2}\\
&\quad -840\cdot q_0^{2}\cdot q_1^{4}\cdot q_2^{1}\cdot q_3^{1}+840\cdot q_0^{4}\cdot q_1^{2}\cdot q_2^{1}\cdot q_3^{1}\end{align*}
\begin{align*}(T_{h}^4)_{1,7}&=28\cdot q_0^{3}\cdot q_3^{5}+28\cdot q_0^{5}\cdot q_3^{3}-28\cdot q_1^{3}\cdot q_2^{5}-28\cdot q_1^{5}\cdot q_2^{3}\\
&\quad -84\cdot q_0^{1}\cdot q_1^{2}\cdot q_3^{5}+140\cdot q_0^{1}\cdot q_1^{4}\cdot q_3^{3}+84\cdot q_0^{2}\cdot q_1^{1}\cdot q_2^{5}\\
&\quad +280\cdot q_0^{2}\cdot q_1^{3}\cdot q_2^{3}-280\cdot q_0^{3}\cdot q_1^{2}\cdot q_3^{3}-280\cdot q_0^{3}\cdot q_2^{2}\cdot q_3^{3}\\
&\quad +140\cdot q_0^{3}\cdot q_2^{4}\cdot q_3^{1}-140\cdot q_0^{4}\cdot q_1^{1}\cdot q_2^{3}-84\cdot q_0^{5}\cdot q_2^{2}\cdot q_3^{1}\\
&\quad -140\cdot q_1^{3}\cdot q_2^{1}\cdot q_3^{4}+280\cdot q_1^{3}\cdot q_2^{3}\cdot q_3^{2}+84\cdot q_1^{5}\cdot q_2^{1}\cdot q_3^{2}\\
&\quad +840\cdot q_0^{1}\cdot q_1^{2}\cdot q_2^{2}\cdot q_3^{3}-420\cdot q_0^{1}\cdot q_1^{2}\cdot q_2^{4}\cdot q_3^{1}\\
&\quad -420\cdot q_0^{1}\cdot q_1^{4}\cdot q_2^{2}\cdot q_3^{1}+420\cdot q_0^{2}\cdot q_1^{1}\cdot q_2^{1}\cdot q_3^{4}\\
&\quad -840\cdot q_0^{2}\cdot q_1^{1}\cdot q_2^{3}\cdot q_3^{2}-840\cdot q_0^{2}\cdot q_1^{3}\cdot q_2^{1}\cdot q_3^{2}\\
&\quad +840\cdot q_0^{3}\cdot q_1^{2}\cdot q_2^{2}\cdot q_3^{1}+420\cdot q_0^{4}\cdot q_1^{1}\cdot q_2^{1}\cdot q_3^{2}\end{align*}
\begin{align*}(T_{h}^4)_{1,8}&=28\cdot q_0^{3}\cdot q_2^{5}+28\cdot q_0^{5}\cdot q_2^{3}+28\cdot q_1^{3}\cdot q_3^{5}+28\cdot q_1^{5}\cdot q_3^{3}\\
&\quad -84\cdot q_0^{1}\cdot q_1^{2}\cdot q_2^{5}+140\cdot q_0^{1}\cdot q_1^{4}\cdot q_2^{3}-84\cdot q_0^{2}\cdot q_1^{1}\cdot q_3^{5}\\
&\quad -280\cdot q_0^{2}\cdot q_1^{3}\cdot q_3^{3}-280\cdot q_0^{3}\cdot q_1^{2}\cdot q_2^{3}+140\cdot q_0^{3}\cdot q_2^{1}\cdot q_3^{4}\\
&\quad -280\cdot q_0^{3}\cdot q_2^{3}\cdot q_3^{2}+140\cdot q_0^{4}\cdot q_1^{1}\cdot q_3^{3}-84\cdot q_0^{5}\cdot q_2^{1}\cdot q_3^{2}\\
&\quad -280\cdot q_1^{3}\cdot q_2^{2}\cdot q_3^{3}+140\cdot q_1^{3}\cdot q_2^{4}\cdot q_3^{1}-84\cdot q_1^{5}\cdot q_2^{2}\cdot q_3^{1}\\
&\quad -420\cdot q_0^{1}\cdot q_1^{2}\cdot q_2^{1}\cdot q_3^{4}+840\cdot q_0^{1}\cdot q_1^{2}\cdot q_2^{3}\cdot q_3^{2}\\
&\quad -420\cdot q_0^{1}\cdot q_1^{4}\cdot q_2^{1}\cdot q_3^{2}+840\cdot q_0^{2}\cdot q_1^{1}\cdot q_2^{2}\cdot q_3^{3}\\
&\quad -420\cdot q_0^{2}\cdot q_1^{1}\cdot q_2^{4}\cdot q_3^{1}+840\cdot q_0^{2}\cdot q_1^{3}\cdot q_2^{2}\cdot q_3^{1}\\
&\quad +840\cdot q_0^{3}\cdot q_1^{2}\cdot q_2^{1}\cdot q_3^{2}-420\cdot q_0^{4}\cdot q_1^{1}\cdot q_2^{2}\cdot q_3^{1}\end{align*}
\begin{align*}(T_{h}^4)_{1,9}&=-1120\cdot q_0^{1}\cdot q_1^{3}\cdot q_2^{4}-1120\cdot q_0^{1}\cdot q_1^{3}\cdot q_3^{4}+1120\cdot q_0^{3}\cdot q_1^{1}\cdot q_2^{4}\\
&\quad +1120\cdot q_0^{3}\cdot q_1^{1}\cdot q_3^{4}-1120\cdot q_0^{4}\cdot q_2^{1}\cdot q_3^{3}+1120\cdot q_0^{4}\cdot q_2^{3}\cdot q_3^{1}\\
&\quad -1120\cdot q_1^{4}\cdot q_2^{1}\cdot q_3^{3}+1120\cdot q_1^{4}\cdot q_2^{3}\cdot q_3^{1}+6720\cdot q_0^{1}\cdot q_1^{3}\cdot q_2^{2}\cdot q_3^{2}\\
&\quad +6720\cdot q_0^{2}\cdot q_1^{2}\cdot q_2^{1}\cdot q_3^{3}-6720\cdot q_0^{2}\cdot q_1^{2}\cdot q_2^{3}\cdot q_3^{1}\\
&\quad -6720\cdot q_0^{3}\cdot q_1^{1}\cdot q_2^{2}\cdot q_3^{2}\end{align*}
\begin{align*}(T_{h}^4)_{2,1}&=2\cdot q_0^{1}\cdot q_1^{7}-14\cdot q_0^{3}\cdot q_1^{5}+14\cdot q_0^{5}\cdot q_1^{3}-2\cdot q_0^{7}\cdot q_1^{1}\\
&\quad -2\cdot q_2^{1}\cdot q_3^{7}+14\cdot q_2^{3}\cdot q_3^{5}-14\cdot q_2^{5}\cdot q_3^{3}+2\cdot q_2^{7}\cdot q_3^{1}\end{align*}
\begin{align*}(T_{h}^4)_{2,2}&=1\cdot q_0^{8}+1\cdot q_1^{8}+1\cdot q_2^{8}+1\cdot q_3^{8}-28\cdot q_0^{2}\cdot q_1^{6}+70\cdot q_0^{4}\cdot q_1^{4}\\
&\quad -28\cdot q_0^{6}\cdot q_1^{2}-28\cdot q_2^{2}\cdot q_3^{6}+70\cdot q_2^{4}\cdot q_3^{4}-28\cdot q_2^{6}\cdot q_3^{2}\end{align*}
\begin{align*}(T_{h}^4)_{2,3}&=-1\cdot q_0^{1}\cdot q_2^{7}+1\cdot q_0^{7}\cdot q_2^{1}-1\cdot q_1^{1}\cdot q_3^{7}+1\cdot q_1^{7}\cdot q_3^{1}\\
&\quad -7\cdot q_0^{1}\cdot q_1^{6}\cdot q_2^{1}+7\cdot q_0^{1}\cdot q_2^{1}\cdot q_3^{6}-35\cdot q_0^{1}\cdot q_2^{3}\cdot q_3^{4}\\
&\quad +21\cdot q_0^{1}\cdot q_2^{5}\cdot q_3^{2}-21\cdot q_0^{2}\cdot q_1^{5}\cdot q_3^{1}+35\cdot q_0^{3}\cdot q_1^{4}\cdot q_2^{1}\\
&\quad +35\cdot q_0^{4}\cdot q_1^{3}\cdot q_3^{1}-21\cdot q_0^{5}\cdot q_1^{2}\cdot q_2^{1}-7\cdot q_0^{6}\cdot q_1^{1}\cdot q_3^{1}\\
&\quad +21\cdot q_1^{1}\cdot q_2^{2}\cdot q_3^{5}-35\cdot q_1^{1}\cdot q_2^{4}\cdot q_3^{3}+7\cdot q_1^{1}\cdot q_2^{6}\cdot q_3^{1}\end{align*}
\begin{align*}(T_{h}^4)_{2,4}&=-1\cdot q_0^{1}\cdot q_3^{7}+1\cdot q_0^{7}\cdot q_3^{1}+1\cdot q_1^{1}\cdot q_2^{7}-1\cdot q_1^{7}\cdot q_2^{1}\\
&\quad -7\cdot q_0^{1}\cdot q_1^{6}\cdot q_3^{1}+21\cdot q_0^{1}\cdot q_2^{2}\cdot q_3^{5}-35\cdot q_0^{1}\cdot q_2^{4}\cdot q_3^{3}\\
&\quad +7\cdot q_0^{1}\cdot q_2^{6}\cdot q_3^{1}+21\cdot q_0^{2}\cdot q_1^{5}\cdot q_2^{1}+35\cdot q_0^{3}\cdot q_1^{4}\cdot q_3^{1}\\
&\quad -35\cdot q_0^{4}\cdot q_1^{3}\cdot q_2^{1}-21\cdot q_0^{5}\cdot q_1^{2}\cdot q_3^{1}+7\cdot q_0^{6}\cdot q_1^{1}\cdot q_2^{1}\\
&\quad -7\cdot q_1^{1}\cdot q_2^{1}\cdot q_3^{6}+35\cdot q_1^{1}\cdot q_2^{3}\cdot q_3^{4}-21\cdot q_1^{1}\cdot q_2^{5}\cdot q_3^{2}\end{align*}
\begin{align*}(T_{h}^4)_{2,5}&=-7\cdot q_0^{1}\cdot q_1^{1}\cdot q_2^{6}+7\cdot q_0^{1}\cdot q_1^{1}\cdot q_3^{6}+21\cdot q_0^{1}\cdot q_1^{5}\cdot q_2^{2}\\
&\quad -21\cdot q_0^{1}\cdot q_1^{5}\cdot q_3^{2}-21\cdot q_0^{2}\cdot q_2^{1}\cdot q_3^{5}+70\cdot q_0^{2}\cdot q_2^{3}\cdot q_3^{3}\\
&\quad -21\cdot q_0^{2}\cdot q_2^{5}\cdot q_3^{1}-70\cdot q_0^{3}\cdot q_1^{3}\cdot q_2^{2}+70\cdot q_0^{3}\cdot q_1^{3}\cdot q_3^{2}\\
&\quad +21\cdot q_0^{5}\cdot q_1^{1}\cdot q_2^{2}-21\cdot q_0^{5}\cdot q_1^{1}\cdot q_3^{2}+7\cdot q_0^{6}\cdot q_2^{1}\cdot q_3^{1}\\
&\quad +21\cdot q_1^{2}\cdot q_2^{1}\cdot q_3^{5}-70\cdot q_1^{2}\cdot q_2^{3}\cdot q_3^{3}+21\cdot q_1^{2}\cdot q_2^{5}\cdot q_3^{1}\\
&\quad -7\cdot q_1^{6}\cdot q_2^{1}\cdot q_3^{1}-105\cdot q_0^{1}\cdot q_1^{1}\cdot q_2^{2}\cdot q_3^{4}+105\cdot q_0^{1}\cdot q_1^{1}\cdot q_2^{4}\cdot q_3^{2}\\
&\quad +105\cdot q_0^{2}\cdot q_1^{4}\cdot q_2^{1}\cdot q_3^{1}-105\cdot q_0^{4}\cdot q_1^{2}\cdot q_2^{1}\cdot q_3^{1}\end{align*}
\begin{align*}(T_{h}^4)_{2,6}&=-7\cdot q_0^{2}\cdot q_2^{6}+7\cdot q_0^{2}\cdot q_3^{6}-7\cdot q_0^{6}\cdot q_2^{2}+7\cdot q_0^{6}\cdot q_3^{2}\\
&\quad +7\cdot q_1^{2}\cdot q_2^{6}-7\cdot q_1^{2}\cdot q_3^{6}+7\cdot q_1^{6}\cdot q_2^{2}-7\cdot q_1^{6}\cdot q_3^{2}\\
&\quad -105\cdot q_0^{2}\cdot q_1^{4}\cdot q_2^{2}+105\cdot q_0^{2}\cdot q_1^{4}\cdot q_3^{2}-105\cdot q_0^{2}\cdot q_2^{2}\cdot q_3^{4}\\
&\quad +105\cdot q_0^{2}\cdot q_2^{4}\cdot q_3^{2}+105\cdot q_0^{4}\cdot q_1^{2}\cdot q_2^{2}-105\cdot q_0^{4}\cdot q_1^{2}\cdot q_3^{2}\\
&\quad +105\cdot q_1^{2}\cdot q_2^{2}\cdot q_3^{4}-105\cdot q_1^{2}\cdot q_2^{4}\cdot q_3^{2}+84\cdot q_0^{1}\cdot q_1^{1}\cdot q_2^{1}\cdot q_3^{5}\\
&\quad -280\cdot q_0^{1}\cdot q_1^{1}\cdot q_2^{3}\cdot q_3^{3}+84\cdot q_0^{1}\cdot q_1^{1}\cdot q_2^{5}\cdot q_3^{1}\\
&\quad +84\cdot q_0^{1}\cdot q_1^{5}\cdot q_2^{1}\cdot q_3^{1}-280\cdot q_0^{3}\cdot q_1^{3}\cdot q_2^{1}\cdot q_3^{1}\\
&\quad +84\cdot q_0^{5}\cdot q_1^{1}\cdot q_2^{1}\cdot q_3^{1}\end{align*}
\begin{align*}(T_{h}^4)_{2,7}&=7\cdot q_0^{3}\cdot q_2^{5}-7\cdot q_0^{5}\cdot q_2^{3}+7\cdot q_1^{3}\cdot q_3^{5}-7\cdot q_1^{5}\cdot q_3^{3}\\
&\quad -21\cdot q_0^{1}\cdot q_1^{2}\cdot q_2^{5}-35\cdot q_0^{1}\cdot q_1^{4}\cdot q_2^{3}-21\cdot q_0^{2}\cdot q_1^{1}\cdot q_3^{5}\\
&\quad +70\cdot q_0^{2}\cdot q_1^{3}\cdot q_3^{3}+70\cdot q_0^{3}\cdot q_1^{2}\cdot q_2^{3}+35\cdot q_0^{3}\cdot q_2^{1}\cdot q_3^{4}\\
&\quad -70\cdot q_0^{3}\cdot q_2^{3}\cdot q_3^{2}-35\cdot q_0^{4}\cdot q_1^{1}\cdot q_3^{3}+21\cdot q_0^{5}\cdot q_2^{1}\cdot q_3^{2}\\
&\quad -70\cdot q_1^{3}\cdot q_2^{2}\cdot q_3^{3}+35\cdot q_1^{3}\cdot q_2^{4}\cdot q_3^{1}+21\cdot q_1^{5}\cdot q_2^{2}\cdot q_3^{1}\\
&\quad -105\cdot q_0^{1}\cdot q_1^{2}\cdot q_2^{1}\cdot q_3^{4}+210\cdot q_0^{1}\cdot q_1^{2}\cdot q_2^{3}\cdot q_3^{2}\\
&\quad +105\cdot q_0^{1}\cdot q_1^{4}\cdot q_2^{1}\cdot q_3^{2}+210\cdot q_0^{2}\cdot q_1^{1}\cdot q_2^{2}\cdot q_3^{3}\\
&\quad -105\cdot q_0^{2}\cdot q_1^{1}\cdot q_2^{4}\cdot q_3^{1}-210\cdot q_0^{2}\cdot q_1^{3}\cdot q_2^{2}\cdot q_3^{1}\\
&\quad -210\cdot q_0^{3}\cdot q_1^{2}\cdot q_2^{1}\cdot q_3^{2}+105\cdot q_0^{4}\cdot q_1^{1}\cdot q_2^{2}\cdot q_3^{1}\end{align*}
\begin{align*}(T_{h}^4)_{2,8}&=-7\cdot q_0^{3}\cdot q_3^{5}+7\cdot q_0^{5}\cdot q_3^{3}+7\cdot q_1^{3}\cdot q_2^{5}-7\cdot q_1^{5}\cdot q_2^{3}\\
&\quad +21\cdot q_0^{1}\cdot q_1^{2}\cdot q_3^{5}+35\cdot q_0^{1}\cdot q_1^{4}\cdot q_3^{3}-21\cdot q_0^{2}\cdot q_1^{1}\cdot q_2^{5}\\
&\quad +70\cdot q_0^{2}\cdot q_1^{3}\cdot q_2^{3}-70\cdot q_0^{3}\cdot q_1^{2}\cdot q_3^{3}+70\cdot q_0^{3}\cdot q_2^{2}\cdot q_3^{3}\\
&\quad -35\cdot q_0^{3}\cdot q_2^{4}\cdot q_3^{1}-35\cdot q_0^{4}\cdot q_1^{1}\cdot q_2^{3}-21\cdot q_0^{5}\cdot q_2^{2}\cdot q_3^{1}\\
&\quad +35\cdot q_1^{3}\cdot q_2^{1}\cdot q_3^{4}-70\cdot q_1^{3}\cdot q_2^{3}\cdot q_3^{2}+21\cdot q_1^{5}\cdot q_2^{1}\cdot q_3^{2}\\
&\quad -210\cdot q_0^{1}\cdot q_1^{2}\cdot q_2^{2}\cdot q_3^{3}+105\cdot q_0^{1}\cdot q_1^{2}\cdot q_2^{4}\cdot q_3^{1}\\
&\quad -105\cdot q_0^{1}\cdot q_1^{4}\cdot q_2^{2}\cdot q_3^{1}-105\cdot q_0^{2}\cdot q_1^{1}\cdot q_2^{1}\cdot q_3^{4}\\
&\quad +210\cdot q_0^{2}\cdot q_1^{1}\cdot q_2^{3}\cdot q_3^{2}-210\cdot q_0^{2}\cdot q_1^{3}\cdot q_2^{1}\cdot q_3^{2}\\
&\quad +210\cdot q_0^{3}\cdot q_1^{2}\cdot q_2^{2}\cdot q_3^{1}+105\cdot q_0^{4}\cdot q_1^{1}\cdot q_2^{1}\cdot q_3^{2}\end{align*}
\begin{align*}(T_{h}^4)_{2,9}&=70\cdot q_0^{4}\cdot q_2^{4}+70\cdot q_0^{4}\cdot q_3^{4}+70\cdot q_1^{4}\cdot q_2^{4}+70\cdot q_1^{4}\cdot q_3^{4}\\
&\quad -420\cdot q_0^{2}\cdot q_1^{2}\cdot q_2^{4}-420\cdot q_0^{2}\cdot q_1^{2}\cdot q_3^{4}-420\cdot q_0^{4}\cdot q_2^{2}\cdot q_3^{2}\\
&\quad -420\cdot q_1^{4}\cdot q_2^{2}\cdot q_3^{2}-1120\cdot q_0^{1}\cdot q_1^{3}\cdot q_2^{1}\cdot q_3^{3}+1120\cdot q_0^{1}\cdot q_1^{3}\cdot q_2^{3}\cdot q_3^{1}\\
&\quad +2520\cdot q_0^{2}\cdot q_1^{2}\cdot q_2^{2}\cdot q_3^{2}+1120\cdot q_0^{3}\cdot q_1^{1}\cdot q_2^{1}\cdot q_3^{3}\\
&\quad -1120\cdot q_0^{3}\cdot q_1^{1}\cdot q_2^{3}\cdot q_3^{1}\end{align*}
\begin{align*}(T_{h}^4)_{3,1}&=-2\cdot q_0^{1}\cdot q_3^{7}-2\cdot q_0^{7}\cdot q_3^{1}-2\cdot q_1^{1}\cdot q_2^{7}-2\cdot q_1^{7}\cdot q_2^{1}\\
&\quad +14\cdot q_0^{1}\cdot q_1^{6}\cdot q_3^{1}+42\cdot q_0^{1}\cdot q_2^{2}\cdot q_3^{5}-70\cdot q_0^{1}\cdot q_2^{4}\cdot q_3^{3}\\
&\quad +14\cdot q_0^{1}\cdot q_2^{6}\cdot q_3^{1}+42\cdot q_0^{2}\cdot q_1^{5}\cdot q_2^{1}-70\cdot q_0^{3}\cdot q_1^{4}\cdot q_3^{1}\\
&\quad -70\cdot q_0^{4}\cdot q_1^{3}\cdot q_2^{1}+42\cdot q_0^{5}\cdot q_1^{2}\cdot q_3^{1}+14\cdot q_0^{6}\cdot q_1^{1}\cdot q_2^{1}\\
&\quad +14\cdot q_1^{1}\cdot q_2^{1}\cdot q_3^{6}-70\cdot q_1^{1}\cdot q_2^{3}\cdot q_3^{4}+42\cdot q_1^{1}\cdot q_2^{5}\cdot q_3^{2}\end{align*}
\begin{align*}(T_{h}^4)_{3,2}&=8\cdot q_0^{1}\cdot q_2^{7}-8\cdot q_0^{7}\cdot q_2^{1}-8\cdot q_1^{1}\cdot q_3^{7}+8\cdot q_1^{7}\cdot q_3^{1}\\
&\quad +56\cdot q_0^{1}\cdot q_1^{6}\cdot q_2^{1}-56\cdot q_0^{1}\cdot q_2^{1}\cdot q_3^{6}+280\cdot q_0^{1}\cdot q_2^{3}\cdot q_3^{4}\\
&\quad -168\cdot q_0^{1}\cdot q_2^{5}\cdot q_3^{2}-168\cdot q_0^{2}\cdot q_1^{5}\cdot q_3^{1}-280\cdot q_0^{3}\cdot q_1^{4}\cdot q_2^{1}\\
&\quad +280\cdot q_0^{4}\cdot q_1^{3}\cdot q_3^{1}+168\cdot q_0^{5}\cdot q_1^{2}\cdot q_2^{1}-56\cdot q_0^{6}\cdot q_1^{1}\cdot q_3^{1}\\
&\quad +168\cdot q_1^{1}\cdot q_2^{2}\cdot q_3^{5}-280\cdot q_1^{1}\cdot q_2^{4}\cdot q_3^{3}+56\cdot q_1^{1}\cdot q_2^{6}\cdot q_3^{1}\end{align*}
\begin{align*}(T_{h}^4)_{3,3}&=1\cdot q_0^{8}-1\cdot q_1^{8}+1\cdot q_2^{8}-1\cdot q_3^{8}+14\cdot q_0^{2}\cdot q_1^{6}-7\cdot q_0^{2}\cdot q_2^{6}\\
&\quad +7\cdot q_0^{2}\cdot q_3^{6}-14\cdot q_0^{6}\cdot q_1^{2}-7\cdot q_0^{6}\cdot q_2^{2}-7\cdot q_0^{6}\cdot q_3^{2}\\
&\quad -7\cdot q_1^{2}\cdot q_2^{6}+7\cdot q_1^{2}\cdot q_3^{6}+7\cdot q_1^{6}\cdot q_2^{2}+7\cdot q_1^{6}\cdot q_3^{2}\\
&\quad +14\cdot q_2^{2}\cdot q_3^{6}-14\cdot q_2^{6}\cdot q_3^{2}-105\cdot q_0^{2}\cdot q_1^{4}\cdot q_2^{2}-105\cdot q_0^{2}\cdot q_1^{4}\cdot q_3^{2}\\
&\quad -105\cdot q_0^{2}\cdot q_2^{2}\cdot q_3^{4}+105\cdot q_0^{2}\cdot q_2^{4}\cdot q_3^{2}+105\cdot q_0^{4}\cdot q_1^{2}\cdot q_2^{2}\\
&\quad +105\cdot q_0^{4}\cdot q_1^{2}\cdot q_3^{2}-105\cdot q_1^{2}\cdot q_2^{2}\cdot q_3^{4}+105\cdot q_1^{2}\cdot q_2^{4}\cdot q_3^{2}\end{align*}
\begin{align*}(T_{h}^4)_{3,4}&=6\cdot q_0^{1}\cdot q_1^{7}-14\cdot q_0^{3}\cdot q_1^{5}-14\cdot q_0^{5}\cdot q_1^{3}+6\cdot q_0^{7}\cdot q_1^{1}\\
&\quad -6\cdot q_2^{1}\cdot q_3^{7}+14\cdot q_2^{3}\cdot q_3^{5}+14\cdot q_2^{5}\cdot q_3^{3}-6\cdot q_2^{7}\cdot q_3^{1}\\
&\quad -42\cdot q_0^{1}\cdot q_1^{5}\cdot q_2^{2}-42\cdot q_0^{1}\cdot q_1^{5}\cdot q_3^{2}+42\cdot q_0^{2}\cdot q_2^{1}\cdot q_3^{5}\\
&\quad -140\cdot q_0^{2}\cdot q_2^{3}\cdot q_3^{3}+42\cdot q_0^{2}\cdot q_2^{5}\cdot q_3^{1}+140\cdot q_0^{3}\cdot q_1^{3}\cdot q_2^{2}\\
&\quad +140\cdot q_0^{3}\cdot q_1^{3}\cdot q_3^{2}-42\cdot q_0^{5}\cdot q_1^{1}\cdot q_2^{2}-42\cdot q_0^{5}\cdot q_1^{1}\cdot q_3^{2}\\
&\quad +42\cdot q_1^{2}\cdot q_2^{1}\cdot q_3^{5}-140\cdot q_1^{2}\cdot q_2^{3}\cdot q_3^{3}+42\cdot q_1^{2}\cdot q_2^{5}\cdot q_3^{1}\end{align*}
\begin{align*}(T_{h}^4)_{3,5}&=7\cdot q_0^{1}\cdot q_3^{7}-21\cdot q_0^{3}\cdot q_3^{5}-21\cdot q_0^{5}\cdot q_3^{3}+7\cdot q_0^{7}\cdot q_3^{1}\\
&\quad +7\cdot q_1^{1}\cdot q_2^{7}-21\cdot q_1^{3}\cdot q_2^{5}-21\cdot q_1^{5}\cdot q_2^{3}+7\cdot q_1^{7}\cdot q_2^{1}\\
&\quad -21\cdot q_0^{1}\cdot q_1^{2}\cdot q_3^{5}-105\cdot q_0^{1}\cdot q_1^{4}\cdot q_3^{3}+35\cdot q_0^{1}\cdot q_1^{6}\cdot q_3^{1}\\
&\quad -63\cdot q_0^{1}\cdot q_2^{2}\cdot q_3^{5}-35\cdot q_0^{1}\cdot q_2^{4}\cdot q_3^{3}+35\cdot q_0^{1}\cdot q_2^{6}\cdot q_3^{1}\\
&\quad -21\cdot q_0^{2}\cdot q_1^{1}\cdot q_2^{5}+210\cdot q_0^{2}\cdot q_1^{3}\cdot q_2^{3}-63\cdot q_0^{2}\cdot q_1^{5}\cdot q_2^{1}\\
&\quad +210\cdot q_0^{3}\cdot q_1^{2}\cdot q_3^{3}-35\cdot q_0^{3}\cdot q_1^{4}\cdot q_3^{1}+210\cdot q_0^{3}\cdot q_2^{2}\cdot q_3^{3}\\
&\quad -105\cdot q_0^{3}\cdot q_2^{4}\cdot q_3^{1}-105\cdot q_0^{4}\cdot q_1^{1}\cdot q_2^{3}-35\cdot q_0^{4}\cdot q_1^{3}\cdot q_2^{1}\\
&\quad -63\cdot q_0^{5}\cdot q_1^{2}\cdot q_3^{1}-21\cdot q_0^{5}\cdot q_2^{2}\cdot q_3^{1}+35\cdot q_0^{6}\cdot q_1^{1}\cdot q_2^{1}\\
&\quad +35\cdot q_1^{1}\cdot q_2^{1}\cdot q_3^{6}-35\cdot q_1^{1}\cdot q_2^{3}\cdot q_3^{4}-63\cdot q_1^{1}\cdot q_2^{5}\cdot q_3^{2}\\
&\quad -105\cdot q_1^{3}\cdot q_2^{1}\cdot q_3^{4}+210\cdot q_1^{3}\cdot q_2^{3}\cdot q_3^{2}-21\cdot q_1^{5}\cdot q_2^{1}\cdot q_3^{2}\\
&\quad +210\cdot q_0^{1}\cdot q_1^{2}\cdot q_2^{2}\cdot q_3^{3}-105\cdot q_0^{1}\cdot q_1^{2}\cdot q_2^{4}\cdot q_3^{1}\\
&\quad -105\cdot q_0^{1}\cdot q_1^{4}\cdot q_2^{2}\cdot q_3^{1}-105\cdot q_0^{2}\cdot q_1^{1}\cdot q_2^{1}\cdot q_3^{4}\\
&\quad +210\cdot q_0^{2}\cdot q_1^{1}\cdot q_2^{3}\cdot q_3^{2}+210\cdot q_0^{2}\cdot q_1^{3}\cdot q_2^{1}\cdot q_3^{2}\\
&\quad +210\cdot q_0^{3}\cdot q_1^{2}\cdot q_2^{2}\cdot q_3^{1}-105\cdot q_0^{4}\cdot q_1^{1}\cdot q_2^{1}\cdot q_3^{2}\end{align*}
\begin{align*}(T_{h}^4)_{3,6}&=14\cdot q_0^{1}\cdot q_2^{7}-42\cdot q_0^{3}\cdot q_2^{5}+42\cdot q_0^{5}\cdot q_2^{3}-14\cdot q_0^{7}\cdot q_2^{1}\\
&\quad -14\cdot q_1^{1}\cdot q_3^{7}+42\cdot q_1^{3}\cdot q_3^{5}-42\cdot q_1^{5}\cdot q_3^{3}+14\cdot q_1^{7}\cdot q_3^{1}\\
&\quad -42\cdot q_0^{1}\cdot q_1^{2}\cdot q_2^{5}+210\cdot q_0^{1}\cdot q_1^{4}\cdot q_2^{3}-70\cdot q_0^{1}\cdot q_1^{6}\cdot q_2^{1}\\
&\quad +70\cdot q_0^{1}\cdot q_2^{1}\cdot q_3^{6}-70\cdot q_0^{1}\cdot q_2^{3}\cdot q_3^{4}-126\cdot q_0^{1}\cdot q_2^{5}\cdot q_3^{2}\\
&\quad +42\cdot q_0^{2}\cdot q_1^{1}\cdot q_3^{5}+420\cdot q_0^{2}\cdot q_1^{3}\cdot q_3^{3}-126\cdot q_0^{2}\cdot q_1^{5}\cdot q_3^{1}\\
&\quad -420\cdot q_0^{3}\cdot q_1^{2}\cdot q_2^{3}+70\cdot q_0^{3}\cdot q_1^{4}\cdot q_2^{1}-210\cdot q_0^{3}\cdot q_2^{1}\cdot q_3^{4}\\
&\quad +420\cdot q_0^{3}\cdot q_2^{3}\cdot q_3^{2}-210\cdot q_0^{4}\cdot q_1^{1}\cdot q_3^{3}-70\cdot q_0^{4}\cdot q_1^{3}\cdot q_3^{1}\\
&\quad +126\cdot q_0^{5}\cdot q_1^{2}\cdot q_2^{1}+42\cdot q_0^{5}\cdot q_2^{1}\cdot q_3^{2}+70\cdot q_0^{6}\cdot q_1^{1}\cdot q_3^{1}\\
&\quad +126\cdot q_1^{1}\cdot q_2^{2}\cdot q_3^{5}+70\cdot q_1^{1}\cdot q_2^{4}\cdot q_3^{3}-70\cdot q_1^{1}\cdot q_2^{6}\cdot q_3^{1}\\
&\quad -420\cdot q_1^{3}\cdot q_2^{2}\cdot q_3^{3}+210\cdot q_1^{3}\cdot q_2^{4}\cdot q_3^{1}-42\cdot q_1^{5}\cdot q_2^{2}\cdot q_3^{1}\\
&\quad -210\cdot q_0^{1}\cdot q_1^{2}\cdot q_2^{1}\cdot q_3^{4}+420\cdot q_0^{1}\cdot q_1^{2}\cdot q_2^{3}\cdot q_3^{2}\\
&\quad +210\cdot q_0^{1}\cdot q_1^{4}\cdot q_2^{1}\cdot q_3^{2}-420\cdot q_0^{2}\cdot q_1^{1}\cdot q_2^{2}\cdot q_3^{3}\\
&\quad +210\cdot q_0^{2}\cdot q_1^{1}\cdot q_2^{4}\cdot q_3^{1}+420\cdot q_0^{2}\cdot q_1^{3}\cdot q_2^{2}\cdot q_3^{1}\\
&\quad -420\cdot q_0^{3}\cdot q_1^{2}\cdot q_2^{1}\cdot q_3^{2}-210\cdot q_0^{4}\cdot q_1^{1}\cdot q_2^{2}\cdot q_3^{1}\end{align*}
\begin{align*}(T_{h}^4)_{3,7}&=-21\cdot q_0^{2}\cdot q_2^{6}-21\cdot q_0^{2}\cdot q_3^{6}+70\cdot q_0^{4}\cdot q_2^{4}-21\cdot q_0^{6}\cdot q_2^{2}\\
&\quad +21\cdot q_0^{6}\cdot q_3^{2}+21\cdot q_1^{2}\cdot q_2^{6}+21\cdot q_1^{2}\cdot q_3^{6}-70\cdot q_1^{4}\cdot q_3^{4}\\
&\quad -21\cdot q_1^{6}\cdot q_2^{2}+21\cdot q_1^{6}\cdot q_3^{2}-210\cdot q_0^{2}\cdot q_1^{2}\cdot q_2^{4}+210\cdot q_0^{2}\cdot q_1^{2}\cdot q_3^{4}\\
&\quad +105\cdot q_0^{2}\cdot q_1^{4}\cdot q_2^{2}-105\cdot q_0^{2}\cdot q_1^{4}\cdot q_3^{2}+105\cdot q_0^{2}\cdot q_2^{2}\cdot q_3^{4}\\
&\quad +105\cdot q_0^{2}\cdot q_2^{4}\cdot q_3^{2}+105\cdot q_0^{4}\cdot q_1^{2}\cdot q_2^{2}-105\cdot q_0^{4}\cdot q_1^{2}\cdot q_3^{2}\\
&\quad -210\cdot q_0^{4}\cdot q_2^{2}\cdot q_3^{2}-105\cdot q_1^{2}\cdot q_2^{2}\cdot q_3^{4}-105\cdot q_1^{2}\cdot q_2^{4}\cdot q_3^{2}\\
&\quad +210\cdot q_1^{4}\cdot q_2^{2}\cdot q_3^{2}-168\cdot q_0^{1}\cdot q_1^{1}\cdot q_2^{1}\cdot q_3^{5}+168\cdot q_0^{1}\cdot q_1^{1}\cdot q_2^{5}\cdot q_3^{1}\\
&\quad +560\cdot q_0^{1}\cdot q_1^{3}\cdot q_2^{1}\cdot q_3^{3}-168\cdot q_0^{1}\cdot q_1^{5}\cdot q_2^{1}\cdot q_3^{1}\\
&\quad -560\cdot q_0^{3}\cdot q_1^{1}\cdot q_2^{3}\cdot q_3^{1}+168\cdot q_0^{5}\cdot q_1^{1}\cdot q_2^{1}\cdot q_3^{1}\end{align*}
\begin{align*}(T_{h}^4)_{3,8}&=42\cdot q_0^{1}\cdot q_1^{1}\cdot q_2^{6}+42\cdot q_0^{1}\cdot q_1^{1}\cdot q_3^{6}-210\cdot q_0^{1}\cdot q_1^{3}\cdot q_2^{4}\\
&\quad +70\cdot q_0^{1}\cdot q_1^{3}\cdot q_3^{4}+84\cdot q_0^{1}\cdot q_1^{5}\cdot q_2^{2}-84\cdot q_0^{1}\cdot q_1^{5}\cdot q_3^{2}\\
&\quad -84\cdot q_0^{2}\cdot q_2^{1}\cdot q_3^{5}+84\cdot q_0^{2}\cdot q_2^{5}\cdot q_3^{1}+70\cdot q_0^{3}\cdot q_1^{1}\cdot q_2^{4}\\
&\quad -210\cdot q_0^{3}\cdot q_1^{1}\cdot q_3^{4}+210\cdot q_0^{4}\cdot q_2^{1}\cdot q_3^{3}-70\cdot q_0^{4}\cdot q_2^{3}\cdot q_3^{1}\\
&\quad -84\cdot q_0^{5}\cdot q_1^{1}\cdot q_2^{2}+84\cdot q_0^{5}\cdot q_1^{1}\cdot q_3^{2}-42\cdot q_0^{6}\cdot q_2^{1}\cdot q_3^{1}\\
&\quad +84\cdot q_1^{2}\cdot q_2^{1}\cdot q_3^{5}-84\cdot q_1^{2}\cdot q_2^{5}\cdot q_3^{1}-70\cdot q_1^{4}\cdot q_2^{1}\cdot q_3^{3}\\
&\quad +210\cdot q_1^{4}\cdot q_2^{3}\cdot q_3^{1}-42\cdot q_1^{6}\cdot q_2^{1}\cdot q_3^{1}-210\cdot q_0^{1}\cdot q_1^{1}\cdot q_2^{2}\cdot q_3^{4}\\
&\quad -210\cdot q_0^{1}\cdot q_1^{1}\cdot q_2^{4}\cdot q_3^{2}+420\cdot q_0^{1}\cdot q_1^{3}\cdot q_2^{2}\cdot q_3^{2}\\
&\quad -420\cdot q_0^{2}\cdot q_1^{2}\cdot q_2^{1}\cdot q_3^{3}-420\cdot q_0^{2}\cdot q_1^{2}\cdot q_2^{3}\cdot q_3^{1}\\
&\quad +210\cdot q_0^{2}\cdot q_1^{4}\cdot q_2^{1}\cdot q_3^{1}+420\cdot q_0^{3}\cdot q_1^{1}\cdot q_2^{2}\cdot q_3^{2}\\
&\quad +210\cdot q_0^{4}\cdot q_1^{2}\cdot q_2^{1}\cdot q_3^{1}\end{align*}
\begin{align*}(T_{h}^4)_{3,9}&=-280\cdot q_0^{3}\cdot q_2^{5}+280\cdot q_0^{5}\cdot q_2^{3}+280\cdot q_1^{3}\cdot q_3^{5}-280\cdot q_1^{5}\cdot q_3^{3}\\
&\quad +840\cdot q_0^{1}\cdot q_1^{2}\cdot q_2^{5}-840\cdot q_0^{1}\cdot q_1^{4}\cdot q_2^{3}-840\cdot q_0^{2}\cdot q_1^{1}\cdot q_3^{5}\\
&\quad +560\cdot q_0^{2}\cdot q_1^{3}\cdot q_3^{3}-560\cdot q_0^{3}\cdot q_1^{2}\cdot q_2^{3}+840\cdot q_0^{3}\cdot q_2^{1}\cdot q_3^{4}\\
&\quad +560\cdot q_0^{3}\cdot q_2^{3}\cdot q_3^{2}+840\cdot q_0^{4}\cdot q_1^{1}\cdot q_3^{3}-840\cdot q_0^{5}\cdot q_2^{1}\cdot q_3^{2}\\
&\quad -560\cdot q_1^{3}\cdot q_2^{2}\cdot q_3^{3}-840\cdot q_1^{3}\cdot q_2^{4}\cdot q_3^{1}+840\cdot q_1^{5}\cdot q_2^{2}\cdot q_3^{1}\\
&\quad -2520\cdot q_0^{1}\cdot q_1^{2}\cdot q_2^{1}\cdot q_3^{4}-1680\cdot q_0^{1}\cdot q_1^{2}\cdot q_2^{3}\cdot q_3^{2}\\
&\quad +2520\cdot q_0^{1}\cdot q_1^{4}\cdot q_2^{1}\cdot q_3^{2}+1680\cdot q_0^{2}\cdot q_1^{1}\cdot q_2^{2}\cdot q_3^{3}\\
&\quad +2520\cdot q_0^{2}\cdot q_1^{1}\cdot q_2^{4}\cdot q_3^{1}-1680\cdot q_0^{2}\cdot q_1^{3}\cdot q_2^{2}\cdot q_3^{1}\\
&\quad +1680\cdot q_0^{3}\cdot q_1^{2}\cdot q_2^{1}\cdot q_3^{2}-2520\cdot q_0^{4}\cdot q_1^{1}\cdot q_2^{2}\cdot q_3^{1}\end{align*}
\begin{align*}(T_{h}^4)_{4,1}&=2\cdot q_0^{1}\cdot q_2^{7}+2\cdot q_0^{7}\cdot q_2^{1}-2\cdot q_1^{1}\cdot q_3^{7}-2\cdot q_1^{7}\cdot q_3^{1}\\
&\quad -14\cdot q_0^{1}\cdot q_1^{6}\cdot q_2^{1}-14\cdot q_0^{1}\cdot q_2^{1}\cdot q_3^{6}+70\cdot q_0^{1}\cdot q_2^{3}\cdot q_3^{4}\\
&\quad -42\cdot q_0^{1}\cdot q_2^{5}\cdot q_3^{2}+42\cdot q_0^{2}\cdot q_1^{5}\cdot q_3^{1}+70\cdot q_0^{3}\cdot q_1^{4}\cdot q_2^{1}\\
&\quad -70\cdot q_0^{4}\cdot q_1^{3}\cdot q_3^{1}-42\cdot q_0^{5}\cdot q_1^{2}\cdot q_2^{1}+14\cdot q_0^{6}\cdot q_1^{1}\cdot q_3^{1}\\
&\quad +42\cdot q_1^{1}\cdot q_2^{2}\cdot q_3^{5}-70\cdot q_1^{1}\cdot q_2^{4}\cdot q_3^{3}+14\cdot q_1^{1}\cdot q_2^{6}\cdot q_3^{1}\end{align*}
\begin{align*}(T_{h}^4)_{4,2}&=8\cdot q_0^{1}\cdot q_3^{7}-8\cdot q_0^{7}\cdot q_3^{1}+8\cdot q_1^{1}\cdot q_2^{7}-8\cdot q_1^{7}\cdot q_2^{1}\\
&\quad +56\cdot q_0^{1}\cdot q_1^{6}\cdot q_3^{1}-168\cdot q_0^{1}\cdot q_2^{2}\cdot q_3^{5}+280\cdot q_0^{1}\cdot q_2^{4}\cdot q_3^{3}\\
&\quad -56\cdot q_0^{1}\cdot q_2^{6}\cdot q_3^{1}+168\cdot q_0^{2}\cdot q_1^{5}\cdot q_2^{1}-280\cdot q_0^{3}\cdot q_1^{4}\cdot q_3^{1}\\
&\quad -280\cdot q_0^{4}\cdot q_1^{3}\cdot q_2^{1}+168\cdot q_0^{5}\cdot q_1^{2}\cdot q_3^{1}+56\cdot q_0^{6}\cdot q_1^{1}\cdot q_2^{1}\\
&\quad -56\cdot q_1^{1}\cdot q_2^{1}\cdot q_3^{6}+280\cdot q_1^{1}\cdot q_2^{3}\cdot q_3^{4}-168\cdot q_1^{1}\cdot q_2^{5}\cdot q_3^{2}\end{align*}
\begin{align*}(T_{h}^4)_{4,3}&=-6\cdot q_0^{1}\cdot q_1^{7}+14\cdot q_0^{3}\cdot q_1^{5}+14\cdot q_0^{5}\cdot q_1^{3}-6\cdot q_0^{7}\cdot q_1^{1}\\
&\quad -6\cdot q_2^{1}\cdot q_3^{7}+14\cdot q_2^{3}\cdot q_3^{5}+14\cdot q_2^{5}\cdot q_3^{3}-6\cdot q_2^{7}\cdot q_3^{1}\\
&\quad +42\cdot q_0^{1}\cdot q_1^{5}\cdot q_2^{2}+42\cdot q_0^{1}\cdot q_1^{5}\cdot q_3^{2}+42\cdot q_0^{2}\cdot q_2^{1}\cdot q_3^{5}\\
&\quad -140\cdot q_0^{2}\cdot q_2^{3}\cdot q_3^{3}+42\cdot q_0^{2}\cdot q_2^{5}\cdot q_3^{1}-140\cdot q_0^{3}\cdot q_1^{3}\cdot q_2^{2}\\
&\quad -140\cdot q_0^{3}\cdot q_1^{3}\cdot q_3^{2}+42\cdot q_0^{5}\cdot q_1^{1}\cdot q_2^{2}+42\cdot q_0^{5}\cdot q_1^{1}\cdot q_3^{2}\\
&\quad +42\cdot q_1^{2}\cdot q_2^{1}\cdot q_3^{5}-140\cdot q_1^{2}\cdot q_2^{3}\cdot q_3^{3}+42\cdot q_1^{2}\cdot q_2^{5}\cdot q_3^{1}\end{align*}
\begin{align*}(T_{h}^4)_{4,4}&=1\cdot q_0^{8}-1\cdot q_1^{8}-1\cdot q_2^{8}+1\cdot q_3^{8}+14\cdot q_0^{2}\cdot q_1^{6}+7\cdot q_0^{2}\cdot q_2^{6}\\
&\quad -7\cdot q_0^{2}\cdot q_3^{6}-14\cdot q_0^{6}\cdot q_1^{2}-7\cdot q_0^{6}\cdot q_2^{2}-7\cdot q_0^{6}\cdot q_3^{2}\\
&\quad +7\cdot q_1^{2}\cdot q_2^{6}-7\cdot q_1^{2}\cdot q_3^{6}+7\cdot q_1^{6}\cdot q_2^{2}+7\cdot q_1^{6}\cdot q_3^{2}\\
&\quad -14\cdot q_2^{2}\cdot q_3^{6}+14\cdot q_2^{6}\cdot q_3^{2}-105\cdot q_0^{2}\cdot q_1^{4}\cdot q_2^{2}-105\cdot q_0^{2}\cdot q_1^{4}\cdot q_3^{2}\\
&\quad +105\cdot q_0^{2}\cdot q_2^{2}\cdot q_3^{4}-105\cdot q_0^{2}\cdot q_2^{4}\cdot q_3^{2}+105\cdot q_0^{4}\cdot q_1^{2}\cdot q_2^{2}\\
&\quad +105\cdot q_0^{4}\cdot q_1^{2}\cdot q_3^{2}+105\cdot q_1^{2}\cdot q_2^{2}\cdot q_3^{4}-105\cdot q_1^{2}\cdot q_2^{4}\cdot q_3^{2}\end{align*}
\begin{align*}(T_{h}^4)_{4,5}&=7\cdot q_0^{1}\cdot q_2^{7}-21\cdot q_0^{3}\cdot q_2^{5}-21\cdot q_0^{5}\cdot q_2^{3}+7\cdot q_0^{7}\cdot q_2^{1}\\
&\quad -7\cdot q_1^{1}\cdot q_3^{7}+21\cdot q_1^{3}\cdot q_3^{5}+21\cdot q_1^{5}\cdot q_3^{3}-7\cdot q_1^{7}\cdot q_3^{1}\\
&\quad -21\cdot q_0^{1}\cdot q_1^{2}\cdot q_2^{5}-105\cdot q_0^{1}\cdot q_1^{4}\cdot q_2^{3}+35\cdot q_0^{1}\cdot q_1^{6}\cdot q_2^{1}\\
&\quad +35\cdot q_0^{1}\cdot q_2^{1}\cdot q_3^{6}-35\cdot q_0^{1}\cdot q_2^{3}\cdot q_3^{4}-63\cdot q_0^{1}\cdot q_2^{5}\cdot q_3^{2}\\
&\quad +21\cdot q_0^{2}\cdot q_1^{1}\cdot q_3^{5}-210\cdot q_0^{2}\cdot q_1^{3}\cdot q_3^{3}+63\cdot q_0^{2}\cdot q_1^{5}\cdot q_3^{1}\\
&\quad +210\cdot q_0^{3}\cdot q_1^{2}\cdot q_2^{3}-35\cdot q_0^{3}\cdot q_1^{4}\cdot q_2^{1}-105\cdot q_0^{3}\cdot q_2^{1}\cdot q_3^{4}\\
&\quad +210\cdot q_0^{3}\cdot q_2^{3}\cdot q_3^{2}+105\cdot q_0^{4}\cdot q_1^{1}\cdot q_3^{3}+35\cdot q_0^{4}\cdot q_1^{3}\cdot q_3^{1}\\
&\quad -63\cdot q_0^{5}\cdot q_1^{2}\cdot q_2^{1}-21\cdot q_0^{5}\cdot q_2^{1}\cdot q_3^{2}-35\cdot q_0^{6}\cdot q_1^{1}\cdot q_3^{1}\\
&\quad +63\cdot q_1^{1}\cdot q_2^{2}\cdot q_3^{5}+35\cdot q_1^{1}\cdot q_2^{4}\cdot q_3^{3}-35\cdot q_1^{1}\cdot q_2^{6}\cdot q_3^{1}\\
&\quad -210\cdot q_1^{3}\cdot q_2^{2}\cdot q_3^{3}+105\cdot q_1^{3}\cdot q_2^{4}\cdot q_3^{1}+21\cdot q_1^{5}\cdot q_2^{2}\cdot q_3^{1}\\
&\quad -105\cdot q_0^{1}\cdot q_1^{2}\cdot q_2^{1}\cdot q_3^{4}+210\cdot q_0^{1}\cdot q_1^{2}\cdot q_2^{3}\cdot q_3^{2}\\
&\quad -105\cdot q_0^{1}\cdot q_1^{4}\cdot q_2^{1}\cdot q_3^{2}-210\cdot q_0^{2}\cdot q_1^{1}\cdot q_2^{2}\cdot q_3^{3}\\
&\quad +105\cdot q_0^{2}\cdot q_1^{1}\cdot q_2^{4}\cdot q_3^{1}-210\cdot q_0^{2}\cdot q_1^{3}\cdot q_2^{2}\cdot q_3^{1}\\
&\quad +210\cdot q_0^{3}\cdot q_1^{2}\cdot q_2^{1}\cdot q_3^{2}+105\cdot q_0^{4}\cdot q_1^{1}\cdot q_2^{2}\cdot q_3^{1}\end{align*}
\begin{align*}(T_{h}^4)_{4,6}&=-14\cdot q_0^{1}\cdot q_3^{7}+42\cdot q_0^{3}\cdot q_3^{5}-42\cdot q_0^{5}\cdot q_3^{3}+14\cdot q_0^{7}\cdot q_3^{1}\\
&\quad -14\cdot q_1^{1}\cdot q_2^{7}+42\cdot q_1^{3}\cdot q_2^{5}-42\cdot q_1^{5}\cdot q_2^{3}+14\cdot q_1^{7}\cdot q_2^{1}\\
&\quad +42\cdot q_0^{1}\cdot q_1^{2}\cdot q_3^{5}-210\cdot q_0^{1}\cdot q_1^{4}\cdot q_3^{3}+70\cdot q_0^{1}\cdot q_1^{6}\cdot q_3^{1}\\
&\quad +126\cdot q_0^{1}\cdot q_2^{2}\cdot q_3^{5}+70\cdot q_0^{1}\cdot q_2^{4}\cdot q_3^{3}-70\cdot q_0^{1}\cdot q_2^{6}\cdot q_3^{1}\\
&\quad +42\cdot q_0^{2}\cdot q_1^{1}\cdot q_2^{5}+420\cdot q_0^{2}\cdot q_1^{3}\cdot q_2^{3}-126\cdot q_0^{2}\cdot q_1^{5}\cdot q_2^{1}\\
&\quad +420\cdot q_0^{3}\cdot q_1^{2}\cdot q_3^{3}-70\cdot q_0^{3}\cdot q_1^{4}\cdot q_3^{1}-420\cdot q_0^{3}\cdot q_2^{2}\cdot q_3^{3}\\
&\quad +210\cdot q_0^{3}\cdot q_2^{4}\cdot q_3^{1}-210\cdot q_0^{4}\cdot q_1^{1}\cdot q_2^{3}-70\cdot q_0^{4}\cdot q_1^{3}\cdot q_2^{1}\\
&\quad -126\cdot q_0^{5}\cdot q_1^{2}\cdot q_3^{1}-42\cdot q_0^{5}\cdot q_2^{2}\cdot q_3^{1}+70\cdot q_0^{6}\cdot q_1^{1}\cdot q_2^{1}\\
&\quad -70\cdot q_1^{1}\cdot q_2^{1}\cdot q_3^{6}+70\cdot q_1^{1}\cdot q_2^{3}\cdot q_3^{4}+126\cdot q_1^{1}\cdot q_2^{5}\cdot q_3^{2}\\
&\quad +210\cdot q_1^{3}\cdot q_2^{1}\cdot q_3^{4}-420\cdot q_1^{3}\cdot q_2^{3}\cdot q_3^{2}-42\cdot q_1^{5}\cdot q_2^{1}\cdot q_3^{2}\\
&\quad -420\cdot q_0^{1}\cdot q_1^{2}\cdot q_2^{2}\cdot q_3^{3}+210\cdot q_0^{1}\cdot q_1^{2}\cdot q_2^{4}\cdot q_3^{1}\\
&\quad -210\cdot q_0^{1}\cdot q_1^{4}\cdot q_2^{2}\cdot q_3^{1}+210\cdot q_0^{2}\cdot q_1^{1}\cdot q_2^{1}\cdot q_3^{4}\\
&\quad -420\cdot q_0^{2}\cdot q_1^{1}\cdot q_2^{3}\cdot q_3^{2}+420\cdot q_0^{2}\cdot q_1^{3}\cdot q_2^{1}\cdot q_3^{2}\\
&\quad +420\cdot q_0^{3}\cdot q_1^{2}\cdot q_2^{2}\cdot q_3^{1}-210\cdot q_0^{4}\cdot q_1^{1}\cdot q_2^{1}\cdot q_3^{2}\end{align*}
\begin{align*}(T_{h}^4)_{4,7}&=42\cdot q_0^{1}\cdot q_1^{1}\cdot q_2^{6}+42\cdot q_0^{1}\cdot q_1^{1}\cdot q_3^{6}+70\cdot q_0^{1}\cdot q_1^{3}\cdot q_2^{4}\\
&\quad -210\cdot q_0^{1}\cdot q_1^{3}\cdot q_3^{4}-84\cdot q_0^{1}\cdot q_1^{5}\cdot q_2^{2}+84\cdot q_0^{1}\cdot q_1^{5}\cdot q_3^{2}\\
&\quad -84\cdot q_0^{2}\cdot q_2^{1}\cdot q_3^{5}+84\cdot q_0^{2}\cdot q_2^{5}\cdot q_3^{1}-210\cdot q_0^{3}\cdot q_1^{1}\cdot q_2^{4}\\
&\quad +70\cdot q_0^{3}\cdot q_1^{1}\cdot q_3^{4}+70\cdot q_0^{4}\cdot q_2^{1}\cdot q_3^{3}-210\cdot q_0^{4}\cdot q_2^{3}\cdot q_3^{1}\\
&\quad +84\cdot q_0^{5}\cdot q_1^{1}\cdot q_2^{2}-84\cdot q_0^{5}\cdot q_1^{1}\cdot q_3^{2}+42\cdot q_0^{6}\cdot q_2^{1}\cdot q_3^{1}\\
&\quad +84\cdot q_1^{2}\cdot q_2^{1}\cdot q_3^{5}-84\cdot q_1^{2}\cdot q_2^{5}\cdot q_3^{1}-210\cdot q_1^{4}\cdot q_2^{1}\cdot q_3^{3}\\
&\quad +70\cdot q_1^{4}\cdot q_2^{3}\cdot q_3^{1}+42\cdot q_1^{6}\cdot q_2^{1}\cdot q_3^{1}-210\cdot q_0^{1}\cdot q_1^{1}\cdot q_2^{2}\cdot q_3^{4}\\
&\quad -210\cdot q_0^{1}\cdot q_1^{1}\cdot q_2^{4}\cdot q_3^{2}+420\cdot q_0^{1}\cdot q_1^{3}\cdot q_2^{2}\cdot q_3^{2}\\
&\quad +420\cdot q_0^{2}\cdot q_1^{2}\cdot q_2^{1}\cdot q_3^{3}+420\cdot q_0^{2}\cdot q_1^{2}\cdot q_2^{3}\cdot q_3^{1}\\
&\quad -210\cdot q_0^{2}\cdot q_1^{4}\cdot q_2^{1}\cdot q_3^{1}+420\cdot q_0^{3}\cdot q_1^{1}\cdot q_2^{2}\cdot q_3^{2}\\
&\quad -210\cdot q_0^{4}\cdot q_1^{2}\cdot q_2^{1}\cdot q_3^{1}\end{align*}
\begin{align*}(T_{h}^4)_{4,8}&=21\cdot q_0^{2}\cdot q_2^{6}+21\cdot q_0^{2}\cdot q_3^{6}-70\cdot q_0^{4}\cdot q_3^{4}-21\cdot q_0^{6}\cdot q_2^{2}\\
&\quad +21\cdot q_0^{6}\cdot q_3^{2}-21\cdot q_1^{2}\cdot q_2^{6}-21\cdot q_1^{2}\cdot q_3^{6}+70\cdot q_1^{4}\cdot q_2^{4}\\
&\quad -21\cdot q_1^{6}\cdot q_2^{2}+21\cdot q_1^{6}\cdot q_3^{2}-210\cdot q_0^{2}\cdot q_1^{2}\cdot q_2^{4}+210\cdot q_0^{2}\cdot q_1^{2}\cdot q_3^{4}\\
&\quad +105\cdot q_0^{2}\cdot q_1^{4}\cdot q_2^{2}-105\cdot q_0^{2}\cdot q_1^{4}\cdot q_3^{2}-105\cdot q_0^{2}\cdot q_2^{2}\cdot q_3^{4}\\
&\quad -105\cdot q_0^{2}\cdot q_2^{4}\cdot q_3^{2}+105\cdot q_0^{4}\cdot q_1^{2}\cdot q_2^{2}-105\cdot q_0^{4}\cdot q_1^{2}\cdot q_3^{2}\\
&\quad +210\cdot q_0^{4}\cdot q_2^{2}\cdot q_3^{2}+105\cdot q_1^{2}\cdot q_2^{2}\cdot q_3^{4}+105\cdot q_1^{2}\cdot q_2^{4}\cdot q_3^{2}\\
&\quad -210\cdot q_1^{4}\cdot q_2^{2}\cdot q_3^{2}+168\cdot q_0^{1}\cdot q_1^{1}\cdot q_2^{1}\cdot q_3^{5}-168\cdot q_0^{1}\cdot q_1^{1}\cdot q_2^{5}\cdot q_3^{1}\\
&\quad +560\cdot q_0^{1}\cdot q_1^{3}\cdot q_2^{3}\cdot q_3^{1}-168\cdot q_0^{1}\cdot q_1^{5}\cdot q_2^{1}\cdot q_3^{1}\\
&\quad -560\cdot q_0^{3}\cdot q_1^{1}\cdot q_2^{1}\cdot q_3^{3}+168\cdot q_0^{5}\cdot q_1^{1}\cdot q_2^{1}\cdot q_3^{1}\end{align*}
\begin{align*}(T_{h}^4)_{4,9}&=-280\cdot q_0^{3}\cdot q_3^{5}+280\cdot q_0^{5}\cdot q_3^{3}-280\cdot q_1^{3}\cdot q_2^{5}+280\cdot q_1^{5}\cdot q_2^{3}\\
&\quad +840\cdot q_0^{1}\cdot q_1^{2}\cdot q_3^{5}-840\cdot q_0^{1}\cdot q_1^{4}\cdot q_3^{3}+840\cdot q_0^{2}\cdot q_1^{1}\cdot q_2^{5}\\
&\quad -560\cdot q_0^{2}\cdot q_1^{3}\cdot q_2^{3}-560\cdot q_0^{3}\cdot q_1^{2}\cdot q_3^{3}+560\cdot q_0^{3}\cdot q_2^{2}\cdot q_3^{3}\\
&\quad +840\cdot q_0^{3}\cdot q_2^{4}\cdot q_3^{1}-840\cdot q_0^{4}\cdot q_1^{1}\cdot q_2^{3}-840\cdot q_0^{5}\cdot q_2^{2}\cdot q_3^{1}\\
&\quad +840\cdot q_1^{3}\cdot q_2^{1}\cdot q_3^{4}+560\cdot q_1^{3}\cdot q_2^{3}\cdot q_3^{2}-840\cdot q_1^{5}\cdot q_2^{1}\cdot q_3^{2}\\
&\quad -1680\cdot q_0^{1}\cdot q_1^{2}\cdot q_2^{2}\cdot q_3^{3}-2520\cdot q_0^{1}\cdot q_1^{2}\cdot q_2^{4}\cdot q_3^{1}\\
&\quad +2520\cdot q_0^{1}\cdot q_1^{4}\cdot q_2^{2}\cdot q_3^{1}-2520\cdot q_0^{2}\cdot q_1^{1}\cdot q_2^{1}\cdot q_3^{4}\\
&\quad -1680\cdot q_0^{2}\cdot q_1^{1}\cdot q_2^{3}\cdot q_3^{2}+1680\cdot q_0^{2}\cdot q_1^{3}\cdot q_2^{1}\cdot q_3^{2}\\
&\quad +1680\cdot q_0^{3}\cdot q_1^{2}\cdot q_2^{2}\cdot q_3^{1}+2520\cdot q_0^{4}\cdot q_1^{1}\cdot q_2^{1}\cdot q_3^{2}\end{align*}
\begin{align*}(T_{h}^4)_{5,1}&=2\cdot q_0^{2}\cdot q_2^{6}-2\cdot q_0^{2}\cdot q_3^{6}-2\cdot q_0^{6}\cdot q_2^{2}+2\cdot q_0^{6}\cdot q_3^{2}\\
&\quad -2\cdot q_1^{2}\cdot q_2^{6}+2\cdot q_1^{2}\cdot q_3^{6}+2\cdot q_1^{6}\cdot q_2^{2}-2\cdot q_1^{6}\cdot q_3^{2}\\
&\quad -30\cdot q_0^{2}\cdot q_1^{4}\cdot q_2^{2}+30\cdot q_0^{2}\cdot q_1^{4}\cdot q_3^{2}+30\cdot q_0^{2}\cdot q_2^{2}\cdot q_3^{4}\\
&\quad -30\cdot q_0^{2}\cdot q_2^{4}\cdot q_3^{2}+30\cdot q_0^{4}\cdot q_1^{2}\cdot q_2^{2}-30\cdot q_0^{4}\cdot q_1^{2}\cdot q_3^{2}\\
&\quad -30\cdot q_1^{2}\cdot q_2^{2}\cdot q_3^{4}+30\cdot q_1^{2}\cdot q_2^{4}\cdot q_3^{2}+24\cdot q_0^{1}\cdot q_1^{1}\cdot q_2^{1}\cdot q_3^{5}\\
&\quad -80\cdot q_0^{1}\cdot q_1^{1}\cdot q_2^{3}\cdot q_3^{3}+24\cdot q_0^{1}\cdot q_1^{1}\cdot q_2^{5}\cdot q_3^{1}\\
&\quad -24\cdot q_0^{1}\cdot q_1^{5}\cdot q_2^{1}\cdot q_3^{1}+80\cdot q_0^{3}\cdot q_1^{3}\cdot q_2^{1}\cdot q_3^{1}\\
&\quad -24\cdot q_0^{5}\cdot q_1^{1}\cdot q_2^{1}\cdot q_3^{1}\end{align*}
\begin{align*}(T_{h}^4)_{5,2}&=16\cdot q_0^{1}\cdot q_1^{1}\cdot q_2^{6}-16\cdot q_0^{1}\cdot q_1^{1}\cdot q_3^{6}-48\cdot q_0^{1}\cdot q_1^{5}\cdot q_2^{2}\\
&\quad +48\cdot q_0^{1}\cdot q_1^{5}\cdot q_3^{2}-48\cdot q_0^{2}\cdot q_2^{1}\cdot q_3^{5}+160\cdot q_0^{2}\cdot q_2^{3}\cdot q_3^{3}\\
&\quad -48\cdot q_0^{2}\cdot q_2^{5}\cdot q_3^{1}+160\cdot q_0^{3}\cdot q_1^{3}\cdot q_2^{2}-160\cdot q_0^{3}\cdot q_1^{3}\cdot q_3^{2}\\
&\quad -48\cdot q_0^{5}\cdot q_1^{1}\cdot q_2^{2}+48\cdot q_0^{5}\cdot q_1^{1}\cdot q_3^{2}+16\cdot q_0^{6}\cdot q_2^{1}\cdot q_3^{1}\\
&\quad +48\cdot q_1^{2}\cdot q_2^{1}\cdot q_3^{5}-160\cdot q_1^{2}\cdot q_2^{3}\cdot q_3^{3}+48\cdot q_1^{2}\cdot q_2^{5}\cdot q_3^{1}\\
&\quad -16\cdot q_1^{6}\cdot q_2^{1}\cdot q_3^{1}+240\cdot q_0^{1}\cdot q_1^{1}\cdot q_2^{2}\cdot q_3^{4}-240\cdot q_0^{1}\cdot q_1^{1}\cdot q_2^{4}\cdot q_3^{2}\\
&\quad +240\cdot q_0^{2}\cdot q_1^{4}\cdot q_2^{1}\cdot q_3^{1}-240\cdot q_0^{4}\cdot q_1^{2}\cdot q_2^{1}\cdot q_3^{1}\end{align*}
\begin{align*}(T_{h}^4)_{5,3}&=-2\cdot q_0^{1}\cdot q_3^{7}+6\cdot q_0^{3}\cdot q_3^{5}+6\cdot q_0^{5}\cdot q_3^{3}-2\cdot q_0^{7}\cdot q_3^{1}\\
&\quad +2\cdot q_1^{1}\cdot q_2^{7}-6\cdot q_1^{3}\cdot q_2^{5}-6\cdot q_1^{5}\cdot q_2^{3}+2\cdot q_1^{7}\cdot q_2^{1}\\
&\quad +6\cdot q_0^{1}\cdot q_1^{2}\cdot q_3^{5}+30\cdot q_0^{1}\cdot q_1^{4}\cdot q_3^{3}-10\cdot q_0^{1}\cdot q_1^{6}\cdot q_3^{1}\\
&\quad +18\cdot q_0^{1}\cdot q_2^{2}\cdot q_3^{5}+10\cdot q_0^{1}\cdot q_2^{4}\cdot q_3^{3}-10\cdot q_0^{1}\cdot q_2^{6}\cdot q_3^{1}\\
&\quad -6\cdot q_0^{2}\cdot q_1^{1}\cdot q_2^{5}+60\cdot q_0^{2}\cdot q_1^{3}\cdot q_2^{3}-18\cdot q_0^{2}\cdot q_1^{5}\cdot q_2^{1}\\
&\quad -60\cdot q_0^{3}\cdot q_1^{2}\cdot q_3^{3}+10\cdot q_0^{3}\cdot q_1^{4}\cdot q_3^{1}-60\cdot q_0^{3}\cdot q_2^{2}\cdot q_3^{3}\\
&\quad +30\cdot q_0^{3}\cdot q_2^{4}\cdot q_3^{1}-30\cdot q_0^{4}\cdot q_1^{1}\cdot q_2^{3}-10\cdot q_0^{4}\cdot q_1^{3}\cdot q_2^{1}\\
&\quad +18\cdot q_0^{5}\cdot q_1^{2}\cdot q_3^{1}+6\cdot q_0^{5}\cdot q_2^{2}\cdot q_3^{1}+10\cdot q_0^{6}\cdot q_1^{1}\cdot q_2^{1}\\
&\quad +10\cdot q_1^{1}\cdot q_2^{1}\cdot q_3^{6}-10\cdot q_1^{1}\cdot q_2^{3}\cdot q_3^{4}-18\cdot q_1^{1}\cdot q_2^{5}\cdot q_3^{2}\\
&\quad -30\cdot q_1^{3}\cdot q_2^{1}\cdot q_3^{4}+60\cdot q_1^{3}\cdot q_2^{3}\cdot q_3^{2}-6\cdot q_1^{5}\cdot q_2^{1}\cdot q_3^{2}\\
&\quad -60\cdot q_0^{1}\cdot q_1^{2}\cdot q_2^{2}\cdot q_3^{3}+30\cdot q_0^{1}\cdot q_1^{2}\cdot q_2^{4}\cdot q_3^{1}\\
&\quad +30\cdot q_0^{1}\cdot q_1^{4}\cdot q_2^{2}\cdot q_3^{1}-30\cdot q_0^{2}\cdot q_1^{1}\cdot q_2^{1}\cdot q_3^{4}\\
&\quad +60\cdot q_0^{2}\cdot q_1^{1}\cdot q_2^{3}\cdot q_3^{2}+60\cdot q_0^{2}\cdot q_1^{3}\cdot q_2^{1}\cdot q_3^{2}\\
&\quad -60\cdot q_0^{3}\cdot q_1^{2}\cdot q_2^{2}\cdot q_3^{1}-30\cdot q_0^{4}\cdot q_1^{1}\cdot q_2^{1}\cdot q_3^{2}\end{align*}
\begin{align*}(T_{h}^4)_{5,4}&=-2\cdot q_0^{1}\cdot q_2^{7}+6\cdot q_0^{3}\cdot q_2^{5}+6\cdot q_0^{5}\cdot q_2^{3}-2\cdot q_0^{7}\cdot q_2^{1}\\
&\quad -2\cdot q_1^{1}\cdot q_3^{7}+6\cdot q_1^{3}\cdot q_3^{5}+6\cdot q_1^{5}\cdot q_3^{3}-2\cdot q_1^{7}\cdot q_3^{1}\\
&\quad +6\cdot q_0^{1}\cdot q_1^{2}\cdot q_2^{5}+30\cdot q_0^{1}\cdot q_1^{4}\cdot q_2^{3}-10\cdot q_0^{1}\cdot q_1^{6}\cdot q_2^{1}\\
&\quad -10\cdot q_0^{1}\cdot q_2^{1}\cdot q_3^{6}+10\cdot q_0^{1}\cdot q_2^{3}\cdot q_3^{4}+18\cdot q_0^{1}\cdot q_2^{5}\cdot q_3^{2}\\
&\quad +6\cdot q_0^{2}\cdot q_1^{1}\cdot q_3^{5}-60\cdot q_0^{2}\cdot q_1^{3}\cdot q_3^{3}+18\cdot q_0^{2}\cdot q_1^{5}\cdot q_3^{1}\\
&\quad -60\cdot q_0^{3}\cdot q_1^{2}\cdot q_2^{3}+10\cdot q_0^{3}\cdot q_1^{4}\cdot q_2^{1}+30\cdot q_0^{3}\cdot q_2^{1}\cdot q_3^{4}\\
&\quad -60\cdot q_0^{3}\cdot q_2^{3}\cdot q_3^{2}+30\cdot q_0^{4}\cdot q_1^{1}\cdot q_3^{3}+10\cdot q_0^{4}\cdot q_1^{3}\cdot q_3^{1}\\
&\quad +18\cdot q_0^{5}\cdot q_1^{2}\cdot q_2^{1}+6\cdot q_0^{5}\cdot q_2^{1}\cdot q_3^{2}-10\cdot q_0^{6}\cdot q_1^{1}\cdot q_3^{1}\\
&\quad +18\cdot q_1^{1}\cdot q_2^{2}\cdot q_3^{5}+10\cdot q_1^{1}\cdot q_2^{4}\cdot q_3^{3}-10\cdot q_1^{1}\cdot q_2^{6}\cdot q_3^{1}\\
&\quad -60\cdot q_1^{3}\cdot q_2^{2}\cdot q_3^{3}+30\cdot q_1^{3}\cdot q_2^{4}\cdot q_3^{1}+6\cdot q_1^{5}\cdot q_2^{2}\cdot q_3^{1}\\
&\quad +30\cdot q_0^{1}\cdot q_1^{2}\cdot q_2^{1}\cdot q_3^{4}-60\cdot q_0^{1}\cdot q_1^{2}\cdot q_2^{3}\cdot q_3^{2}\\
&\quad +30\cdot q_0^{1}\cdot q_1^{4}\cdot q_2^{1}\cdot q_3^{2}-60\cdot q_0^{2}\cdot q_1^{1}\cdot q_2^{2}\cdot q_3^{3}\\
&\quad +30\cdot q_0^{2}\cdot q_1^{1}\cdot q_2^{4}\cdot q_3^{1}-60\cdot q_0^{2}\cdot q_1^{3}\cdot q_2^{2}\cdot q_3^{1}\\
&\quad -60\cdot q_0^{3}\cdot q_1^{2}\cdot q_2^{1}\cdot q_3^{2}+30\cdot q_0^{4}\cdot q_1^{1}\cdot q_2^{2}\cdot q_3^{1}\end{align*}
\begin{align*}(T_{h}^4)_{5,5}&=1\cdot q_0^{8}+1\cdot q_1^{8}-1\cdot q_2^{8}-1\cdot q_3^{8}-4\cdot q_0^{2}\cdot q_1^{6}+12\cdot q_0^{2}\cdot q_2^{6}\\
&\quad +12\cdot q_0^{2}\cdot q_3^{6}-10\cdot q_0^{4}\cdot q_1^{4}-4\cdot q_0^{6}\cdot q_1^{2}-12\cdot q_0^{6}\cdot q_2^{2}\\
&\quad -12\cdot q_0^{6}\cdot q_3^{2}+12\cdot q_1^{2}\cdot q_2^{6}+12\cdot q_1^{2}\cdot q_3^{6}-12\cdot q_1^{6}\cdot q_2^{2}\\
&\quad -12\cdot q_1^{6}\cdot q_3^{2}+4\cdot q_2^{2}\cdot q_3^{6}+10\cdot q_2^{4}\cdot q_3^{4}+4\cdot q_2^{6}\cdot q_3^{2}\\
&\quad -120\cdot q_0^{2}\cdot q_1^{2}\cdot q_2^{4}-120\cdot q_0^{2}\cdot q_1^{2}\cdot q_3^{4}+60\cdot q_0^{2}\cdot q_1^{4}\cdot q_2^{2}\\
&\quad +60\cdot q_0^{2}\cdot q_1^{4}\cdot q_3^{2}-60\cdot q_0^{2}\cdot q_2^{2}\cdot q_3^{4}-60\cdot q_0^{2}\cdot q_2^{4}\cdot q_3^{2}\\
&\quad +60\cdot q_0^{4}\cdot q_1^{2}\cdot q_2^{2}+60\cdot q_0^{4}\cdot q_1^{2}\cdot q_3^{2}+120\cdot q_0^{4}\cdot q_2^{2}\cdot q_3^{2}\\
&\quad -60\cdot q_1^{2}\cdot q_2^{2}\cdot q_3^{4}-60\cdot q_1^{2}\cdot q_2^{4}\cdot q_3^{2}+120\cdot q_1^{4}\cdot q_2^{2}\cdot q_3^{2}\end{align*}
\begin{align*}(T_{h}^4)_{5,6}&=-8\cdot q_0^{1}\cdot q_1^{7}-8\cdot q_0^{3}\cdot q_1^{5}+8\cdot q_0^{5}\cdot q_1^{3}+8\cdot q_0^{7}\cdot q_1^{1}\\
&\quad -8\cdot q_2^{1}\cdot q_3^{7}-8\cdot q_2^{3}\cdot q_3^{5}+8\cdot q_2^{5}\cdot q_3^{3}+8\cdot q_2^{7}\cdot q_3^{1}\\
&\quad -120\cdot q_0^{1}\cdot q_1^{3}\cdot q_2^{4}-120\cdot q_0^{1}\cdot q_1^{3}\cdot q_3^{4}+96\cdot q_0^{1}\cdot q_1^{5}\cdot q_2^{2}\\
&\quad +96\cdot q_0^{1}\cdot q_1^{5}\cdot q_3^{2}+96\cdot q_0^{2}\cdot q_2^{1}\cdot q_3^{5}-96\cdot q_0^{2}\cdot q_2^{5}\cdot q_3^{1}\\
&\quad +120\cdot q_0^{3}\cdot q_1^{1}\cdot q_2^{4}+120\cdot q_0^{3}\cdot q_1^{1}\cdot q_3^{4}-120\cdot q_0^{4}\cdot q_2^{1}\cdot q_3^{3}\\
&\quad +120\cdot q_0^{4}\cdot q_2^{3}\cdot q_3^{1}-96\cdot q_0^{5}\cdot q_1^{1}\cdot q_2^{2}-96\cdot q_0^{5}\cdot q_1^{1}\cdot q_3^{2}\\
&\quad +96\cdot q_1^{2}\cdot q_2^{1}\cdot q_3^{5}-96\cdot q_1^{2}\cdot q_2^{5}\cdot q_3^{1}-120\cdot q_1^{4}\cdot q_2^{1}\cdot q_3^{3}\\
&\quad +120\cdot q_1^{4}\cdot q_2^{3}\cdot q_3^{1}-240\cdot q_0^{1}\cdot q_1^{3}\cdot q_2^{2}\cdot q_3^{2}-240\cdot q_0^{2}\cdot q_1^{2}\cdot q_2^{1}\cdot q_3^{3}\\
&\quad +240\cdot q_0^{2}\cdot q_1^{2}\cdot q_2^{3}\cdot q_3^{1}+240\cdot q_0^{3}\cdot q_1^{1}\cdot q_2^{2}\cdot q_3^{2}\end{align*}
\begin{align*}(T_{h}^4)_{5,7}&=6\cdot q_0^{1}\cdot q_3^{7}-10\cdot q_0^{3}\cdot q_3^{5}-10\cdot q_0^{5}\cdot q_3^{3}+6\cdot q_0^{7}\cdot q_3^{1}\\
&\quad -6\cdot q_1^{1}\cdot q_2^{7}+10\cdot q_1^{3}\cdot q_2^{5}+10\cdot q_1^{5}\cdot q_2^{3}-6\cdot q_1^{7}\cdot q_2^{1}\\
&\quad -90\cdot q_0^{1}\cdot q_1^{2}\cdot q_3^{5}+110\cdot q_0^{1}\cdot q_1^{4}\cdot q_3^{3}-18\cdot q_0^{1}\cdot q_1^{6}\cdot q_3^{1}\\
&\quad -6\cdot q_0^{1}\cdot q_2^{2}\cdot q_3^{5}-30\cdot q_0^{1}\cdot q_2^{4}\cdot q_3^{3}-18\cdot q_0^{1}\cdot q_2^{6}\cdot q_3^{1}\\
&\quad +90\cdot q_0^{2}\cdot q_1^{1}\cdot q_2^{5}-100\cdot q_0^{2}\cdot q_1^{3}\cdot q_2^{3}+6\cdot q_0^{2}\cdot q_1^{5}\cdot q_2^{1}\\
&\quad +100\cdot q_0^{3}\cdot q_1^{2}\cdot q_3^{3}-30\cdot q_0^{3}\cdot q_1^{4}\cdot q_3^{1}+100\cdot q_0^{3}\cdot q_2^{2}\cdot q_3^{3}\\
&\quad +110\cdot q_0^{3}\cdot q_2^{4}\cdot q_3^{1}-110\cdot q_0^{4}\cdot q_1^{1}\cdot q_2^{3}+30\cdot q_0^{4}\cdot q_1^{3}\cdot q_2^{1}\\
&\quad -6\cdot q_0^{5}\cdot q_1^{2}\cdot q_3^{1}-90\cdot q_0^{5}\cdot q_2^{2}\cdot q_3^{1}+18\cdot q_0^{6}\cdot q_1^{1}\cdot q_2^{1}\\
&\quad +18\cdot q_1^{1}\cdot q_2^{1}\cdot q_3^{6}+30\cdot q_1^{1}\cdot q_2^{3}\cdot q_3^{4}+6\cdot q_1^{1}\cdot q_2^{5}\cdot q_3^{2}\\
&\quad -110\cdot q_1^{3}\cdot q_2^{1}\cdot q_3^{4}-100\cdot q_1^{3}\cdot q_2^{3}\cdot q_3^{2}+90\cdot q_1^{5}\cdot q_2^{1}\cdot q_3^{2}\\
&\quad -60\cdot q_0^{1}\cdot q_1^{2}\cdot q_2^{2}\cdot q_3^{3}+30\cdot q_0^{1}\cdot q_1^{2}\cdot q_2^{4}\cdot q_3^{1}\\
&\quad +30\cdot q_0^{1}\cdot q_1^{4}\cdot q_2^{2}\cdot q_3^{1}-30\cdot q_0^{2}\cdot q_1^{1}\cdot q_2^{1}\cdot q_3^{4}\\
&\quad +60\cdot q_0^{2}\cdot q_1^{1}\cdot q_2^{3}\cdot q_3^{2}+60\cdot q_0^{2}\cdot q_1^{3}\cdot q_2^{1}\cdot q_3^{2}\\
&\quad -60\cdot q_0^{3}\cdot q_1^{2}\cdot q_2^{2}\cdot q_3^{1}-30\cdot q_0^{4}\cdot q_1^{1}\cdot q_2^{1}\cdot q_3^{2}\end{align*}
\begin{align*}(T_{h}^4)_{5,8}&=-6\cdot q_0^{1}\cdot q_2^{7}+10\cdot q_0^{3}\cdot q_2^{5}+10\cdot q_0^{5}\cdot q_2^{3}-6\cdot q_0^{7}\cdot q_2^{1}\\
&\quad -6\cdot q_1^{1}\cdot q_3^{7}+10\cdot q_1^{3}\cdot q_3^{5}+10\cdot q_1^{5}\cdot q_3^{3}-6\cdot q_1^{7}\cdot q_3^{1}\\
&\quad +90\cdot q_0^{1}\cdot q_1^{2}\cdot q_2^{5}-110\cdot q_0^{1}\cdot q_1^{4}\cdot q_2^{3}+18\cdot q_0^{1}\cdot q_1^{6}\cdot q_2^{1}\\
&\quad +18\cdot q_0^{1}\cdot q_2^{1}\cdot q_3^{6}+30\cdot q_0^{1}\cdot q_2^{3}\cdot q_3^{4}+6\cdot q_0^{1}\cdot q_2^{5}\cdot q_3^{2}\\
&\quad +90\cdot q_0^{2}\cdot q_1^{1}\cdot q_3^{5}-100\cdot q_0^{2}\cdot q_1^{3}\cdot q_3^{3}+6\cdot q_0^{2}\cdot q_1^{5}\cdot q_3^{1}\\
&\quad -100\cdot q_0^{3}\cdot q_1^{2}\cdot q_2^{3}+30\cdot q_0^{3}\cdot q_1^{4}\cdot q_2^{1}-110\cdot q_0^{3}\cdot q_2^{1}\cdot q_3^{4}\\
&\quad -100\cdot q_0^{3}\cdot q_2^{3}\cdot q_3^{2}-110\cdot q_0^{4}\cdot q_1^{1}\cdot q_3^{3}+30\cdot q_0^{4}\cdot q_1^{3}\cdot q_3^{1}\\
&\quad +6\cdot q_0^{5}\cdot q_1^{2}\cdot q_2^{1}+90\cdot q_0^{5}\cdot q_2^{1}\cdot q_3^{2}+18\cdot q_0^{6}\cdot q_1^{1}\cdot q_3^{1}\\
&\quad +6\cdot q_1^{1}\cdot q_2^{2}\cdot q_3^{5}+30\cdot q_1^{1}\cdot q_2^{4}\cdot q_3^{3}+18\cdot q_1^{1}\cdot q_2^{6}\cdot q_3^{1}\\
&\quad -100\cdot q_1^{3}\cdot q_2^{2}\cdot q_3^{3}-110\cdot q_1^{3}\cdot q_2^{4}\cdot q_3^{1}+90\cdot q_1^{5}\cdot q_2^{2}\cdot q_3^{1}\\
&\quad -30\cdot q_0^{1}\cdot q_1^{2}\cdot q_2^{1}\cdot q_3^{4}+60\cdot q_0^{1}\cdot q_1^{2}\cdot q_2^{3}\cdot q_3^{2}\\
&\quad -30\cdot q_0^{1}\cdot q_1^{4}\cdot q_2^{1}\cdot q_3^{2}+60\cdot q_0^{2}\cdot q_1^{1}\cdot q_2^{2}\cdot q_3^{3}\\
&\quad -30\cdot q_0^{2}\cdot q_1^{1}\cdot q_2^{4}\cdot q_3^{1}+60\cdot q_0^{2}\cdot q_1^{3}\cdot q_2^{2}\cdot q_3^{1}\\
&\quad +60\cdot q_0^{3}\cdot q_1^{2}\cdot q_2^{1}\cdot q_3^{2}-30\cdot q_0^{4}\cdot q_1^{1}\cdot q_2^{2}\cdot q_3^{1}\end{align*}
\begin{align*}(T_{h}^4)_{5,9}&=-240\cdot q_0^{1}\cdot q_1^{1}\cdot q_2^{6}+240\cdot q_0^{1}\cdot q_1^{1}\cdot q_3^{6}+640\cdot q_0^{1}\cdot q_1^{3}\cdot q_2^{4}\\
&\quad -640\cdot q_0^{1}\cdot q_1^{3}\cdot q_3^{4}-240\cdot q_0^{1}\cdot q_1^{5}\cdot q_2^{2}+240\cdot q_0^{1}\cdot q_1^{5}\cdot q_3^{2}\\
&\quad -240\cdot q_0^{2}\cdot q_2^{1}\cdot q_3^{5}-480\cdot q_0^{2}\cdot q_2^{3}\cdot q_3^{3}-240\cdot q_0^{2}\cdot q_2^{5}\cdot q_3^{1}\\
&\quad +640\cdot q_0^{3}\cdot q_1^{1}\cdot q_2^{4}-640\cdot q_0^{3}\cdot q_1^{1}\cdot q_3^{4}-480\cdot q_0^{3}\cdot q_1^{3}\cdot q_2^{2}\\
&\quad +480\cdot q_0^{3}\cdot q_1^{3}\cdot q_3^{2}+640\cdot q_0^{4}\cdot q_2^{1}\cdot q_3^{3}+640\cdot q_0^{4}\cdot q_2^{3}\cdot q_3^{1}\\
&\quad -240\cdot q_0^{5}\cdot q_1^{1}\cdot q_2^{2}+240\cdot q_0^{5}\cdot q_1^{1}\cdot q_3^{2}-240\cdot q_0^{6}\cdot q_2^{1}\cdot q_3^{1}\\
&\quad +240\cdot q_1^{2}\cdot q_2^{1}\cdot q_3^{5}+480\cdot q_1^{2}\cdot q_2^{3}\cdot q_3^{3}+240\cdot q_1^{2}\cdot q_2^{5}\cdot q_3^{1}\\
&\quad -640\cdot q_1^{4}\cdot q_2^{1}\cdot q_3^{3}-640\cdot q_1^{4}\cdot q_2^{3}\cdot q_3^{1}+240\cdot q_1^{6}\cdot q_2^{1}\cdot q_3^{1}\\
&\quad +240\cdot q_0^{1}\cdot q_1^{1}\cdot q_2^{2}\cdot q_3^{4}-240\cdot q_0^{1}\cdot q_1^{1}\cdot q_2^{4}\cdot q_3^{2}\\
&\quad +240\cdot q_0^{2}\cdot q_1^{4}\cdot q_2^{1}\cdot q_3^{1}-240\cdot q_0^{4}\cdot q_1^{2}\cdot q_2^{1}\cdot q_3^{1}\end{align*}
\begin{align*}(T_{h}^4)_{6,1}&=2\cdot q_0^{1}\cdot q_1^{1}\cdot q_2^{6}-2\cdot q_0^{1}\cdot q_1^{1}\cdot q_3^{6}+6\cdot q_0^{1}\cdot q_1^{5}\cdot q_2^{2}\\
&\quad -6\cdot q_0^{1}\cdot q_1^{5}\cdot q_3^{2}-6\cdot q_0^{2}\cdot q_2^{1}\cdot q_3^{5}+20\cdot q_0^{2}\cdot q_2^{3}\cdot q_3^{3}\\
&\quad -6\cdot q_0^{2}\cdot q_2^{5}\cdot q_3^{1}-20\cdot q_0^{3}\cdot q_1^{3}\cdot q_2^{2}+20\cdot q_0^{3}\cdot q_1^{3}\cdot q_3^{2}\\
&\quad +6\cdot q_0^{5}\cdot q_1^{1}\cdot q_2^{2}-6\cdot q_0^{5}\cdot q_1^{1}\cdot q_3^{2}-2\cdot q_0^{6}\cdot q_2^{1}\cdot q_3^{1}\\
&\quad +6\cdot q_1^{2}\cdot q_2^{1}\cdot q_3^{5}-20\cdot q_1^{2}\cdot q_2^{3}\cdot q_3^{3}+6\cdot q_1^{2}\cdot q_2^{5}\cdot q_3^{1}\\
&\quad +2\cdot q_1^{6}\cdot q_2^{1}\cdot q_3^{1}+30\cdot q_0^{1}\cdot q_1^{1}\cdot q_2^{2}\cdot q_3^{4}-30\cdot q_0^{1}\cdot q_1^{1}\cdot q_2^{4}\cdot q_3^{2}\\
&\quad -30\cdot q_0^{2}\cdot q_1^{4}\cdot q_2^{1}\cdot q_3^{1}+30\cdot q_0^{4}\cdot q_1^{2}\cdot q_2^{1}\cdot q_3^{1}\end{align*}
\begin{align*}(T_{h}^4)_{6,2}&=-4\cdot q_0^{2}\cdot q_2^{6}+4\cdot q_0^{2}\cdot q_3^{6}-4\cdot q_0^{6}\cdot q_2^{2}+4\cdot q_0^{6}\cdot q_3^{2}\\
&\quad +4\cdot q_1^{2}\cdot q_2^{6}-4\cdot q_1^{2}\cdot q_3^{6}+4\cdot q_1^{6}\cdot q_2^{2}-4\cdot q_1^{6}\cdot q_3^{2}\\
&\quad -60\cdot q_0^{2}\cdot q_1^{4}\cdot q_2^{2}+60\cdot q_0^{2}\cdot q_1^{4}\cdot q_3^{2}-60\cdot q_0^{2}\cdot q_2^{2}\cdot q_3^{4}\\
&\quad +60\cdot q_0^{2}\cdot q_2^{4}\cdot q_3^{2}+60\cdot q_0^{4}\cdot q_1^{2}\cdot q_2^{2}-60\cdot q_0^{4}\cdot q_1^{2}\cdot q_3^{2}\\
&\quad +60\cdot q_1^{2}\cdot q_2^{2}\cdot q_3^{4}-60\cdot q_1^{2}\cdot q_2^{4}\cdot q_3^{2}-48\cdot q_0^{1}\cdot q_1^{1}\cdot q_2^{1}\cdot q_3^{5}\\
&\quad +160\cdot q_0^{1}\cdot q_1^{1}\cdot q_2^{3}\cdot q_3^{3}-48\cdot q_0^{1}\cdot q_1^{1}\cdot q_2^{5}\cdot q_3^{1}\\
&\quad -48\cdot q_0^{1}\cdot q_1^{5}\cdot q_2^{1}\cdot q_3^{1}+160\cdot q_0^{3}\cdot q_1^{3}\cdot q_2^{1}\cdot q_3^{1}\\
&\quad -48\cdot q_0^{5}\cdot q_1^{1}\cdot q_2^{1}\cdot q_3^{1}\end{align*}
\begin{align*}(T_{h}^4)_{6,3}&=-1\cdot q_0^{1}\cdot q_2^{7}+3\cdot q_0^{3}\cdot q_2^{5}-3\cdot q_0^{5}\cdot q_2^{3}+1\cdot q_0^{7}\cdot q_2^{1}\\
&\quad -1\cdot q_1^{1}\cdot q_3^{7}+3\cdot q_1^{3}\cdot q_3^{5}-3\cdot q_1^{5}\cdot q_3^{3}+1\cdot q_1^{7}\cdot q_3^{1}\\
&\quad +3\cdot q_0^{1}\cdot q_1^{2}\cdot q_2^{5}-15\cdot q_0^{1}\cdot q_1^{4}\cdot q_2^{3}+5\cdot q_0^{1}\cdot q_1^{6}\cdot q_2^{1}\\
&\quad -5\cdot q_0^{1}\cdot q_2^{1}\cdot q_3^{6}+5\cdot q_0^{1}\cdot q_2^{3}\cdot q_3^{4}+9\cdot q_0^{1}\cdot q_2^{5}\cdot q_3^{2}\\
&\quad +3\cdot q_0^{2}\cdot q_1^{1}\cdot q_3^{5}+30\cdot q_0^{2}\cdot q_1^{3}\cdot q_3^{3}-9\cdot q_0^{2}\cdot q_1^{5}\cdot q_3^{1}\\
&\quad +30\cdot q_0^{3}\cdot q_1^{2}\cdot q_2^{3}-5\cdot q_0^{3}\cdot q_1^{4}\cdot q_2^{1}+15\cdot q_0^{3}\cdot q_2^{1}\cdot q_3^{4}\\
&\quad -30\cdot q_0^{3}\cdot q_2^{3}\cdot q_3^{2}-15\cdot q_0^{4}\cdot q_1^{1}\cdot q_3^{3}-5\cdot q_0^{4}\cdot q_1^{3}\cdot q_3^{1}\\
&\quad -9\cdot q_0^{5}\cdot q_1^{2}\cdot q_2^{1}-3\cdot q_0^{5}\cdot q_2^{1}\cdot q_3^{2}+5\cdot q_0^{6}\cdot q_1^{1}\cdot q_3^{1}\\
&\quad +9\cdot q_1^{1}\cdot q_2^{2}\cdot q_3^{5}+5\cdot q_1^{1}\cdot q_2^{4}\cdot q_3^{3}-5\cdot q_1^{1}\cdot q_2^{6}\cdot q_3^{1}\\
&\quad -30\cdot q_1^{3}\cdot q_2^{2}\cdot q_3^{3}+15\cdot q_1^{3}\cdot q_2^{4}\cdot q_3^{1}-3\cdot q_1^{5}\cdot q_2^{2}\cdot q_3^{1}\\
&\quad +15\cdot q_0^{1}\cdot q_1^{2}\cdot q_2^{1}\cdot q_3^{4}-30\cdot q_0^{1}\cdot q_1^{2}\cdot q_2^{3}\cdot q_3^{2}\\
&\quad -15\cdot q_0^{1}\cdot q_1^{4}\cdot q_2^{1}\cdot q_3^{2}-30\cdot q_0^{2}\cdot q_1^{1}\cdot q_2^{2}\cdot q_3^{3}\\
&\quad +15\cdot q_0^{2}\cdot q_1^{1}\cdot q_2^{4}\cdot q_3^{1}+30\cdot q_0^{2}\cdot q_1^{3}\cdot q_2^{2}\cdot q_3^{1}\\
&\quad +30\cdot q_0^{3}\cdot q_1^{2}\cdot q_2^{1}\cdot q_3^{2}-15\cdot q_0^{4}\cdot q_1^{1}\cdot q_2^{2}\cdot q_3^{1}\end{align*}
\begin{align*}(T_{h}^4)_{6,4}&=1\cdot q_0^{1}\cdot q_3^{7}-3\cdot q_0^{3}\cdot q_3^{5}+3\cdot q_0^{5}\cdot q_3^{3}-1\cdot q_0^{7}\cdot q_3^{1}\\
&\quad -1\cdot q_1^{1}\cdot q_2^{7}+3\cdot q_1^{3}\cdot q_2^{5}-3\cdot q_1^{5}\cdot q_2^{3}+1\cdot q_1^{7}\cdot q_2^{1}\\
&\quad -3\cdot q_0^{1}\cdot q_1^{2}\cdot q_3^{5}+15\cdot q_0^{1}\cdot q_1^{4}\cdot q_3^{3}-5\cdot q_0^{1}\cdot q_1^{6}\cdot q_3^{1}\\
&\quad -9\cdot q_0^{1}\cdot q_2^{2}\cdot q_3^{5}-5\cdot q_0^{1}\cdot q_2^{4}\cdot q_3^{3}+5\cdot q_0^{1}\cdot q_2^{6}\cdot q_3^{1}\\
&\quad +3\cdot q_0^{2}\cdot q_1^{1}\cdot q_2^{5}+30\cdot q_0^{2}\cdot q_1^{3}\cdot q_2^{3}-9\cdot q_0^{2}\cdot q_1^{5}\cdot q_2^{1}\\
&\quad -30\cdot q_0^{3}\cdot q_1^{2}\cdot q_3^{3}+5\cdot q_0^{3}\cdot q_1^{4}\cdot q_3^{1}+30\cdot q_0^{3}\cdot q_2^{2}\cdot q_3^{3}\\
&\quad -15\cdot q_0^{3}\cdot q_2^{4}\cdot q_3^{1}-15\cdot q_0^{4}\cdot q_1^{1}\cdot q_2^{3}-5\cdot q_0^{4}\cdot q_1^{3}\cdot q_2^{1}\\
&\quad +9\cdot q_0^{5}\cdot q_1^{2}\cdot q_3^{1}+3\cdot q_0^{5}\cdot q_2^{2}\cdot q_3^{1}+5\cdot q_0^{6}\cdot q_1^{1}\cdot q_2^{1}\\
&\quad -5\cdot q_1^{1}\cdot q_2^{1}\cdot q_3^{6}+5\cdot q_1^{1}\cdot q_2^{3}\cdot q_3^{4}+9\cdot q_1^{1}\cdot q_2^{5}\cdot q_3^{2}\\
&\quad +15\cdot q_1^{3}\cdot q_2^{1}\cdot q_3^{4}-30\cdot q_1^{3}\cdot q_2^{3}\cdot q_3^{2}-3\cdot q_1^{5}\cdot q_2^{1}\cdot q_3^{2}\\
&\quad +30\cdot q_0^{1}\cdot q_1^{2}\cdot q_2^{2}\cdot q_3^{3}-15\cdot q_0^{1}\cdot q_1^{2}\cdot q_2^{4}\cdot q_3^{1}\\
&\quad +15\cdot q_0^{1}\cdot q_1^{4}\cdot q_2^{2}\cdot q_3^{1}+15\cdot q_0^{2}\cdot q_1^{1}\cdot q_2^{1}\cdot q_3^{4}\\
&\quad -30\cdot q_0^{2}\cdot q_1^{1}\cdot q_2^{3}\cdot q_3^{2}+30\cdot q_0^{2}\cdot q_1^{3}\cdot q_2^{1}\cdot q_3^{2}\\
&\quad -30\cdot q_0^{3}\cdot q_1^{2}\cdot q_2^{2}\cdot q_3^{1}-15\cdot q_0^{4}\cdot q_1^{1}\cdot q_2^{1}\cdot q_3^{2}\end{align*}
\begin{align*}(T_{h}^4)_{6,5}&=2\cdot q_0^{1}\cdot q_1^{7}+2\cdot q_0^{3}\cdot q_1^{5}-2\cdot q_0^{5}\cdot q_1^{3}-2\cdot q_0^{7}\cdot q_1^{1}\\
&\quad -2\cdot q_2^{1}\cdot q_3^{7}-2\cdot q_2^{3}\cdot q_3^{5}+2\cdot q_2^{5}\cdot q_3^{3}+2\cdot q_2^{7}\cdot q_3^{1}\\
&\quad +30\cdot q_0^{1}\cdot q_1^{3}\cdot q_2^{4}+30\cdot q_0^{1}\cdot q_1^{3}\cdot q_3^{4}-24\cdot q_0^{1}\cdot q_1^{5}\cdot q_2^{2}\\
&\quad -24\cdot q_0^{1}\cdot q_1^{5}\cdot q_3^{2}+24\cdot q_0^{2}\cdot q_2^{1}\cdot q_3^{5}-24\cdot q_0^{2}\cdot q_2^{5}\cdot q_3^{1}\\
&\quad -30\cdot q_0^{3}\cdot q_1^{1}\cdot q_2^{4}-30\cdot q_0^{3}\cdot q_1^{1}\cdot q_3^{4}-30\cdot q_0^{4}\cdot q_2^{1}\cdot q_3^{3}\\
&\quad +30\cdot q_0^{4}\cdot q_2^{3}\cdot q_3^{1}+24\cdot q_0^{5}\cdot q_1^{1}\cdot q_2^{2}+24\cdot q_0^{5}\cdot q_1^{1}\cdot q_3^{2}\\
&\quad +24\cdot q_1^{2}\cdot q_2^{1}\cdot q_3^{5}-24\cdot q_1^{2}\cdot q_2^{5}\cdot q_3^{1}-30\cdot q_1^{4}\cdot q_2^{1}\cdot q_3^{3}\\
&\quad +30\cdot q_1^{4}\cdot q_2^{3}\cdot q_3^{1}+60\cdot q_0^{1}\cdot q_1^{3}\cdot q_2^{2}\cdot q_3^{2}-60\cdot q_0^{2}\cdot q_1^{2}\cdot q_2^{1}\cdot q_3^{3}\\
&\quad +60\cdot q_0^{2}\cdot q_1^{2}\cdot q_2^{3}\cdot q_3^{1}-60\cdot q_0^{3}\cdot q_1^{1}\cdot q_2^{2}\cdot q_3^{2}\end{align*}
\begin{align*}(T_{h}^4)_{6,6}&=1\cdot q_0^{8}+1\cdot q_1^{8}+1\cdot q_2^{8}+1\cdot q_3^{8}-4\cdot q_0^{2}\cdot q_1^{6}-12\cdot q_0^{2}\cdot q_2^{6}\\
&\quad -12\cdot q_0^{2}\cdot q_3^{6}-10\cdot q_0^{4}\cdot q_1^{4}+30\cdot q_0^{4}\cdot q_2^{4}+30\cdot q_0^{4}\cdot q_3^{4}\\
&\quad -4\cdot q_0^{6}\cdot q_1^{2}-12\cdot q_0^{6}\cdot q_2^{2}-12\cdot q_0^{6}\cdot q_3^{2}-12\cdot q_1^{2}\cdot q_2^{6}\\
&\quad -12\cdot q_1^{2}\cdot q_3^{6}+30\cdot q_1^{4}\cdot q_2^{4}+30\cdot q_1^{4}\cdot q_3^{4}-12\cdot q_1^{6}\cdot q_2^{2}\\
&\quad -12\cdot q_1^{6}\cdot q_3^{2}-4\cdot q_2^{2}\cdot q_3^{6}-10\cdot q_2^{4}\cdot q_3^{4}-4\cdot q_2^{6}\cdot q_3^{2}\\
&\quad -60\cdot q_0^{2}\cdot q_1^{2}\cdot q_2^{4}-60\cdot q_0^{2}\cdot q_1^{2}\cdot q_3^{4}+60\cdot q_0^{2}\cdot q_1^{4}\cdot q_2^{2}\\
&\quad +60\cdot q_0^{2}\cdot q_1^{4}\cdot q_3^{2}+60\cdot q_0^{2}\cdot q_2^{2}\cdot q_3^{4}+60\cdot q_0^{2}\cdot q_2^{4}\cdot q_3^{2}\\
&\quad +60\cdot q_0^{4}\cdot q_1^{2}\cdot q_2^{2}+60\cdot q_0^{4}\cdot q_1^{2}\cdot q_3^{2}-60\cdot q_0^{4}\cdot q_2^{2}\cdot q_3^{2}\\
&\quad +60\cdot q_1^{2}\cdot q_2^{2}\cdot q_3^{4}+60\cdot q_1^{2}\cdot q_2^{4}\cdot q_3^{2}-60\cdot q_1^{4}\cdot q_2^{2}\cdot q_3^{2}\\
&\quad -360\cdot q_0^{2}\cdot q_1^{2}\cdot q_2^{2}\cdot q_3^{2}\end{align*}
\begin{align*}(T_{h}^4)_{6,7}&=-3\cdot q_0^{1}\cdot q_2^{7}+25\cdot q_0^{3}\cdot q_2^{5}-25\cdot q_0^{5}\cdot q_2^{3}+3\cdot q_0^{7}\cdot q_2^{1}\\
&\quad -3\cdot q_1^{1}\cdot q_3^{7}+25\cdot q_1^{3}\cdot q_3^{5}-25\cdot q_1^{5}\cdot q_3^{3}+3\cdot q_1^{7}\cdot q_3^{1}\\
&\quad -15\cdot q_0^{1}\cdot q_1^{2}\cdot q_2^{5}+35\cdot q_0^{1}\cdot q_1^{4}\cdot q_2^{3}-9\cdot q_0^{1}\cdot q_1^{6}\cdot q_2^{1}\\
&\quad +9\cdot q_0^{1}\cdot q_2^{1}\cdot q_3^{6}+15\cdot q_0^{1}\cdot q_2^{3}\cdot q_3^{4}+3\cdot q_0^{1}\cdot q_2^{5}\cdot q_3^{2}\\
&\quad -15\cdot q_0^{2}\cdot q_1^{1}\cdot q_3^{5}+10\cdot q_0^{2}\cdot q_1^{3}\cdot q_3^{3}-3\cdot q_0^{2}\cdot q_1^{5}\cdot q_3^{1}\\
&\quad +10\cdot q_0^{3}\cdot q_1^{2}\cdot q_2^{3}-15\cdot q_0^{3}\cdot q_1^{4}\cdot q_2^{1}-35\cdot q_0^{3}\cdot q_2^{1}\cdot q_3^{4}\\
&\quad -10\cdot q_0^{3}\cdot q_2^{3}\cdot q_3^{2}+35\cdot q_0^{4}\cdot q_1^{1}\cdot q_3^{3}-15\cdot q_0^{4}\cdot q_1^{3}\cdot q_3^{1}\\
&\quad -3\cdot q_0^{5}\cdot q_1^{2}\cdot q_2^{1}+15\cdot q_0^{5}\cdot q_2^{1}\cdot q_3^{2}-9\cdot q_0^{6}\cdot q_1^{1}\cdot q_3^{1}\\
&\quad +3\cdot q_1^{1}\cdot q_2^{2}\cdot q_3^{5}+15\cdot q_1^{1}\cdot q_2^{4}\cdot q_3^{3}+9\cdot q_1^{1}\cdot q_2^{6}\cdot q_3^{1}\\
&\quad -10\cdot q_1^{3}\cdot q_2^{2}\cdot q_3^{3}-35\cdot q_1^{3}\cdot q_2^{4}\cdot q_3^{1}+15\cdot q_1^{5}\cdot q_2^{2}\cdot q_3^{1}\\
&\quad -75\cdot q_0^{1}\cdot q_1^{2}\cdot q_2^{1}\cdot q_3^{4}-90\cdot q_0^{1}\cdot q_1^{2}\cdot q_2^{3}\cdot q_3^{2}\\
&\quad +75\cdot q_0^{1}\cdot q_1^{4}\cdot q_2^{1}\cdot q_3^{2}-90\cdot q_0^{2}\cdot q_1^{1}\cdot q_2^{2}\cdot q_3^{3}\\
&\quad -75\cdot q_0^{2}\cdot q_1^{1}\cdot q_2^{4}\cdot q_3^{1}+90\cdot q_0^{2}\cdot q_1^{3}\cdot q_2^{2}\cdot q_3^{1}\\
&\quad +90\cdot q_0^{3}\cdot q_1^{2}\cdot q_2^{1}\cdot q_3^{2}+75\cdot q_0^{4}\cdot q_1^{1}\cdot q_2^{2}\cdot q_3^{1}\end{align*}
\begin{align*}(T_{h}^4)_{6,8}&=-3\cdot q_0^{1}\cdot q_3^{7}+25\cdot q_0^{3}\cdot q_3^{5}-25\cdot q_0^{5}\cdot q_3^{3}+3\cdot q_0^{7}\cdot q_3^{1}\\
&\quad +3\cdot q_1^{1}\cdot q_2^{7}-25\cdot q_1^{3}\cdot q_2^{5}+25\cdot q_1^{5}\cdot q_2^{3}-3\cdot q_1^{7}\cdot q_2^{1}\\
&\quad -15\cdot q_0^{1}\cdot q_1^{2}\cdot q_3^{5}+35\cdot q_0^{1}\cdot q_1^{4}\cdot q_3^{3}-9\cdot q_0^{1}\cdot q_1^{6}\cdot q_3^{1}\\
&\quad +3\cdot q_0^{1}\cdot q_2^{2}\cdot q_3^{5}+15\cdot q_0^{1}\cdot q_2^{4}\cdot q_3^{3}+9\cdot q_0^{1}\cdot q_2^{6}\cdot q_3^{1}\\
&\quad +15\cdot q_0^{2}\cdot q_1^{1}\cdot q_2^{5}-10\cdot q_0^{2}\cdot q_1^{3}\cdot q_2^{3}+3\cdot q_0^{2}\cdot q_1^{5}\cdot q_2^{1}\\
&\quad +10\cdot q_0^{3}\cdot q_1^{2}\cdot q_3^{3}-15\cdot q_0^{3}\cdot q_1^{4}\cdot q_3^{1}-10\cdot q_0^{3}\cdot q_2^{2}\cdot q_3^{3}\\
&\quad -35\cdot q_0^{3}\cdot q_2^{4}\cdot q_3^{1}-35\cdot q_0^{4}\cdot q_1^{1}\cdot q_2^{3}+15\cdot q_0^{4}\cdot q_1^{3}\cdot q_2^{1}\\
&\quad -3\cdot q_0^{5}\cdot q_1^{2}\cdot q_3^{1}+15\cdot q_0^{5}\cdot q_2^{2}\cdot q_3^{1}+9\cdot q_0^{6}\cdot q_1^{1}\cdot q_2^{1}\\
&\quad -9\cdot q_1^{1}\cdot q_2^{1}\cdot q_3^{6}-15\cdot q_1^{1}\cdot q_2^{3}\cdot q_3^{4}-3\cdot q_1^{1}\cdot q_2^{5}\cdot q_3^{2}\\
&\quad +35\cdot q_1^{3}\cdot q_2^{1}\cdot q_3^{4}+10\cdot q_1^{3}\cdot q_2^{3}\cdot q_3^{2}-15\cdot q_1^{5}\cdot q_2^{1}\cdot q_3^{2}\\
&\quad -90\cdot q_0^{1}\cdot q_1^{2}\cdot q_2^{2}\cdot q_3^{3}-75\cdot q_0^{1}\cdot q_1^{2}\cdot q_2^{4}\cdot q_3^{1}\\
&\quad +75\cdot q_0^{1}\cdot q_1^{4}\cdot q_2^{2}\cdot q_3^{1}+75\cdot q_0^{2}\cdot q_1^{1}\cdot q_2^{1}\cdot q_3^{4}\\
&\quad +90\cdot q_0^{2}\cdot q_1^{1}\cdot q_2^{3}\cdot q_3^{2}-90\cdot q_0^{2}\cdot q_1^{3}\cdot q_2^{1}\cdot q_3^{2}\\
&\quad +90\cdot q_0^{3}\cdot q_1^{2}\cdot q_2^{2}\cdot q_3^{1}-75\cdot q_0^{4}\cdot q_1^{1}\cdot q_2^{1}\cdot q_3^{2}\end{align*}
\begin{align*}(T_{h}^4)_{6,9}&=-60\cdot q_0^{2}\cdot q_2^{6}+60\cdot q_0^{2}\cdot q_3^{6}+160\cdot q_0^{4}\cdot q_2^{4}-160\cdot q_0^{4}\cdot q_3^{4}\\
&\quad -60\cdot q_0^{6}\cdot q_2^{2}+60\cdot q_0^{6}\cdot q_3^{2}+60\cdot q_1^{2}\cdot q_2^{6}-60\cdot q_1^{2}\cdot q_3^{6}\\
&\quad -160\cdot q_1^{4}\cdot q_2^{4}+160\cdot q_1^{4}\cdot q_3^{4}+60\cdot q_1^{6}\cdot q_2^{2}-60\cdot q_1^{6}\cdot q_3^{2}\\
&\quad +60\cdot q_0^{2}\cdot q_1^{4}\cdot q_2^{2}-60\cdot q_0^{2}\cdot q_1^{4}\cdot q_3^{2}+60\cdot q_0^{2}\cdot q_2^{2}\cdot q_3^{4}\\
&\quad -60\cdot q_0^{2}\cdot q_2^{4}\cdot q_3^{2}-60\cdot q_0^{4}\cdot q_1^{2}\cdot q_2^{2}+60\cdot q_0^{4}\cdot q_1^{2}\cdot q_3^{2}\\
&\quad -60\cdot q_1^{2}\cdot q_2^{2}\cdot q_3^{4}+60\cdot q_1^{2}\cdot q_2^{4}\cdot q_3^{2}+240\cdot q_0^{1}\cdot q_1^{1}\cdot q_2^{1}\cdot q_3^{5}\\
&\quad +480\cdot q_0^{1}\cdot q_1^{1}\cdot q_2^{3}\cdot q_3^{3}+240\cdot q_0^{1}\cdot q_1^{1}\cdot q_2^{5}\cdot q_3^{1}\\
&\quad -640\cdot q_0^{1}\cdot q_1^{3}\cdot q_2^{1}\cdot q_3^{3}-640\cdot q_0^{1}\cdot q_1^{3}\cdot q_2^{3}\cdot q_3^{1}\\
&\quad +240\cdot q_0^{1}\cdot q_1^{5}\cdot q_2^{1}\cdot q_3^{1}-640\cdot q_0^{3}\cdot q_1^{1}\cdot q_2^{1}\cdot q_3^{3}\\
&\quad -640\cdot q_0^{3}\cdot q_1^{1}\cdot q_2^{3}\cdot q_3^{1}+480\cdot q_0^{3}\cdot q_1^{3}\cdot q_2^{1}\cdot q_3^{1}\\
&\quad +240\cdot q_0^{5}\cdot q_1^{1}\cdot q_2^{1}\cdot q_3^{1}\end{align*}
\begin{align*}(T_{h}^4)_{7,1}&=-2\cdot q_0^{3}\cdot q_3^{5}-2\cdot q_0^{5}\cdot q_3^{3}-2\cdot q_1^{3}\cdot q_2^{5}-2\cdot q_1^{5}\cdot q_2^{3}\\
&\quad +6\cdot q_0^{1}\cdot q_1^{2}\cdot q_3^{5}-10\cdot q_0^{1}\cdot q_1^{4}\cdot q_3^{3}+6\cdot q_0^{2}\cdot q_1^{1}\cdot q_2^{5}\\
&\quad +20\cdot q_0^{2}\cdot q_1^{3}\cdot q_2^{3}+20\cdot q_0^{3}\cdot q_1^{2}\cdot q_3^{3}+20\cdot q_0^{3}\cdot q_2^{2}\cdot q_3^{3}\\
&\quad -10\cdot q_0^{3}\cdot q_2^{4}\cdot q_3^{1}-10\cdot q_0^{4}\cdot q_1^{1}\cdot q_2^{3}+6\cdot q_0^{5}\cdot q_2^{2}\cdot q_3^{1}\\
&\quad -10\cdot q_1^{3}\cdot q_2^{1}\cdot q_3^{4}+20\cdot q_1^{3}\cdot q_2^{3}\cdot q_3^{2}+6\cdot q_1^{5}\cdot q_2^{1}\cdot q_3^{2}\\
&\quad -60\cdot q_0^{1}\cdot q_1^{2}\cdot q_2^{2}\cdot q_3^{3}+30\cdot q_0^{1}\cdot q_1^{2}\cdot q_2^{4}\cdot q_3^{1}\\
&\quad +30\cdot q_0^{1}\cdot q_1^{4}\cdot q_2^{2}\cdot q_3^{1}+30\cdot q_0^{2}\cdot q_1^{1}\cdot q_2^{1}\cdot q_3^{4}\\
&\quad -60\cdot q_0^{2}\cdot q_1^{1}\cdot q_2^{3}\cdot q_3^{2}-60\cdot q_0^{2}\cdot q_1^{3}\cdot q_2^{1}\cdot q_3^{2}\\
&\quad -60\cdot q_0^{3}\cdot q_1^{2}\cdot q_2^{2}\cdot q_3^{1}+30\cdot q_0^{4}\cdot q_1^{1}\cdot q_2^{1}\cdot q_3^{2}\end{align*}
\begin{align*}(T_{h}^4)_{7,2}&=-8\cdot q_0^{3}\cdot q_2^{5}+8\cdot q_0^{5}\cdot q_2^{3}+8\cdot q_1^{3}\cdot q_3^{5}-8\cdot q_1^{5}\cdot q_3^{3}\\
&\quad +24\cdot q_0^{1}\cdot q_1^{2}\cdot q_2^{5}+40\cdot q_0^{1}\cdot q_1^{4}\cdot q_2^{3}-24\cdot q_0^{2}\cdot q_1^{1}\cdot q_3^{5}\\
&\quad +80\cdot q_0^{2}\cdot q_1^{3}\cdot q_3^{3}-80\cdot q_0^{3}\cdot q_1^{2}\cdot q_2^{3}-40\cdot q_0^{3}\cdot q_2^{1}\cdot q_3^{4}\\
&\quad +80\cdot q_0^{3}\cdot q_2^{3}\cdot q_3^{2}-40\cdot q_0^{4}\cdot q_1^{1}\cdot q_3^{3}-24\cdot q_0^{5}\cdot q_2^{1}\cdot q_3^{2}\\
&\quad -80\cdot q_1^{3}\cdot q_2^{2}\cdot q_3^{3}+40\cdot q_1^{3}\cdot q_2^{4}\cdot q_3^{1}+24\cdot q_1^{5}\cdot q_2^{2}\cdot q_3^{1}\\
&\quad +120\cdot q_0^{1}\cdot q_1^{2}\cdot q_2^{1}\cdot q_3^{4}-240\cdot q_0^{1}\cdot q_1^{2}\cdot q_2^{3}\cdot q_3^{2}\\
&\quad -120\cdot q_0^{1}\cdot q_1^{4}\cdot q_2^{1}\cdot q_3^{2}+240\cdot q_0^{2}\cdot q_1^{1}\cdot q_2^{2}\cdot q_3^{3}\\
&\quad -120\cdot q_0^{2}\cdot q_1^{1}\cdot q_2^{4}\cdot q_3^{1}-240\cdot q_0^{2}\cdot q_1^{3}\cdot q_2^{2}\cdot q_3^{1}\\
&\quad +240\cdot q_0^{3}\cdot q_1^{2}\cdot q_2^{1}\cdot q_3^{2}+120\cdot q_0^{4}\cdot q_1^{1}\cdot q_2^{2}\cdot q_3^{1}\end{align*}
\begin{align*}(T_{h}^4)_{7,3}&=-3\cdot q_0^{2}\cdot q_2^{6}-3\cdot q_0^{2}\cdot q_3^{6}+10\cdot q_0^{4}\cdot q_2^{4}-3\cdot q_0^{6}\cdot q_2^{2}\\
&\quad +3\cdot q_0^{6}\cdot q_3^{2}+3\cdot q_1^{2}\cdot q_2^{6}+3\cdot q_1^{2}\cdot q_3^{6}-10\cdot q_1^{4}\cdot q_3^{4}\\
&\quad -3\cdot q_1^{6}\cdot q_2^{2}+3\cdot q_1^{6}\cdot q_3^{2}-30\cdot q_0^{2}\cdot q_1^{2}\cdot q_2^{4}+30\cdot q_0^{2}\cdot q_1^{2}\cdot q_3^{4}\\
&\quad +15\cdot q_0^{2}\cdot q_1^{4}\cdot q_2^{2}-15\cdot q_0^{2}\cdot q_1^{4}\cdot q_3^{2}+15\cdot q_0^{2}\cdot q_2^{2}\cdot q_3^{4}\\
&\quad +15\cdot q_0^{2}\cdot q_2^{4}\cdot q_3^{2}+15\cdot q_0^{4}\cdot q_1^{2}\cdot q_2^{2}-15\cdot q_0^{4}\cdot q_1^{2}\cdot q_3^{2}\\
&\quad -30\cdot q_0^{4}\cdot q_2^{2}\cdot q_3^{2}-15\cdot q_1^{2}\cdot q_2^{2}\cdot q_3^{4}-15\cdot q_1^{2}\cdot q_2^{4}\cdot q_3^{2}\\
&\quad +30\cdot q_1^{4}\cdot q_2^{2}\cdot q_3^{2}+24\cdot q_0^{1}\cdot q_1^{1}\cdot q_2^{1}\cdot q_3^{5}-24\cdot q_0^{1}\cdot q_1^{1}\cdot q_2^{5}\cdot q_3^{1}\\
&\quad -80\cdot q_0^{1}\cdot q_1^{3}\cdot q_2^{1}\cdot q_3^{3}+24\cdot q_0^{1}\cdot q_1^{5}\cdot q_2^{1}\cdot q_3^{1}\\
&\quad +80\cdot q_0^{3}\cdot q_1^{1}\cdot q_2^{3}\cdot q_3^{1}-24\cdot q_0^{5}\cdot q_1^{1}\cdot q_2^{1}\cdot q_3^{1}\end{align*}
\begin{align*}(T_{h}^4)_{7,4}&=-6\cdot q_0^{1}\cdot q_1^{1}\cdot q_2^{6}-6\cdot q_0^{1}\cdot q_1^{1}\cdot q_3^{6}-10\cdot q_0^{1}\cdot q_1^{3}\cdot q_2^{4}\\
&\quad +30\cdot q_0^{1}\cdot q_1^{3}\cdot q_3^{4}+12\cdot q_0^{1}\cdot q_1^{5}\cdot q_2^{2}-12\cdot q_0^{1}\cdot q_1^{5}\cdot q_3^{2}\\
&\quad -12\cdot q_0^{2}\cdot q_2^{1}\cdot q_3^{5}+12\cdot q_0^{2}\cdot q_2^{5}\cdot q_3^{1}+30\cdot q_0^{3}\cdot q_1^{1}\cdot q_2^{4}\\
&\quad -10\cdot q_0^{3}\cdot q_1^{1}\cdot q_3^{4}+10\cdot q_0^{4}\cdot q_2^{1}\cdot q_3^{3}-30\cdot q_0^{4}\cdot q_2^{3}\cdot q_3^{1}\\
&\quad -12\cdot q_0^{5}\cdot q_1^{1}\cdot q_2^{2}+12\cdot q_0^{5}\cdot q_1^{1}\cdot q_3^{2}+6\cdot q_0^{6}\cdot q_2^{1}\cdot q_3^{1}\\
&\quad +12\cdot q_1^{2}\cdot q_2^{1}\cdot q_3^{5}-12\cdot q_1^{2}\cdot q_2^{5}\cdot q_3^{1}-30\cdot q_1^{4}\cdot q_2^{1}\cdot q_3^{3}\\
&\quad +10\cdot q_1^{4}\cdot q_2^{3}\cdot q_3^{1}+6\cdot q_1^{6}\cdot q_2^{1}\cdot q_3^{1}+30\cdot q_0^{1}\cdot q_1^{1}\cdot q_2^{2}\cdot q_3^{4}\\
&\quad +30\cdot q_0^{1}\cdot q_1^{1}\cdot q_2^{4}\cdot q_3^{2}-60\cdot q_0^{1}\cdot q_1^{3}\cdot q_2^{2}\cdot q_3^{2}\\
&\quad +60\cdot q_0^{2}\cdot q_1^{2}\cdot q_2^{1}\cdot q_3^{3}+60\cdot q_0^{2}\cdot q_1^{2}\cdot q_2^{3}\cdot q_3^{1}\\
&\quad -30\cdot q_0^{2}\cdot q_1^{4}\cdot q_2^{1}\cdot q_3^{1}-60\cdot q_0^{3}\cdot q_1^{1}\cdot q_2^{2}\cdot q_3^{2}\\
&\quad -30\cdot q_0^{4}\cdot q_1^{2}\cdot q_2^{1}\cdot q_3^{1}\end{align*}
\begin{align*}(T_{h}^4)_{7,5}&=-3\cdot q_0^{1}\cdot q_3^{7}+5\cdot q_0^{3}\cdot q_3^{5}+5\cdot q_0^{5}\cdot q_3^{3}-3\cdot q_0^{7}\cdot q_3^{1}\\
&\quad -3\cdot q_1^{1}\cdot q_2^{7}+5\cdot q_1^{3}\cdot q_2^{5}+5\cdot q_1^{5}\cdot q_2^{3}-3\cdot q_1^{7}\cdot q_2^{1}\\
&\quad +45\cdot q_0^{1}\cdot q_1^{2}\cdot q_3^{5}-55\cdot q_0^{1}\cdot q_1^{4}\cdot q_3^{3}+9\cdot q_0^{1}\cdot q_1^{6}\cdot q_3^{1}\\
&\quad +3\cdot q_0^{1}\cdot q_2^{2}\cdot q_3^{5}+15\cdot q_0^{1}\cdot q_2^{4}\cdot q_3^{3}+9\cdot q_0^{1}\cdot q_2^{6}\cdot q_3^{1}\\
&\quad +45\cdot q_0^{2}\cdot q_1^{1}\cdot q_2^{5}-50\cdot q_0^{2}\cdot q_1^{3}\cdot q_2^{3}+3\cdot q_0^{2}\cdot q_1^{5}\cdot q_2^{1}\\
&\quad -50\cdot q_0^{3}\cdot q_1^{2}\cdot q_3^{3}+15\cdot q_0^{3}\cdot q_1^{4}\cdot q_3^{1}-50\cdot q_0^{3}\cdot q_2^{2}\cdot q_3^{3}\\
&\quad -55\cdot q_0^{3}\cdot q_2^{4}\cdot q_3^{1}-55\cdot q_0^{4}\cdot q_1^{1}\cdot q_2^{3}+15\cdot q_0^{4}\cdot q_1^{3}\cdot q_2^{1}\\
&\quad +3\cdot q_0^{5}\cdot q_1^{2}\cdot q_3^{1}+45\cdot q_0^{5}\cdot q_2^{2}\cdot q_3^{1}+9\cdot q_0^{6}\cdot q_1^{1}\cdot q_2^{1}\\
&\quad +9\cdot q_1^{1}\cdot q_2^{1}\cdot q_3^{6}+15\cdot q_1^{1}\cdot q_2^{3}\cdot q_3^{4}+3\cdot q_1^{1}\cdot q_2^{5}\cdot q_3^{2}\\
&\quad -55\cdot q_1^{3}\cdot q_2^{1}\cdot q_3^{4}-50\cdot q_1^{3}\cdot q_2^{3}\cdot q_3^{2}+45\cdot q_1^{5}\cdot q_2^{1}\cdot q_3^{2}\\
&\quad +30\cdot q_0^{1}\cdot q_1^{2}\cdot q_2^{2}\cdot q_3^{3}-15\cdot q_0^{1}\cdot q_1^{2}\cdot q_2^{4}\cdot q_3^{1}\\
&\quad -15\cdot q_0^{1}\cdot q_1^{4}\cdot q_2^{2}\cdot q_3^{1}-15\cdot q_0^{2}\cdot q_1^{1}\cdot q_2^{1}\cdot q_3^{4}\\
&\quad +30\cdot q_0^{2}\cdot q_1^{1}\cdot q_2^{3}\cdot q_3^{2}+30\cdot q_0^{2}\cdot q_1^{3}\cdot q_2^{1}\cdot q_3^{2}\\
&\quad +30\cdot q_0^{3}\cdot q_1^{2}\cdot q_2^{2}\cdot q_3^{1}-15\cdot q_0^{4}\cdot q_1^{1}\cdot q_2^{1}\cdot q_3^{2}\end{align*}
\begin{align*}(T_{h}^4)_{7,6}&=6\cdot q_0^{1}\cdot q_2^{7}-50\cdot q_0^{3}\cdot q_2^{5}+50\cdot q_0^{5}\cdot q_2^{3}-6\cdot q_0^{7}\cdot q_2^{1}\\
&\quad -6\cdot q_1^{1}\cdot q_3^{7}+50\cdot q_1^{3}\cdot q_3^{5}-50\cdot q_1^{5}\cdot q_3^{3}+6\cdot q_1^{7}\cdot q_3^{1}\\
&\quad +30\cdot q_0^{1}\cdot q_1^{2}\cdot q_2^{5}-70\cdot q_0^{1}\cdot q_1^{4}\cdot q_2^{3}+18\cdot q_0^{1}\cdot q_1^{6}\cdot q_2^{1}\\
&\quad -18\cdot q_0^{1}\cdot q_2^{1}\cdot q_3^{6}-30\cdot q_0^{1}\cdot q_2^{3}\cdot q_3^{4}-6\cdot q_0^{1}\cdot q_2^{5}\cdot q_3^{2}\\
&\quad -30\cdot q_0^{2}\cdot q_1^{1}\cdot q_3^{5}+20\cdot q_0^{2}\cdot q_1^{3}\cdot q_3^{3}-6\cdot q_0^{2}\cdot q_1^{5}\cdot q_3^{1}\\
&\quad -20\cdot q_0^{3}\cdot q_1^{2}\cdot q_2^{3}+30\cdot q_0^{3}\cdot q_1^{4}\cdot q_2^{1}+70\cdot q_0^{3}\cdot q_2^{1}\cdot q_3^{4}\\
&\quad +20\cdot q_0^{3}\cdot q_2^{3}\cdot q_3^{2}+70\cdot q_0^{4}\cdot q_1^{1}\cdot q_3^{3}-30\cdot q_0^{4}\cdot q_1^{3}\cdot q_3^{1}\\
&\quad +6\cdot q_0^{5}\cdot q_1^{2}\cdot q_2^{1}-30\cdot q_0^{5}\cdot q_2^{1}\cdot q_3^{2}-18\cdot q_0^{6}\cdot q_1^{1}\cdot q_3^{1}\\
&\quad +6\cdot q_1^{1}\cdot q_2^{2}\cdot q_3^{5}+30\cdot q_1^{1}\cdot q_2^{4}\cdot q_3^{3}+18\cdot q_1^{1}\cdot q_2^{6}\cdot q_3^{1}\\
&\quad -20\cdot q_1^{3}\cdot q_2^{2}\cdot q_3^{3}-70\cdot q_1^{3}\cdot q_2^{4}\cdot q_3^{1}+30\cdot q_1^{5}\cdot q_2^{2}\cdot q_3^{1}\\
&\quad +150\cdot q_0^{1}\cdot q_1^{2}\cdot q_2^{1}\cdot q_3^{4}+180\cdot q_0^{1}\cdot q_1^{2}\cdot q_2^{3}\cdot q_3^{2}\\
&\quad -150\cdot q_0^{1}\cdot q_1^{4}\cdot q_2^{1}\cdot q_3^{2}-180\cdot q_0^{2}\cdot q_1^{1}\cdot q_2^{2}\cdot q_3^{3}\\
&\quad -150\cdot q_0^{2}\cdot q_1^{1}\cdot q_2^{4}\cdot q_3^{1}+180\cdot q_0^{2}\cdot q_1^{3}\cdot q_2^{2}\cdot q_3^{1}\\
&\quad -180\cdot q_0^{3}\cdot q_1^{2}\cdot q_2^{1}\cdot q_3^{2}+150\cdot q_0^{4}\cdot q_1^{1}\cdot q_2^{2}\cdot q_3^{1}\end{align*}
\begin{align*}(T_{h}^4)_{7,7}&=1\cdot q_0^{8}-1\cdot q_1^{8}+1\cdot q_2^{8}-1\cdot q_3^{8}-2\cdot q_0^{2}\cdot q_1^{6}-25\cdot q_0^{2}\cdot q_2^{6}\\
&\quad +5\cdot q_0^{2}\cdot q_3^{6}+60\cdot q_0^{4}\cdot q_2^{4}+2\cdot q_0^{6}\cdot q_1^{2}-25\cdot q_0^{6}\cdot q_2^{2}\\
&\quad -5\cdot q_0^{6}\cdot q_3^{2}-5\cdot q_1^{2}\cdot q_2^{6}+25\cdot q_1^{2}\cdot q_3^{6}-60\cdot q_1^{4}\cdot q_3^{4}\\
&\quad +5\cdot q_1^{6}\cdot q_2^{2}+25\cdot q_1^{6}\cdot q_3^{2}-2\cdot q_2^{2}\cdot q_3^{6}+2\cdot q_2^{6}\cdot q_3^{2}\\
&\quad +60\cdot q_0^{2}\cdot q_1^{2}\cdot q_2^{4}-60\cdot q_0^{2}\cdot q_1^{2}\cdot q_3^{4}-15\cdot q_0^{2}\cdot q_1^{4}\cdot q_2^{2}\\
&\quad +45\cdot q_0^{2}\cdot q_1^{4}\cdot q_3^{2}-15\cdot q_0^{2}\cdot q_2^{2}\cdot q_3^{4}-45\cdot q_0^{2}\cdot q_2^{4}\cdot q_3^{2}\\
&\quad -45\cdot q_0^{4}\cdot q_1^{2}\cdot q_2^{2}+15\cdot q_0^{4}\cdot q_1^{2}\cdot q_3^{2}+60\cdot q_0^{4}\cdot q_2^{2}\cdot q_3^{2}\\
&\quad +45\cdot q_1^{2}\cdot q_2^{2}\cdot q_3^{4}+15\cdot q_1^{2}\cdot q_2^{4}\cdot q_3^{2}-60\cdot q_1^{4}\cdot q_2^{2}\cdot q_3^{2}\end{align*}
\begin{align*}(T_{h}^4)_{7,8}&=2\cdot q_0^{1}\cdot q_1^{7}+6\cdot q_0^{3}\cdot q_1^{5}+6\cdot q_0^{5}\cdot q_1^{3}+2\cdot q_0^{7}\cdot q_1^{1}\\
&\quad -2\cdot q_2^{1}\cdot q_3^{7}-6\cdot q_2^{3}\cdot q_3^{5}-6\cdot q_2^{5}\cdot q_3^{3}-2\cdot q_2^{7}\cdot q_3^{1}\\
&\quad -20\cdot q_0^{1}\cdot q_1^{1}\cdot q_2^{6}-20\cdot q_0^{1}\cdot q_1^{1}\cdot q_3^{6}+60\cdot q_0^{1}\cdot q_1^{3}\cdot q_2^{4}\\
&\quad +60\cdot q_0^{1}\cdot q_1^{3}\cdot q_3^{4}-30\cdot q_0^{1}\cdot q_1^{5}\cdot q_2^{2}-30\cdot q_0^{1}\cdot q_1^{5}\cdot q_3^{2}\\
&\quad +30\cdot q_0^{2}\cdot q_2^{1}\cdot q_3^{5}+60\cdot q_0^{2}\cdot q_2^{3}\cdot q_3^{3}+30\cdot q_0^{2}\cdot q_2^{5}\cdot q_3^{1}\\
&\quad +60\cdot q_0^{3}\cdot q_1^{1}\cdot q_2^{4}+60\cdot q_0^{3}\cdot q_1^{1}\cdot q_3^{4}-60\cdot q_0^{3}\cdot q_1^{3}\cdot q_2^{2}\\
&\quad -60\cdot q_0^{3}\cdot q_1^{3}\cdot q_3^{2}-60\cdot q_0^{4}\cdot q_2^{1}\cdot q_3^{3}-60\cdot q_0^{4}\cdot q_2^{3}\cdot q_3^{1}\\
&\quad -30\cdot q_0^{5}\cdot q_1^{1}\cdot q_2^{2}-30\cdot q_0^{5}\cdot q_1^{1}\cdot q_3^{2}+20\cdot q_0^{6}\cdot q_2^{1}\cdot q_3^{1}\\
&\quad +30\cdot q_1^{2}\cdot q_2^{1}\cdot q_3^{5}+60\cdot q_1^{2}\cdot q_2^{3}\cdot q_3^{3}+30\cdot q_1^{2}\cdot q_2^{5}\cdot q_3^{1}\\
&\quad -60\cdot q_1^{4}\cdot q_2^{1}\cdot q_3^{3}-60\cdot q_1^{4}\cdot q_2^{3}\cdot q_3^{1}+20\cdot q_1^{6}\cdot q_2^{1}\cdot q_3^{1}\\
&\quad -60\cdot q_0^{1}\cdot q_1^{1}\cdot q_2^{2}\cdot q_3^{4}-60\cdot q_0^{1}\cdot q_1^{1}\cdot q_2^{4}\cdot q_3^{2}\\
&\quad +120\cdot q_0^{1}\cdot q_1^{3}\cdot q_2^{2}\cdot q_3^{2}-120\cdot q_0^{2}\cdot q_1^{2}\cdot q_2^{1}\cdot q_3^{3}\\
&\quad -120\cdot q_0^{2}\cdot q_1^{2}\cdot q_2^{3}\cdot q_3^{1}+60\cdot q_0^{2}\cdot q_1^{4}\cdot q_2^{1}\cdot q_3^{1}\\
&\quad +120\cdot q_0^{3}\cdot q_1^{1}\cdot q_2^{2}\cdot q_3^{2}+60\cdot q_0^{4}\cdot q_1^{2}\cdot q_2^{1}\cdot q_3^{1}\end{align*}
\begin{align*}(T_{h}^4)_{7,9}&=40\cdot q_0^{1}\cdot q_2^{7}-240\cdot q_0^{3}\cdot q_2^{5}+240\cdot q_0^{5}\cdot q_2^{3}-40\cdot q_0^{7}\cdot q_2^{1}\\
&\quad -40\cdot q_1^{1}\cdot q_3^{7}+240\cdot q_1^{3}\cdot q_3^{5}-240\cdot q_1^{5}\cdot q_3^{3}+40\cdot q_1^{7}\cdot q_3^{1}\\
&\quad -240\cdot q_0^{1}\cdot q_1^{2}\cdot q_2^{5}+240\cdot q_0^{1}\cdot q_1^{4}\cdot q_2^{3}-40\cdot q_0^{1}\cdot q_1^{6}\cdot q_2^{1}\\
&\quad +40\cdot q_0^{1}\cdot q_2^{1}\cdot q_3^{6}+120\cdot q_0^{1}\cdot q_2^{3}\cdot q_3^{4}+120\cdot q_0^{1}\cdot q_2^{5}\cdot q_3^{2}\\
&\quad +240\cdot q_0^{2}\cdot q_1^{1}\cdot q_3^{5}-480\cdot q_0^{2}\cdot q_1^{3}\cdot q_3^{3}+120\cdot q_0^{2}\cdot q_1^{5}\cdot q_3^{1}\\
&\quad +480\cdot q_0^{3}\cdot q_1^{2}\cdot q_2^{3}-120\cdot q_0^{3}\cdot q_1^{4}\cdot q_2^{1}-240\cdot q_0^{3}\cdot q_2^{1}\cdot q_3^{4}\\
&\quad -480\cdot q_0^{3}\cdot q_2^{3}\cdot q_3^{2}-240\cdot q_0^{4}\cdot q_1^{1}\cdot q_3^{3}+120\cdot q_0^{4}\cdot q_1^{3}\cdot q_3^{1}\\
&\quad -120\cdot q_0^{5}\cdot q_1^{2}\cdot q_2^{1}+240\cdot q_0^{5}\cdot q_2^{1}\cdot q_3^{2}+40\cdot q_0^{6}\cdot q_1^{1}\cdot q_3^{1}\\
&\quad -120\cdot q_1^{1}\cdot q_2^{2}\cdot q_3^{5}-120\cdot q_1^{1}\cdot q_2^{4}\cdot q_3^{3}-40\cdot q_1^{1}\cdot q_2^{6}\cdot q_3^{1}\\
&\quad +480\cdot q_1^{3}\cdot q_2^{2}\cdot q_3^{3}+240\cdot q_1^{3}\cdot q_2^{4}\cdot q_3^{1}-240\cdot q_1^{5}\cdot q_2^{2}\cdot q_3^{1}\\
&\quad -240\cdot q_0^{1}\cdot q_1^{2}\cdot q_2^{1}\cdot q_3^{4}-480\cdot q_0^{1}\cdot q_1^{2}\cdot q_2^{3}\cdot q_3^{2}\\
&\quad +240\cdot q_0^{1}\cdot q_1^{4}\cdot q_2^{1}\cdot q_3^{2}+480\cdot q_0^{2}\cdot q_1^{1}\cdot q_2^{2}\cdot q_3^{3}\\
&\quad +240\cdot q_0^{2}\cdot q_1^{1}\cdot q_2^{4}\cdot q_3^{1}-480\cdot q_0^{2}\cdot q_1^{3}\cdot q_2^{2}\cdot q_3^{1}\\
&\quad +480\cdot q_0^{3}\cdot q_1^{2}\cdot q_2^{1}\cdot q_3^{2}-240\cdot q_0^{4}\cdot q_1^{1}\cdot q_2^{2}\cdot q_3^{1}\end{align*}
\begin{align*}(T_{h}^4)_{8,1}&=-2\cdot q_0^{3}\cdot q_2^{5}-2\cdot q_0^{5}\cdot q_2^{3}+2\cdot q_1^{3}\cdot q_3^{5}+2\cdot q_1^{5}\cdot q_3^{3}\\
&\quad +6\cdot q_0^{1}\cdot q_1^{2}\cdot q_2^{5}-10\cdot q_0^{1}\cdot q_1^{4}\cdot q_2^{3}-6\cdot q_0^{2}\cdot q_1^{1}\cdot q_3^{5}\\
&\quad -20\cdot q_0^{2}\cdot q_1^{3}\cdot q_3^{3}+20\cdot q_0^{3}\cdot q_1^{2}\cdot q_2^{3}-10\cdot q_0^{3}\cdot q_2^{1}\cdot q_3^{4}\\
&\quad +20\cdot q_0^{3}\cdot q_2^{3}\cdot q_3^{2}+10\cdot q_0^{4}\cdot q_1^{1}\cdot q_3^{3}+6\cdot q_0^{5}\cdot q_2^{1}\cdot q_3^{2}\\
&\quad -20\cdot q_1^{3}\cdot q_2^{2}\cdot q_3^{3}+10\cdot q_1^{3}\cdot q_2^{4}\cdot q_3^{1}-6\cdot q_1^{5}\cdot q_2^{2}\cdot q_3^{1}\\
&\quad +30\cdot q_0^{1}\cdot q_1^{2}\cdot q_2^{1}\cdot q_3^{4}-60\cdot q_0^{1}\cdot q_1^{2}\cdot q_2^{3}\cdot q_3^{2}\\
&\quad +30\cdot q_0^{1}\cdot q_1^{4}\cdot q_2^{1}\cdot q_3^{2}+60\cdot q_0^{2}\cdot q_1^{1}\cdot q_2^{2}\cdot q_3^{3}\\
&\quad -30\cdot q_0^{2}\cdot q_1^{1}\cdot q_2^{4}\cdot q_3^{1}+60\cdot q_0^{2}\cdot q_1^{3}\cdot q_2^{2}\cdot q_3^{1}\\
&\quad -60\cdot q_0^{3}\cdot q_1^{2}\cdot q_2^{1}\cdot q_3^{2}-30\cdot q_0^{4}\cdot q_1^{1}\cdot q_2^{2}\cdot q_3^{1}\end{align*}
\begin{align*}(T_{h}^4)_{8,2}&=8\cdot q_0^{3}\cdot q_3^{5}-8\cdot q_0^{5}\cdot q_3^{3}+8\cdot q_1^{3}\cdot q_2^{5}-8\cdot q_1^{5}\cdot q_2^{3}\\
&\quad -24\cdot q_0^{1}\cdot q_1^{2}\cdot q_3^{5}-40\cdot q_0^{1}\cdot q_1^{4}\cdot q_3^{3}-24\cdot q_0^{2}\cdot q_1^{1}\cdot q_2^{5}\\
&\quad +80\cdot q_0^{2}\cdot q_1^{3}\cdot q_2^{3}+80\cdot q_0^{3}\cdot q_1^{2}\cdot q_3^{3}-80\cdot q_0^{3}\cdot q_2^{2}\cdot q_3^{3}\\
&\quad +40\cdot q_0^{3}\cdot q_2^{4}\cdot q_3^{1}-40\cdot q_0^{4}\cdot q_1^{1}\cdot q_2^{3}+24\cdot q_0^{5}\cdot q_2^{2}\cdot q_3^{1}\\
&\quad +40\cdot q_1^{3}\cdot q_2^{1}\cdot q_3^{4}-80\cdot q_1^{3}\cdot q_2^{3}\cdot q_3^{2}+24\cdot q_1^{5}\cdot q_2^{1}\cdot q_3^{2}\\
&\quad +240\cdot q_0^{1}\cdot q_1^{2}\cdot q_2^{2}\cdot q_3^{3}-120\cdot q_0^{1}\cdot q_1^{2}\cdot q_2^{4}\cdot q_3^{1}\\
&\quad +120\cdot q_0^{1}\cdot q_1^{4}\cdot q_2^{2}\cdot q_3^{1}-120\cdot q_0^{2}\cdot q_1^{1}\cdot q_2^{1}\cdot q_3^{4}\\
&\quad +240\cdot q_0^{2}\cdot q_1^{1}\cdot q_2^{3}\cdot q_3^{2}-240\cdot q_0^{2}\cdot q_1^{3}\cdot q_2^{1}\cdot q_3^{2}\\
&\quad -240\cdot q_0^{3}\cdot q_1^{2}\cdot q_2^{2}\cdot q_3^{1}+120\cdot q_0^{4}\cdot q_1^{1}\cdot q_2^{1}\cdot q_3^{2}\end{align*}
\begin{align*}(T_{h}^4)_{8,3}&=-6\cdot q_0^{1}\cdot q_1^{1}\cdot q_2^{6}-6\cdot q_0^{1}\cdot q_1^{1}\cdot q_3^{6}+30\cdot q_0^{1}\cdot q_1^{3}\cdot q_2^{4}\\
&\quad -10\cdot q_0^{1}\cdot q_1^{3}\cdot q_3^{4}-12\cdot q_0^{1}\cdot q_1^{5}\cdot q_2^{2}+12\cdot q_0^{1}\cdot q_1^{5}\cdot q_3^{2}\\
&\quad -12\cdot q_0^{2}\cdot q_2^{1}\cdot q_3^{5}+12\cdot q_0^{2}\cdot q_2^{5}\cdot q_3^{1}-10\cdot q_0^{3}\cdot q_1^{1}\cdot q_2^{4}\\
&\quad +30\cdot q_0^{3}\cdot q_1^{1}\cdot q_3^{4}+30\cdot q_0^{4}\cdot q_2^{1}\cdot q_3^{3}-10\cdot q_0^{4}\cdot q_2^{3}\cdot q_3^{1}\\
&\quad +12\cdot q_0^{5}\cdot q_1^{1}\cdot q_2^{2}-12\cdot q_0^{5}\cdot q_1^{1}\cdot q_3^{2}-6\cdot q_0^{6}\cdot q_2^{1}\cdot q_3^{1}\\
&\quad +12\cdot q_1^{2}\cdot q_2^{1}\cdot q_3^{5}-12\cdot q_1^{2}\cdot q_2^{5}\cdot q_3^{1}-10\cdot q_1^{4}\cdot q_2^{1}\cdot q_3^{3}\\
&\quad +30\cdot q_1^{4}\cdot q_2^{3}\cdot q_3^{1}-6\cdot q_1^{6}\cdot q_2^{1}\cdot q_3^{1}+30\cdot q_0^{1}\cdot q_1^{1}\cdot q_2^{2}\cdot q_3^{4}\\
&\quad +30\cdot q_0^{1}\cdot q_1^{1}\cdot q_2^{4}\cdot q_3^{2}-60\cdot q_0^{1}\cdot q_1^{3}\cdot q_2^{2}\cdot q_3^{2}\\
&\quad -60\cdot q_0^{2}\cdot q_1^{2}\cdot q_2^{1}\cdot q_3^{3}-60\cdot q_0^{2}\cdot q_1^{2}\cdot q_2^{3}\cdot q_3^{1}\\
&\quad +30\cdot q_0^{2}\cdot q_1^{4}\cdot q_2^{1}\cdot q_3^{1}-60\cdot q_0^{3}\cdot q_1^{1}\cdot q_2^{2}\cdot q_3^{2}\\
&\quad +30\cdot q_0^{4}\cdot q_1^{2}\cdot q_2^{1}\cdot q_3^{1}\end{align*}
\begin{align*}(T_{h}^4)_{8,4}&=3\cdot q_0^{2}\cdot q_2^{6}+3\cdot q_0^{2}\cdot q_3^{6}-10\cdot q_0^{4}\cdot q_3^{4}-3\cdot q_0^{6}\cdot q_2^{2}\\
&\quad +3\cdot q_0^{6}\cdot q_3^{2}-3\cdot q_1^{2}\cdot q_2^{6}-3\cdot q_1^{2}\cdot q_3^{6}+10\cdot q_1^{4}\cdot q_2^{4}\\
&\quad -3\cdot q_1^{6}\cdot q_2^{2}+3\cdot q_1^{6}\cdot q_3^{2}-30\cdot q_0^{2}\cdot q_1^{2}\cdot q_2^{4}+30\cdot q_0^{2}\cdot q_1^{2}\cdot q_3^{4}\\
&\quad +15\cdot q_0^{2}\cdot q_1^{4}\cdot q_2^{2}-15\cdot q_0^{2}\cdot q_1^{4}\cdot q_3^{2}-15\cdot q_0^{2}\cdot q_2^{2}\cdot q_3^{4}\\
&\quad -15\cdot q_0^{2}\cdot q_2^{4}\cdot q_3^{2}+15\cdot q_0^{4}\cdot q_1^{2}\cdot q_2^{2}-15\cdot q_0^{4}\cdot q_1^{2}\cdot q_3^{2}\\
&\quad +30\cdot q_0^{4}\cdot q_2^{2}\cdot q_3^{2}+15\cdot q_1^{2}\cdot q_2^{2}\cdot q_3^{4}+15\cdot q_1^{2}\cdot q_2^{4}\cdot q_3^{2}\\
&\quad -30\cdot q_1^{4}\cdot q_2^{2}\cdot q_3^{2}-24\cdot q_0^{1}\cdot q_1^{1}\cdot q_2^{1}\cdot q_3^{5}+24\cdot q_0^{1}\cdot q_1^{1}\cdot q_2^{5}\cdot q_3^{1}\\
&\quad -80\cdot q_0^{1}\cdot q_1^{3}\cdot q_2^{3}\cdot q_3^{1}+24\cdot q_0^{1}\cdot q_1^{5}\cdot q_2^{1}\cdot q_3^{1}\\
&\quad +80\cdot q_0^{3}\cdot q_1^{1}\cdot q_2^{1}\cdot q_3^{3}-24\cdot q_0^{5}\cdot q_1^{1}\cdot q_2^{1}\cdot q_3^{1}\end{align*}
\begin{align*}(T_{h}^4)_{8,5}&=3\cdot q_0^{1}\cdot q_2^{7}-5\cdot q_0^{3}\cdot q_2^{5}-5\cdot q_0^{5}\cdot q_2^{3}+3\cdot q_0^{7}\cdot q_2^{1}\\
&\quad -3\cdot q_1^{1}\cdot q_3^{7}+5\cdot q_1^{3}\cdot q_3^{5}+5\cdot q_1^{5}\cdot q_3^{3}-3\cdot q_1^{7}\cdot q_3^{1}\\
&\quad -45\cdot q_0^{1}\cdot q_1^{2}\cdot q_2^{5}+55\cdot q_0^{1}\cdot q_1^{4}\cdot q_2^{3}-9\cdot q_0^{1}\cdot q_1^{6}\cdot q_2^{1}\\
&\quad -9\cdot q_0^{1}\cdot q_2^{1}\cdot q_3^{6}-15\cdot q_0^{1}\cdot q_2^{3}\cdot q_3^{4}-3\cdot q_0^{1}\cdot q_2^{5}\cdot q_3^{2}\\
&\quad +45\cdot q_0^{2}\cdot q_1^{1}\cdot q_3^{5}-50\cdot q_0^{2}\cdot q_1^{3}\cdot q_3^{3}+3\cdot q_0^{2}\cdot q_1^{5}\cdot q_3^{1}\\
&\quad +50\cdot q_0^{3}\cdot q_1^{2}\cdot q_2^{3}-15\cdot q_0^{3}\cdot q_1^{4}\cdot q_2^{1}+55\cdot q_0^{3}\cdot q_2^{1}\cdot q_3^{4}\\
&\quad +50\cdot q_0^{3}\cdot q_2^{3}\cdot q_3^{2}-55\cdot q_0^{4}\cdot q_1^{1}\cdot q_3^{3}+15\cdot q_0^{4}\cdot q_1^{3}\cdot q_3^{1}\\
&\quad -3\cdot q_0^{5}\cdot q_1^{2}\cdot q_2^{1}-45\cdot q_0^{5}\cdot q_2^{1}\cdot q_3^{2}+9\cdot q_0^{6}\cdot q_1^{1}\cdot q_3^{1}\\
&\quad +3\cdot q_1^{1}\cdot q_2^{2}\cdot q_3^{5}+15\cdot q_1^{1}\cdot q_2^{4}\cdot q_3^{3}+9\cdot q_1^{1}\cdot q_2^{6}\cdot q_3^{1}\\
&\quad -50\cdot q_1^{3}\cdot q_2^{2}\cdot q_3^{3}-55\cdot q_1^{3}\cdot q_2^{4}\cdot q_3^{1}+45\cdot q_1^{5}\cdot q_2^{2}\cdot q_3^{1}\\
&\quad +15\cdot q_0^{1}\cdot q_1^{2}\cdot q_2^{1}\cdot q_3^{4}-30\cdot q_0^{1}\cdot q_1^{2}\cdot q_2^{3}\cdot q_3^{2}\\
&\quad +15\cdot q_0^{1}\cdot q_1^{4}\cdot q_2^{1}\cdot q_3^{2}+30\cdot q_0^{2}\cdot q_1^{1}\cdot q_2^{2}\cdot q_3^{3}\\
&\quad -15\cdot q_0^{2}\cdot q_1^{1}\cdot q_2^{4}\cdot q_3^{1}+30\cdot q_0^{2}\cdot q_1^{3}\cdot q_2^{2}\cdot q_3^{1}\\
&\quad -30\cdot q_0^{3}\cdot q_1^{2}\cdot q_2^{1}\cdot q_3^{2}-15\cdot q_0^{4}\cdot q_1^{1}\cdot q_2^{2}\cdot q_3^{1}\end{align*}
\begin{align*}(T_{h}^4)_{8,6}&=6\cdot q_0^{1}\cdot q_3^{7}-50\cdot q_0^{3}\cdot q_3^{5}+50\cdot q_0^{5}\cdot q_3^{3}-6\cdot q_0^{7}\cdot q_3^{1}\\
&\quad +6\cdot q_1^{1}\cdot q_2^{7}-50\cdot q_1^{3}\cdot q_2^{5}+50\cdot q_1^{5}\cdot q_2^{3}-6\cdot q_1^{7}\cdot q_2^{1}\\
&\quad +30\cdot q_0^{1}\cdot q_1^{2}\cdot q_3^{5}-70\cdot q_0^{1}\cdot q_1^{4}\cdot q_3^{3}+18\cdot q_0^{1}\cdot q_1^{6}\cdot q_3^{1}\\
&\quad -6\cdot q_0^{1}\cdot q_2^{2}\cdot q_3^{5}-30\cdot q_0^{1}\cdot q_2^{4}\cdot q_3^{3}-18\cdot q_0^{1}\cdot q_2^{6}\cdot q_3^{1}\\
&\quad +30\cdot q_0^{2}\cdot q_1^{1}\cdot q_2^{5}-20\cdot q_0^{2}\cdot q_1^{3}\cdot q_2^{3}+6\cdot q_0^{2}\cdot q_1^{5}\cdot q_2^{1}\\
&\quad -20\cdot q_0^{3}\cdot q_1^{2}\cdot q_3^{3}+30\cdot q_0^{3}\cdot q_1^{4}\cdot q_3^{1}+20\cdot q_0^{3}\cdot q_2^{2}\cdot q_3^{3}\\
&\quad +70\cdot q_0^{3}\cdot q_2^{4}\cdot q_3^{1}-70\cdot q_0^{4}\cdot q_1^{1}\cdot q_2^{3}+30\cdot q_0^{4}\cdot q_1^{3}\cdot q_2^{1}\\
&\quad +6\cdot q_0^{5}\cdot q_1^{2}\cdot q_3^{1}-30\cdot q_0^{5}\cdot q_2^{2}\cdot q_3^{1}+18\cdot q_0^{6}\cdot q_1^{1}\cdot q_2^{1}\\
&\quad -18\cdot q_1^{1}\cdot q_2^{1}\cdot q_3^{6}-30\cdot q_1^{1}\cdot q_2^{3}\cdot q_3^{4}-6\cdot q_1^{1}\cdot q_2^{5}\cdot q_3^{2}\\
&\quad +70\cdot q_1^{3}\cdot q_2^{1}\cdot q_3^{4}+20\cdot q_1^{3}\cdot q_2^{3}\cdot q_3^{2}-30\cdot q_1^{5}\cdot q_2^{1}\cdot q_3^{2}\\
&\quad +180\cdot q_0^{1}\cdot q_1^{2}\cdot q_2^{2}\cdot q_3^{3}+150\cdot q_0^{1}\cdot q_1^{2}\cdot q_2^{4}\cdot q_3^{1}\\
&\quad -150\cdot q_0^{1}\cdot q_1^{4}\cdot q_2^{2}\cdot q_3^{1}+150\cdot q_0^{2}\cdot q_1^{1}\cdot q_2^{1}\cdot q_3^{4}\\
&\quad +180\cdot q_0^{2}\cdot q_1^{1}\cdot q_2^{3}\cdot q_3^{2}-180\cdot q_0^{2}\cdot q_1^{3}\cdot q_2^{1}\cdot q_3^{2}\\
&\quad -180\cdot q_0^{3}\cdot q_1^{2}\cdot q_2^{2}\cdot q_3^{1}-150\cdot q_0^{4}\cdot q_1^{1}\cdot q_2^{1}\cdot q_3^{2}\end{align*}
\begin{align*}(T_{h}^4)_{8,7}&=-2\cdot q_0^{1}\cdot q_1^{7}-6\cdot q_0^{3}\cdot q_1^{5}-6\cdot q_0^{5}\cdot q_1^{3}-2\cdot q_0^{7}\cdot q_1^{1}\\
&\quad -2\cdot q_2^{1}\cdot q_3^{7}-6\cdot q_2^{3}\cdot q_3^{5}-6\cdot q_2^{5}\cdot q_3^{3}-2\cdot q_2^{7}\cdot q_3^{1}\\
&\quad +20\cdot q_0^{1}\cdot q_1^{1}\cdot q_2^{6}+20\cdot q_0^{1}\cdot q_1^{1}\cdot q_3^{6}-60\cdot q_0^{1}\cdot q_1^{3}\cdot q_2^{4}\\
&\quad -60\cdot q_0^{1}\cdot q_1^{3}\cdot q_3^{4}+30\cdot q_0^{1}\cdot q_1^{5}\cdot q_2^{2}+30\cdot q_0^{1}\cdot q_1^{5}\cdot q_3^{2}\\
&\quad +30\cdot q_0^{2}\cdot q_2^{1}\cdot q_3^{5}+60\cdot q_0^{2}\cdot q_2^{3}\cdot q_3^{3}+30\cdot q_0^{2}\cdot q_2^{5}\cdot q_3^{1}\\
&\quad -60\cdot q_0^{3}\cdot q_1^{1}\cdot q_2^{4}-60\cdot q_0^{3}\cdot q_1^{1}\cdot q_3^{4}+60\cdot q_0^{3}\cdot q_1^{3}\cdot q_2^{2}\\
&\quad +60\cdot q_0^{3}\cdot q_1^{3}\cdot q_3^{2}-60\cdot q_0^{4}\cdot q_2^{1}\cdot q_3^{3}-60\cdot q_0^{4}\cdot q_2^{3}\cdot q_3^{1}\\
&\quad +30\cdot q_0^{5}\cdot q_1^{1}\cdot q_2^{2}+30\cdot q_0^{5}\cdot q_1^{1}\cdot q_3^{2}+20\cdot q_0^{6}\cdot q_2^{1}\cdot q_3^{1}\\
&\quad +30\cdot q_1^{2}\cdot q_2^{1}\cdot q_3^{5}+60\cdot q_1^{2}\cdot q_2^{3}\cdot q_3^{3}+30\cdot q_1^{2}\cdot q_2^{5}\cdot q_3^{1}\\
&\quad -60\cdot q_1^{4}\cdot q_2^{1}\cdot q_3^{3}-60\cdot q_1^{4}\cdot q_2^{3}\cdot q_3^{1}+20\cdot q_1^{6}\cdot q_2^{1}\cdot q_3^{1}\\
&\quad +60\cdot q_0^{1}\cdot q_1^{1}\cdot q_2^{2}\cdot q_3^{4}+60\cdot q_0^{1}\cdot q_1^{1}\cdot q_2^{4}\cdot q_3^{2}\\
&\quad -120\cdot q_0^{1}\cdot q_1^{3}\cdot q_2^{2}\cdot q_3^{2}-120\cdot q_0^{2}\cdot q_1^{2}\cdot q_2^{1}\cdot q_3^{3}\\
&\quad -120\cdot q_0^{2}\cdot q_1^{2}\cdot q_2^{3}\cdot q_3^{1}+60\cdot q_0^{2}\cdot q_1^{4}\cdot q_2^{1}\cdot q_3^{1}\\
&\quad -120\cdot q_0^{3}\cdot q_1^{1}\cdot q_2^{2}\cdot q_3^{2}+60\cdot q_0^{4}\cdot q_1^{2}\cdot q_2^{1}\cdot q_3^{1}\end{align*}
\begin{align*}(T_{h}^4)_{8,8}&=1\cdot q_0^{8}-1\cdot q_1^{8}-1\cdot q_2^{8}+1\cdot q_3^{8}-2\cdot q_0^{2}\cdot q_1^{6}+5\cdot q_0^{2}\cdot q_2^{6}\\
&\quad -25\cdot q_0^{2}\cdot q_3^{6}+60\cdot q_0^{4}\cdot q_3^{4}+2\cdot q_0^{6}\cdot q_1^{2}-5\cdot q_0^{6}\cdot q_2^{2}\\
&\quad -25\cdot q_0^{6}\cdot q_3^{2}+25\cdot q_1^{2}\cdot q_2^{6}-5\cdot q_1^{2}\cdot q_3^{6}-60\cdot q_1^{4}\cdot q_2^{4}\\
&\quad +25\cdot q_1^{6}\cdot q_2^{2}+5\cdot q_1^{6}\cdot q_3^{2}+2\cdot q_2^{2}\cdot q_3^{6}-2\cdot q_2^{6}\cdot q_3^{2}\\
&\quad -60\cdot q_0^{2}\cdot q_1^{2}\cdot q_2^{4}+60\cdot q_0^{2}\cdot q_1^{2}\cdot q_3^{4}+45\cdot q_0^{2}\cdot q_1^{4}\cdot q_2^{2}\\
&\quad -15\cdot q_0^{2}\cdot q_1^{4}\cdot q_3^{2}-45\cdot q_0^{2}\cdot q_2^{2}\cdot q_3^{4}-15\cdot q_0^{2}\cdot q_2^{4}\cdot q_3^{2}\\
&\quad +15\cdot q_0^{4}\cdot q_1^{2}\cdot q_2^{2}-45\cdot q_0^{4}\cdot q_1^{2}\cdot q_3^{2}+60\cdot q_0^{4}\cdot q_2^{2}\cdot q_3^{2}\\
&\quad +15\cdot q_1^{2}\cdot q_2^{2}\cdot q_3^{4}+45\cdot q_1^{2}\cdot q_2^{4}\cdot q_3^{2}-60\cdot q_1^{4}\cdot q_2^{2}\cdot q_3^{2}\end{align*}
\begin{align*}(T_{h}^4)_{8,9}&=-40\cdot q_0^{1}\cdot q_3^{7}+240\cdot q_0^{3}\cdot q_3^{5}-240\cdot q_0^{5}\cdot q_3^{3}+40\cdot q_0^{7}\cdot q_3^{1}\\
&\quad -40\cdot q_1^{1}\cdot q_2^{7}+240\cdot q_1^{3}\cdot q_2^{5}-240\cdot q_1^{5}\cdot q_2^{3}+40\cdot q_1^{7}\cdot q_2^{1}\\
&\quad +240\cdot q_0^{1}\cdot q_1^{2}\cdot q_3^{5}-240\cdot q_0^{1}\cdot q_1^{4}\cdot q_3^{3}+40\cdot q_0^{1}\cdot q_1^{6}\cdot q_3^{1}\\
&\quad -120\cdot q_0^{1}\cdot q_2^{2}\cdot q_3^{5}-120\cdot q_0^{1}\cdot q_2^{4}\cdot q_3^{3}-40\cdot q_0^{1}\cdot q_2^{6}\cdot q_3^{1}\\
&\quad +240\cdot q_0^{2}\cdot q_1^{1}\cdot q_2^{5}-480\cdot q_0^{2}\cdot q_1^{3}\cdot q_2^{3}+120\cdot q_0^{2}\cdot q_1^{5}\cdot q_2^{1}\\
&\quad -480\cdot q_0^{3}\cdot q_1^{2}\cdot q_3^{3}+120\cdot q_0^{3}\cdot q_1^{4}\cdot q_3^{1}+480\cdot q_0^{3}\cdot q_2^{2}\cdot q_3^{3}\\
&\quad +240\cdot q_0^{3}\cdot q_2^{4}\cdot q_3^{1}-240\cdot q_0^{4}\cdot q_1^{1}\cdot q_2^{3}+120\cdot q_0^{4}\cdot q_1^{3}\cdot q_2^{1}\\
&\quad +120\cdot q_0^{5}\cdot q_1^{2}\cdot q_3^{1}-240\cdot q_0^{5}\cdot q_2^{2}\cdot q_3^{1}+40\cdot q_0^{6}\cdot q_1^{1}\cdot q_2^{1}\\
&\quad -40\cdot q_1^{1}\cdot q_2^{1}\cdot q_3^{6}-120\cdot q_1^{1}\cdot q_2^{3}\cdot q_3^{4}-120\cdot q_1^{1}\cdot q_2^{5}\cdot q_3^{2}\\
&\quad +240\cdot q_1^{3}\cdot q_2^{1}\cdot q_3^{4}+480\cdot q_1^{3}\cdot q_2^{3}\cdot q_3^{2}-240\cdot q_1^{5}\cdot q_2^{1}\cdot q_3^{2}\\
&\quad +480\cdot q_0^{1}\cdot q_1^{2}\cdot q_2^{2}\cdot q_3^{3}+240\cdot q_0^{1}\cdot q_1^{2}\cdot q_2^{4}\cdot q_3^{1}\\
&\quad -240\cdot q_0^{1}\cdot q_1^{4}\cdot q_2^{2}\cdot q_3^{1}+240\cdot q_0^{2}\cdot q_1^{1}\cdot q_2^{1}\cdot q_3^{4}\\
&\quad +480\cdot q_0^{2}\cdot q_1^{1}\cdot q_2^{3}\cdot q_3^{2}-480\cdot q_0^{2}\cdot q_1^{3}\cdot q_2^{1}\cdot q_3^{2}\\
&\quad -480\cdot q_0^{3}\cdot q_1^{2}\cdot q_2^{2}\cdot q_3^{1}-240\cdot q_0^{4}\cdot q_1^{1}\cdot q_2^{1}\cdot q_3^{2}\end{align*}
\begin{align*}(T_{h}^4)_{9,1}&=2\cdot q_0^{1}\cdot q_1^{3}\cdot q_2^{4}+2\cdot q_0^{1}\cdot q_1^{3}\cdot q_3^{4}-2\cdot q_0^{3}\cdot q_1^{1}\cdot q_2^{4}\\
&\quad -2\cdot q_0^{3}\cdot q_1^{1}\cdot q_3^{4}-2\cdot q_0^{4}\cdot q_2^{1}\cdot q_3^{3}+2\cdot q_0^{4}\cdot q_2^{3}\cdot q_3^{1}\\
&\quad -2\cdot q_1^{4}\cdot q_2^{1}\cdot q_3^{3}+2\cdot q_1^{4}\cdot q_2^{3}\cdot q_3^{1}-12\cdot q_0^{1}\cdot q_1^{3}\cdot q_2^{2}\cdot q_3^{2}\\
&\quad +12\cdot q_0^{2}\cdot q_1^{2}\cdot q_2^{1}\cdot q_3^{3}-12\cdot q_0^{2}\cdot q_1^{2}\cdot q_2^{3}\cdot q_3^{1}\\
&\quad +12\cdot q_0^{3}\cdot q_1^{1}\cdot q_2^{2}\cdot q_3^{2}\end{align*}
\begin{align*}(T_{h}^4)_{9,2}&=2\cdot q_0^{4}\cdot q_2^{4}+2\cdot q_0^{4}\cdot q_3^{4}+2\cdot q_1^{4}\cdot q_2^{4}+2\cdot q_1^{4}\cdot q_3^{4}\\
&\quad -12\cdot q_0^{2}\cdot q_1^{2}\cdot q_2^{4}-12\cdot q_0^{2}\cdot q_1^{2}\cdot q_3^{4}-12\cdot q_0^{4}\cdot q_2^{2}\cdot q_3^{2}\\
&\quad -12\cdot q_1^{4}\cdot q_2^{2}\cdot q_3^{2}+32\cdot q_0^{1}\cdot q_1^{3}\cdot q_2^{1}\cdot q_3^{3}-32\cdot q_0^{1}\cdot q_1^{3}\cdot q_2^{3}\cdot q_3^{1}\\
&\quad +72\cdot q_0^{2}\cdot q_1^{2}\cdot q_2^{2}\cdot q_3^{2}-32\cdot q_0^{3}\cdot q_1^{1}\cdot q_2^{1}\cdot q_3^{3}\\
&\quad +32\cdot q_0^{3}\cdot q_1^{1}\cdot q_2^{3}\cdot q_3^{1}\end{align*}
\begin{align*}(T_{h}^4)_{9,3}&=1\cdot q_0^{3}\cdot q_2^{5}-1\cdot q_0^{5}\cdot q_2^{3}+1\cdot q_1^{3}\cdot q_3^{5}-1\cdot q_1^{5}\cdot q_3^{3}\\
&\quad -3\cdot q_0^{1}\cdot q_1^{2}\cdot q_2^{5}+3\cdot q_0^{1}\cdot q_1^{4}\cdot q_2^{3}-3\cdot q_0^{2}\cdot q_1^{1}\cdot q_3^{5}\\
&\quad +2\cdot q_0^{2}\cdot q_1^{3}\cdot q_3^{3}+2\cdot q_0^{3}\cdot q_1^{2}\cdot q_2^{3}-3\cdot q_0^{3}\cdot q_2^{1}\cdot q_3^{4}\\
&\quad -2\cdot q_0^{3}\cdot q_2^{3}\cdot q_3^{2}+3\cdot q_0^{4}\cdot q_1^{1}\cdot q_3^{3}+3\cdot q_0^{5}\cdot q_2^{1}\cdot q_3^{2}\\
&\quad -2\cdot q_1^{3}\cdot q_2^{2}\cdot q_3^{3}-3\cdot q_1^{3}\cdot q_2^{4}\cdot q_3^{1}+3\cdot q_1^{5}\cdot q_2^{2}\cdot q_3^{1}\\
&\quad +9\cdot q_0^{1}\cdot q_1^{2}\cdot q_2^{1}\cdot q_3^{4}+6\cdot q_0^{1}\cdot q_1^{2}\cdot q_2^{3}\cdot q_3^{2}\\
&\quad -9\cdot q_0^{1}\cdot q_1^{4}\cdot q_2^{1}\cdot q_3^{2}+6\cdot q_0^{2}\cdot q_1^{1}\cdot q_2^{2}\cdot q_3^{3}\\
&\quad +9\cdot q_0^{2}\cdot q_1^{1}\cdot q_2^{4}\cdot q_3^{1}-6\cdot q_0^{2}\cdot q_1^{3}\cdot q_2^{2}\cdot q_3^{1}\\
&\quad -6\cdot q_0^{3}\cdot q_1^{2}\cdot q_2^{1}\cdot q_3^{2}-9\cdot q_0^{4}\cdot q_1^{1}\cdot q_2^{2}\cdot q_3^{1}\end{align*}
\begin{align*}(T_{h}^4)_{9,4}&=1\cdot q_0^{3}\cdot q_3^{5}-1\cdot q_0^{5}\cdot q_3^{3}-1\cdot q_1^{3}\cdot q_2^{5}+1\cdot q_1^{5}\cdot q_2^{3}\\
&\quad -3\cdot q_0^{1}\cdot q_1^{2}\cdot q_3^{5}+3\cdot q_0^{1}\cdot q_1^{4}\cdot q_3^{3}+3\cdot q_0^{2}\cdot q_1^{1}\cdot q_2^{5}\\
&\quad -2\cdot q_0^{2}\cdot q_1^{3}\cdot q_2^{3}+2\cdot q_0^{3}\cdot q_1^{2}\cdot q_3^{3}-2\cdot q_0^{3}\cdot q_2^{2}\cdot q_3^{3}\\
&\quad -3\cdot q_0^{3}\cdot q_2^{4}\cdot q_3^{1}-3\cdot q_0^{4}\cdot q_1^{1}\cdot q_2^{3}+3\cdot q_0^{5}\cdot q_2^{2}\cdot q_3^{1}\\
&\quad +3\cdot q_1^{3}\cdot q_2^{1}\cdot q_3^{4}+2\cdot q_1^{3}\cdot q_2^{3}\cdot q_3^{2}-3\cdot q_1^{5}\cdot q_2^{1}\cdot q_3^{2}\\
&\quad +6\cdot q_0^{1}\cdot q_1^{2}\cdot q_2^{2}\cdot q_3^{3}+9\cdot q_0^{1}\cdot q_1^{2}\cdot q_2^{4}\cdot q_3^{1}\\
&\quad -9\cdot q_0^{1}\cdot q_1^{4}\cdot q_2^{2}\cdot q_3^{1}-9\cdot q_0^{2}\cdot q_1^{1}\cdot q_2^{1}\cdot q_3^{4}\\
&\quad -6\cdot q_0^{2}\cdot q_1^{1}\cdot q_2^{3}\cdot q_3^{2}+6\cdot q_0^{2}\cdot q_1^{3}\cdot q_2^{1}\cdot q_3^{2}\\
&\quad -6\cdot q_0^{3}\cdot q_1^{2}\cdot q_2^{2}\cdot q_3^{1}+9\cdot q_0^{4}\cdot q_1^{1}\cdot q_2^{1}\cdot q_3^{2}\end{align*}
\begin{align*}(T_{h}^4)_{9,5}&=3\cdot q_0^{1}\cdot q_1^{1}\cdot q_2^{6}-3\cdot q_0^{1}\cdot q_1^{1}\cdot q_3^{6}-8\cdot q_0^{1}\cdot q_1^{3}\cdot q_2^{4}\\
&\quad +8\cdot q_0^{1}\cdot q_1^{3}\cdot q_3^{4}+3\cdot q_0^{1}\cdot q_1^{5}\cdot q_2^{2}-3\cdot q_0^{1}\cdot q_1^{5}\cdot q_3^{2}\\
&\quad -3\cdot q_0^{2}\cdot q_2^{1}\cdot q_3^{5}-6\cdot q_0^{2}\cdot q_2^{3}\cdot q_3^{3}-3\cdot q_0^{2}\cdot q_2^{5}\cdot q_3^{1}\\
&\quad -8\cdot q_0^{3}\cdot q_1^{1}\cdot q_2^{4}+8\cdot q_0^{3}\cdot q_1^{1}\cdot q_3^{4}+6\cdot q_0^{3}\cdot q_1^{3}\cdot q_2^{2}\\
&\quad -6\cdot q_0^{3}\cdot q_1^{3}\cdot q_3^{2}+8\cdot q_0^{4}\cdot q_2^{1}\cdot q_3^{3}+8\cdot q_0^{4}\cdot q_2^{3}\cdot q_3^{1}\\
&\quad +3\cdot q_0^{5}\cdot q_1^{1}\cdot q_2^{2}-3\cdot q_0^{5}\cdot q_1^{1}\cdot q_3^{2}-3\cdot q_0^{6}\cdot q_2^{1}\cdot q_3^{1}\\
&\quad +3\cdot q_1^{2}\cdot q_2^{1}\cdot q_3^{5}+6\cdot q_1^{2}\cdot q_2^{3}\cdot q_3^{3}+3\cdot q_1^{2}\cdot q_2^{5}\cdot q_3^{1}\\
&\quad -8\cdot q_1^{4}\cdot q_2^{1}\cdot q_3^{3}-8\cdot q_1^{4}\cdot q_2^{3}\cdot q_3^{1}+3\cdot q_1^{6}\cdot q_2^{1}\cdot q_3^{1}\\
&\quad -3\cdot q_0^{1}\cdot q_1^{1}\cdot q_2^{2}\cdot q_3^{4}+3\cdot q_0^{1}\cdot q_1^{1}\cdot q_2^{4}\cdot q_3^{2}\\
&\quad +3\cdot q_0^{2}\cdot q_1^{4}\cdot q_2^{1}\cdot q_3^{1}-3\cdot q_0^{4}\cdot q_1^{2}\cdot q_2^{1}\cdot q_3^{1}\end{align*}
\begin{align*}(T_{h}^4)_{9,6}&=-3\cdot q_0^{2}\cdot q_2^{6}+3\cdot q_0^{2}\cdot q_3^{6}+8\cdot q_0^{4}\cdot q_2^{4}-8\cdot q_0^{4}\cdot q_3^{4}\\
&\quad -3\cdot q_0^{6}\cdot q_2^{2}+3\cdot q_0^{6}\cdot q_3^{2}+3\cdot q_1^{2}\cdot q_2^{6}-3\cdot q_1^{2}\cdot q_3^{6}\\
&\quad -8\cdot q_1^{4}\cdot q_2^{4}+8\cdot q_1^{4}\cdot q_3^{4}+3\cdot q_1^{6}\cdot q_2^{2}-3\cdot q_1^{6}\cdot q_3^{2}\\
&\quad +3\cdot q_0^{2}\cdot q_1^{4}\cdot q_2^{2}-3\cdot q_0^{2}\cdot q_1^{4}\cdot q_3^{2}+3\cdot q_0^{2}\cdot q_2^{2}\cdot q_3^{4}\\
&\quad -3\cdot q_0^{2}\cdot q_2^{4}\cdot q_3^{2}-3\cdot q_0^{4}\cdot q_1^{2}\cdot q_2^{2}+3\cdot q_0^{4}\cdot q_1^{2}\cdot q_3^{2}\\
&\quad -3\cdot q_1^{2}\cdot q_2^{2}\cdot q_3^{4}+3\cdot q_1^{2}\cdot q_2^{4}\cdot q_3^{2}-12\cdot q_0^{1}\cdot q_1^{1}\cdot q_2^{1}\cdot q_3^{5}\\
&\quad -24\cdot q_0^{1}\cdot q_1^{1}\cdot q_2^{3}\cdot q_3^{3}-12\cdot q_0^{1}\cdot q_1^{1}\cdot q_2^{5}\cdot q_3^{1}\\
&\quad +32\cdot q_0^{1}\cdot q_1^{3}\cdot q_2^{1}\cdot q_3^{3}+32\cdot q_0^{1}\cdot q_1^{3}\cdot q_2^{3}\cdot q_3^{1}\\
&\quad -12\cdot q_0^{1}\cdot q_1^{5}\cdot q_2^{1}\cdot q_3^{1}+32\cdot q_0^{3}\cdot q_1^{1}\cdot q_2^{1}\cdot q_3^{3}\\
&\quad +32\cdot q_0^{3}\cdot q_1^{1}\cdot q_2^{3}\cdot q_3^{1}-24\cdot q_0^{3}\cdot q_1^{3}\cdot q_2^{1}\cdot q_3^{1}\\
&\quad -12\cdot q_0^{5}\cdot q_1^{1}\cdot q_2^{1}\cdot q_3^{1}\end{align*}
\begin{align*}(T_{h}^4)_{9,7}&=-1\cdot q_0^{1}\cdot q_2^{7}+6\cdot q_0^{3}\cdot q_2^{5}-6\cdot q_0^{5}\cdot q_2^{3}+1\cdot q_0^{7}\cdot q_2^{1}\\
&\quad -1\cdot q_1^{1}\cdot q_3^{7}+6\cdot q_1^{3}\cdot q_3^{5}-6\cdot q_1^{5}\cdot q_3^{3}+1\cdot q_1^{7}\cdot q_3^{1}\\
&\quad +6\cdot q_0^{1}\cdot q_1^{2}\cdot q_2^{5}-6\cdot q_0^{1}\cdot q_1^{4}\cdot q_2^{3}+1\cdot q_0^{1}\cdot q_1^{6}\cdot q_2^{1}\\
&\quad -1\cdot q_0^{1}\cdot q_2^{1}\cdot q_3^{6}-3\cdot q_0^{1}\cdot q_2^{3}\cdot q_3^{4}-3\cdot q_0^{1}\cdot q_2^{5}\cdot q_3^{2}\\
&\quad +6\cdot q_0^{2}\cdot q_1^{1}\cdot q_3^{5}-12\cdot q_0^{2}\cdot q_1^{3}\cdot q_3^{3}+3\cdot q_0^{2}\cdot q_1^{5}\cdot q_3^{1}\\
&\quad -12\cdot q_0^{3}\cdot q_1^{2}\cdot q_2^{3}+3\cdot q_0^{3}\cdot q_1^{4}\cdot q_2^{1}+6\cdot q_0^{3}\cdot q_2^{1}\cdot q_3^{4}\\
&\quad +12\cdot q_0^{3}\cdot q_2^{3}\cdot q_3^{2}-6\cdot q_0^{4}\cdot q_1^{1}\cdot q_3^{3}+3\cdot q_0^{4}\cdot q_1^{3}\cdot q_3^{1}\\
&\quad +3\cdot q_0^{5}\cdot q_1^{2}\cdot q_2^{1}-6\cdot q_0^{5}\cdot q_2^{1}\cdot q_3^{2}+1\cdot q_0^{6}\cdot q_1^{1}\cdot q_3^{1}\\
&\quad -3\cdot q_1^{1}\cdot q_2^{2}\cdot q_3^{5}-3\cdot q_1^{1}\cdot q_2^{4}\cdot q_3^{3}-1\cdot q_1^{1}\cdot q_2^{6}\cdot q_3^{1}\\
&\quad +12\cdot q_1^{3}\cdot q_2^{2}\cdot q_3^{3}+6\cdot q_1^{3}\cdot q_2^{4}\cdot q_3^{1}-6\cdot q_1^{5}\cdot q_2^{2}\cdot q_3^{1}\\
&\quad +6\cdot q_0^{1}\cdot q_1^{2}\cdot q_2^{1}\cdot q_3^{4}+12\cdot q_0^{1}\cdot q_1^{2}\cdot q_2^{3}\cdot q_3^{2}\\
&\quad -6\cdot q_0^{1}\cdot q_1^{4}\cdot q_2^{1}\cdot q_3^{2}+12\cdot q_0^{2}\cdot q_1^{1}\cdot q_2^{2}\cdot q_3^{3}\\
&\quad +6\cdot q_0^{2}\cdot q_1^{1}\cdot q_2^{4}\cdot q_3^{1}-12\cdot q_0^{2}\cdot q_1^{3}\cdot q_2^{2}\cdot q_3^{1}\\
&\quad -12\cdot q_0^{3}\cdot q_1^{2}\cdot q_2^{1}\cdot q_3^{2}-6\cdot q_0^{4}\cdot q_1^{1}\cdot q_2^{2}\cdot q_3^{1}\end{align*}
\begin{align*}(T_{h}^4)_{9,8}&=1\cdot q_0^{1}\cdot q_3^{7}-6\cdot q_0^{3}\cdot q_3^{5}+6\cdot q_0^{5}\cdot q_3^{3}-1\cdot q_0^{7}\cdot q_3^{1}\\
&\quad -1\cdot q_1^{1}\cdot q_2^{7}+6\cdot q_1^{3}\cdot q_2^{5}-6\cdot q_1^{5}\cdot q_2^{3}+1\cdot q_1^{7}\cdot q_2^{1}\\
&\quad -6\cdot q_0^{1}\cdot q_1^{2}\cdot q_3^{5}+6\cdot q_0^{1}\cdot q_1^{4}\cdot q_3^{3}-1\cdot q_0^{1}\cdot q_1^{6}\cdot q_3^{1}\\
&\quad +3\cdot q_0^{1}\cdot q_2^{2}\cdot q_3^{5}+3\cdot q_0^{1}\cdot q_2^{4}\cdot q_3^{3}+1\cdot q_0^{1}\cdot q_2^{6}\cdot q_3^{1}\\
&\quad +6\cdot q_0^{2}\cdot q_1^{1}\cdot q_2^{5}-12\cdot q_0^{2}\cdot q_1^{3}\cdot q_2^{3}+3\cdot q_0^{2}\cdot q_1^{5}\cdot q_2^{1}\\
&\quad +12\cdot q_0^{3}\cdot q_1^{2}\cdot q_3^{3}-3\cdot q_0^{3}\cdot q_1^{4}\cdot q_3^{1}-12\cdot q_0^{3}\cdot q_2^{2}\cdot q_3^{3}\\
&\quad -6\cdot q_0^{3}\cdot q_2^{4}\cdot q_3^{1}-6\cdot q_0^{4}\cdot q_1^{1}\cdot q_2^{3}+3\cdot q_0^{4}\cdot q_1^{3}\cdot q_2^{1}\\
&\quad -3\cdot q_0^{5}\cdot q_1^{2}\cdot q_3^{1}+6\cdot q_0^{5}\cdot q_2^{2}\cdot q_3^{1}+1\cdot q_0^{6}\cdot q_1^{1}\cdot q_2^{1}\\
&\quad -1\cdot q_1^{1}\cdot q_2^{1}\cdot q_3^{6}-3\cdot q_1^{1}\cdot q_2^{3}\cdot q_3^{4}-3\cdot q_1^{1}\cdot q_2^{5}\cdot q_3^{2}\\
&\quad +6\cdot q_1^{3}\cdot q_2^{1}\cdot q_3^{4}+12\cdot q_1^{3}\cdot q_2^{3}\cdot q_3^{2}-6\cdot q_1^{5}\cdot q_2^{1}\cdot q_3^{2}\\
&\quad -12\cdot q_0^{1}\cdot q_1^{2}\cdot q_2^{2}\cdot q_3^{3}-6\cdot q_0^{1}\cdot q_1^{2}\cdot q_2^{4}\cdot q_3^{1}\\
&\quad +6\cdot q_0^{1}\cdot q_1^{4}\cdot q_2^{2}\cdot q_3^{1}+6\cdot q_0^{2}\cdot q_1^{1}\cdot q_2^{1}\cdot q_3^{4}\\
&\quad +12\cdot q_0^{2}\cdot q_1^{1}\cdot q_2^{3}\cdot q_3^{2}-12\cdot q_0^{2}\cdot q_1^{3}\cdot q_2^{1}\cdot q_3^{2}\\
&\quad +12\cdot q_0^{3}\cdot q_1^{2}\cdot q_2^{2}\cdot q_3^{1}-6\cdot q_0^{4}\cdot q_1^{1}\cdot q_2^{1}\cdot q_3^{2}\end{align*}
\begin{align*}(T_{h}^4)_{9,9}&=1\cdot q_0^{8}+1\cdot q_1^{8}+1\cdot q_2^{8}+1\cdot q_3^{8}+4\cdot q_0^{2}\cdot q_1^{6}-16\cdot q_0^{2}\cdot q_2^{6}\\
&\quad -16\cdot q_0^{2}\cdot q_3^{6}+6\cdot q_0^{4}\cdot q_1^{4}+36\cdot q_0^{4}\cdot q_2^{4}+36\cdot q_0^{4}\cdot q_3^{4}\\
&\quad +4\cdot q_0^{6}\cdot q_1^{2}-16\cdot q_0^{6}\cdot q_2^{2}-16\cdot q_0^{6}\cdot q_3^{2}-16\cdot q_1^{2}\cdot q_2^{6}\\
&\quad -16\cdot q_1^{2}\cdot q_3^{6}+36\cdot q_1^{4}\cdot q_2^{4}+36\cdot q_1^{4}\cdot q_3^{4}-16\cdot q_1^{6}\cdot q_2^{2}\\
&\quad -16\cdot q_1^{6}\cdot q_3^{2}+4\cdot q_2^{2}\cdot q_3^{6}+6\cdot q_2^{4}\cdot q_3^{4}+4\cdot q_2^{6}\cdot q_3^{2}\\
&\quad +72\cdot q_0^{2}\cdot q_1^{2}\cdot q_2^{4}+72\cdot q_0^{2}\cdot q_1^{2}\cdot q_3^{4}-48\cdot q_0^{2}\cdot q_1^{4}\cdot q_2^{2}\\
&\quad -48\cdot q_0^{2}\cdot q_1^{4}\cdot q_3^{2}-48\cdot q_0^{2}\cdot q_2^{2}\cdot q_3^{4}-48\cdot q_0^{2}\cdot q_2^{4}\cdot q_3^{2}\\
&\quad -48\cdot q_0^{4}\cdot q_1^{2}\cdot q_2^{2}-48\cdot q_0^{4}\cdot q_1^{2}\cdot q_3^{2}+72\cdot q_0^{4}\cdot q_2^{2}\cdot q_3^{2}\\
&\quad -48\cdot q_1^{2}\cdot q_2^{2}\cdot q_3^{4}-48\cdot q_1^{2}\cdot q_2^{4}\cdot q_3^{2}+72\cdot q_1^{4}\cdot q_2^{2}\cdot q_3^{2}\\
&\quad +144\cdot q_0^{2}\cdot q_1^{2}\cdot q_2^{2}\cdot q_3^{2}\end{align*}
\section{Polynomial $U_h$}
\begin{align*}U_h(q_0,q_1,q_2,q_3)&=-4\cdot q_0^{4}+12\cdot q_1^{4}+12\cdot q_2^{4}-4\cdot q_3^{4}+8\cdot q_0^{2}\cdot q_1^{2}-8\cdot q_0^{2}\cdot q_2^{2}\\
&\quad -8\cdot q_0^{2}\cdot q_3^{2}+24\cdot q_1^{2}\cdot q_2^{2}-8\cdot q_1^{2}\cdot q_3^{2}+8\cdot q_2^{2}\cdot q_3^{2}\\
&\quad -32\cdot q_0^{1}\cdot q_1^{1}\cdot q_2^{1}\cdot q_3^{1}\end{align*}
\section{Polynomial $Q_h$}
\begin{align*}Q_h(q_0,q_1,q_2,q_3)&=-16\cdot q_0^{2}\cdot q_2^{6}-16\cdot q_0^{2}\cdot q_3^{6}-32\cdot q_0^{4}\cdot q_2^{4}-32\cdot q_0^{4}\cdot q_3^{4}\\
&\quad -16\cdot q_0^{6}\cdot q_2^{2}-16\cdot q_0^{6}\cdot q_3^{2}+48\cdot q_1^{2}\cdot q_2^{6}-16\cdot q_1^{2}\cdot q_3^{6}\\
&\quad +96\cdot q_1^{4}\cdot q_2^{4}-32\cdot q_1^{4}\cdot q_3^{4}+48\cdot q_1^{6}\cdot q_2^{2}-16\cdot q_1^{6}\cdot q_3^{2}\\
&\quad +64\cdot q_0^{2}\cdot q_1^{2}\cdot q_2^{4}-64\cdot q_0^{2}\cdot q_1^{2}\cdot q_3^{4}+80\cdot q_0^{2}\cdot q_1^{4}\cdot q_2^{2}\\
&\quad -48\cdot q_0^{2}\cdot q_1^{4}\cdot q_3^{2}-48\cdot q_0^{2}\cdot q_2^{2}\cdot q_3^{4}-48\cdot q_0^{2}\cdot q_2^{4}\cdot q_3^{2}\\
&\quad +16\cdot q_0^{4}\cdot q_1^{2}\cdot q_2^{2}-48\cdot q_0^{4}\cdot q_1^{2}\cdot q_3^{2}-64\cdot q_0^{4}\cdot q_2^{2}\cdot q_3^{2}\\
&\quad +16\cdot q_1^{2}\cdot q_2^{2}\cdot q_3^{4}+80\cdot q_1^{2}\cdot q_2^{4}\cdot q_3^{2}+64\cdot q_1^{4}\cdot q_2^{2}\cdot q_3^{2}\end{align*}
\section{Polynomial $G$}
We store $20G$ instead of $G$ in the memory because $20G$ can have all of its coefficients stored with full precision, without rounding. Recall that computers work in binary and for them, $\frac15$ is an infinite expression that must be rounded. Computers can only tolerate denominators that are powers of two.
\begin{align*}20G(q_0,q_1,q_2,q_3)&=-\frac{134541}{2^{10}}\cdot q_0^{8}+\frac{162281}{2^{16}}\cdot q_1^{8}+\frac{648739}{2^{18}}\cdot q_2^{8}\\
&\quad -\frac{17158373}{2^{17}}\cdot q_3^{8}+\frac{9649177}{2^{19}}\cdot q_0^{2}\cdot q_1^{6}+\frac{14823259}{2^{20}}\cdot q_0^{2}\cdot q_2^{6}\\
&\quad -\frac{324641461}{2^{20}}\cdot q_0^{2}\cdot q_3^{6}-\frac{21419991}{2^{18}}\cdot q_0^{4}\cdot q_1^{4}\\
&\quad -\frac{387529743}{2^{21}}\cdot q_0^{4}\cdot q_2^{4}-\frac{1069566183}{2^{21}}\cdot q_0^{4}\cdot q_3^{4}\\
&\quad +\frac{223173297}{2^{19}}\cdot q_0^{6}\cdot q_1^{2}-\frac{256000221}{2^{20}}\cdot q_0^{6}\cdot q_2^{2}\\
&\quad -\frac{322962261}{2^{20}}\cdot q_0^{6}\cdot q_3^{2}+\frac{44438839}{2^{20}}\cdot q_1^{2}\cdot q_2^{6}-\frac{258181681}{2^{20}}\cdot q_1^{2}\cdot q_3^{6}\\
&\quad -\frac{146332863}{2^{21}}\cdot q_1^{4}\cdot q_2^{4}-\frac{383178903}{2^{21}}\cdot q_1^{4}\cdot q_3^{4}\\
&\quad +\frac{44434919}{2^{20}}\cdot q_1^{6}\cdot q_2^{2}+\frac{14317079}{2^{20}}\cdot q_1^{6}\cdot q_3^{2}+\frac{222092227}{2^{19}}\cdot q_2^{2}\cdot q_3^{6}\\
&\quad -\frac{326439}{2^{12}}\cdot q_2^{4}\cdot q_3^{4}+\frac{9415967}{2^{19}}\cdot q_2^{6}\cdot q_3^{2}+\frac{430091247}{2^{20}}\cdot q_0^{2}\cdot q_1^{2}\cdot q_2^{4}\\
&\quad -\frac{449022393}{2^{20}}\cdot q_0^{2}\cdot q_1^{2}\cdot q_3^{4}-\frac{446437383}{2^{20}}\cdot q_0^{2}\cdot q_1^{4}\cdot q_2^{2}\\
&\quad +\frac{134241297}{2^{20}}\cdot q_0^{2}\cdot q_1^{4}\cdot q_3^{2}-\frac{627763863}{2^{20}}\cdot q_0^{2}\cdot q_2^{2}\cdot q_3^{4}\\
&\quad +\frac{134173257}{2^{20}}\cdot q_0^{2}\cdot q_2^{4}\cdot q_3^{2}+\frac{1077571077}{2^{20}}\cdot q_0^{4}\cdot q_1^{2}\cdot q_2^{2}\\
&\quad -\frac{633692043}{2^{20}}\cdot q_0^{4}\cdot q_1^{2}\cdot q_3^{2}-\frac{454875813}{2^{20}}\cdot q_0^{4}\cdot q_2^{2}\cdot q_3^{2}\\
&\quad +\frac{1081992477}{2^{20}}\cdot q_1^{2}\cdot q_2^{2}\cdot q_3^{4}-\frac{448012203}{2^{20}}\cdot q_1^{2}\cdot q_2^{4}\cdot q_3^{2}\\
&\quad +\frac{428588667}{2^{20}}\cdot q_1^{4}\cdot q_2^{2}\cdot q_3^{2}+\frac{795393177}{2^{19}}\cdot q_0^{2}\cdot q_1^{2}\cdot q_2^{2}\cdot q_3^{2}\end{align*}
\section{Substitution $q_0=1$}
\subsection{Polynomial $U_h(x,y,z)$}
\begin{align*}U_h(x,y,z)&=-4+8\cdot x^{2}+12\cdot x^{4}-8\cdot y^{2}+12\cdot y^{4}-8\cdot z^{2}-4\cdot z^{4}+24\cdot x^{2}\cdot y^{2}\\
&\quad -8\cdot x^{2}\cdot z^{2}+8\cdot y^{2}\cdot z^{2}-32\cdot x^{1}\cdot y^{1}\cdot z^{1}\end{align*}\subsection{Polynomial $Q_h(x,y,z)$}
\begin{align*}Q_h(x,y,z)&=-16\cdot y^{2}-32\cdot y^{4}-16\cdot y^{6}-16\cdot z^{2}-32\cdot z^{4}-16\cdot z^{6}+16\cdot x^{2}\cdot y^{2}\\
&\quad +64\cdot x^{2}\cdot y^{4}+48\cdot x^{2}\cdot y^{6}-48\cdot x^{2}\cdot z^{2}-64\cdot x^{2}\cdot z^{4}-16\cdot x^{2}\cdot z^{6}\\
&\quad +80\cdot x^{4}\cdot y^{2}+96\cdot x^{4}\cdot y^{4}-48\cdot x^{4}\cdot z^{2}-32\cdot x^{4}\cdot z^{4}+48\cdot x^{6}\cdot y^{2}\\
&\quad -16\cdot x^{6}\cdot z^{2}-64\cdot y^{2}\cdot z^{2}-48\cdot y^{2}\cdot z^{4}-48\cdot y^{4}\cdot z^{2}+16\cdot x^{2}\cdot y^{2}\cdot z^{4}\\
&\quad +80\cdot x^{2}\cdot y^{4}\cdot z^{2}+64\cdot x^{4}\cdot y^{2}\cdot z^{2}\end{align*}\subsection{Polynomial $20G(x,y,z)$}
\begin{align*}20G(x,y,z)&=-\frac{134541}{2^{10}}+\frac{223173297}{2^{19}}\cdot x^{2}-\frac{21419991}{2^{18}}\cdot x^{4}+\frac{9649177}{2^{19}}\cdot x^{6}\\
&\quad +\frac{162281}{2^{16}}\cdot x^{8}-\frac{256000221}{2^{20}}\cdot y^{2}-\frac{387529743}{2^{21}}\cdot y^{4}\\
&\quad +\frac{14823259}{2^{20}}\cdot y^{6}+\frac{648739}{2^{18}}\cdot y^{8}-\frac{322962261}{2^{20}}\cdot z^{2}\\
&\quad -\frac{1069566183}{2^{21}}\cdot z^{4}-\frac{324641461}{2^{20}}\cdot z^{6}-\frac{17158373}{2^{17}}\cdot z^{8}\\
&\quad +\frac{1077571077}{2^{20}}\cdot x^{2}\cdot y^{2}+\frac{430091247}{2^{20}}\cdot x^{2}\cdot y^{4}+\frac{44438839}{2^{20}}\cdot x^{2}\cdot y^{6}\\
&\quad -\frac{633692043}{2^{20}}\cdot x^{2}\cdot z^{2}-\frac{449022393}{2^{20}}\cdot x^{2}\cdot z^{4}-\frac{258181681}{2^{20}}\cdot x^{2}\cdot z^{6}\\
&\quad -\frac{446437383}{2^{20}}\cdot x^{4}\cdot y^{2}-\frac{146332863}{2^{21}}\cdot x^{4}\cdot y^{4}+\frac{134241297}{2^{20}}\cdot x^{4}\cdot z^{2}\\
&\quad -\frac{383178903}{2^{21}}\cdot x^{4}\cdot z^{4}+\frac{44434919}{2^{20}}\cdot x^{6}\cdot y^{2}+\frac{14317079}{2^{20}}\cdot x^{6}\cdot z^{2}\\
&\quad -\frac{454875813}{2^{20}}\cdot y^{2}\cdot z^{2}-\frac{627763863}{2^{20}}\cdot y^{2}\cdot z^{4}+\frac{222092227}{2^{19}}\cdot y^{2}\cdot z^{6}\\
&\quad +\frac{134173257}{2^{20}}\cdot y^{4}\cdot z^{2}-\frac{326439}{2^{12}}\cdot y^{4}\cdot z^{4}+\frac{9415967}{2^{19}}\cdot y^{6}\cdot z^{2}\\
&\quad +\frac{795393177}{2^{19}}\cdot x^{2}\cdot y^{2}\cdot z^{2}+\frac{1081992477}{2^{20}}\cdot x^{2}\cdot y^{2}\cdot z^{4}\\
&\quad -\frac{448012203}{2^{20}}\cdot x^{2}\cdot y^{4}\cdot z^{2}+\frac{428588667}{2^{20}}\cdot x^{4}\cdot y^{2}\cdot z^{2}\end{align*}\section{Substitution $q_1=1$}
\subsection{Polynomial $U_h(x,y,z)$}
\begin{align*}U_h(x,y,z)&=12+8\cdot x^{2}-4\cdot x^{4}+24\cdot y^{2}+12\cdot y^{4}-8\cdot z^{2}-4\cdot z^{4}-8\cdot x^{2}\cdot y^{2}\\
&\quad -8\cdot x^{2}\cdot z^{2}+8\cdot y^{2}\cdot z^{2}-32\cdot x^{1}\cdot y^{1}\cdot z^{1}\end{align*}\subsection{Polynomial $Q_h(x,y,z)$}
\begin{align*}Q_h(x,y,z)&=48\cdot y^{2}+96\cdot y^{4}+48\cdot y^{6}-16\cdot z^{2}-32\cdot z^{4}-16\cdot z^{6}+80\cdot x^{2}\cdot y^{2}\\
&\quad +64\cdot x^{2}\cdot y^{4}-16\cdot x^{2}\cdot y^{6}-48\cdot x^{2}\cdot z^{2}-64\cdot x^{2}\cdot z^{4}-16\cdot x^{2}\cdot z^{6}\\
&\quad +16\cdot x^{4}\cdot y^{2}-32\cdot x^{4}\cdot y^{4}-48\cdot x^{4}\cdot z^{2}-32\cdot x^{4}\cdot z^{4}-16\cdot x^{6}\cdot y^{2}\\
&\quad -16\cdot x^{6}\cdot z^{2}+64\cdot y^{2}\cdot z^{2}+16\cdot y^{2}\cdot z^{4}+80\cdot y^{4}\cdot z^{2}-48\cdot x^{2}\cdot y^{2}\cdot z^{4}\\
&\quad -48\cdot x^{2}\cdot y^{4}\cdot z^{2}-64\cdot x^{4}\cdot y^{2}\cdot z^{2}\end{align*}\subsection{Polynomial $20G(x,y,z)$}
\begin{align*}20G(x,y,z)&=\frac{162281}{2^{16}}+\frac{9649177}{2^{19}}\cdot x^{2}-\frac{21419991}{2^{18}}\cdot x^{4}+\frac{223173297}{2^{19}}\cdot x^{6}\\
&\quad -\frac{134541}{2^{10}}\cdot x^{8}+\frac{44434919}{2^{20}}\cdot y^{2}-\frac{146332863}{2^{21}}\cdot y^{4}\\
&\quad +\frac{44438839}{2^{20}}\cdot y^{6}+\frac{648739}{2^{18}}\cdot y^{8}+\frac{14317079}{2^{20}}\cdot z^{2}-\frac{383178903}{2^{21}}\cdot z^{4}\\
&\quad -\frac{258181681}{2^{20}}\cdot z^{6}-\frac{17158373}{2^{17}}\cdot z^{8}-\frac{446437383}{2^{20}}\cdot x^{2}\cdot y^{2}\\
&\quad +\frac{430091247}{2^{20}}\cdot x^{2}\cdot y^{4}+\frac{14823259}{2^{20}}\cdot x^{2}\cdot y^{6}+\frac{134241297}{2^{20}}\cdot x^{2}\cdot z^{2}\\
&\quad -\frac{449022393}{2^{20}}\cdot x^{2}\cdot z^{4}-\frac{324641461}{2^{20}}\cdot x^{2}\cdot z^{6}+\frac{1077571077}{2^{20}}\cdot x^{4}\cdot y^{2}\\
&\quad -\frac{387529743}{2^{21}}\cdot x^{4}\cdot y^{4}-\frac{633692043}{2^{20}}\cdot x^{4}\cdot z^{2}-\frac{1069566183}{2^{21}}\cdot x^{4}\cdot z^{4}\\
&\quad -\frac{256000221}{2^{20}}\cdot x^{6}\cdot y^{2}-\frac{322962261}{2^{20}}\cdot x^{6}\cdot z^{2}+\frac{428588667}{2^{20}}\cdot y^{2}\cdot z^{2}\\
&\quad +\frac{1081992477}{2^{20}}\cdot y^{2}\cdot z^{4}+\frac{222092227}{2^{19}}\cdot y^{2}\cdot z^{6}-\frac{448012203}{2^{20}}\cdot y^{4}\cdot z^{2}\\
&\quad -\frac{326439}{2^{12}}\cdot y^{4}\cdot z^{4}+\frac{9415967}{2^{19}}\cdot y^{6}\cdot z^{2}+\frac{795393177}{2^{19}}\cdot x^{2}\cdot y^{2}\cdot z^{2}\\
&\quad -\frac{627763863}{2^{20}}\cdot x^{2}\cdot y^{2}\cdot z^{4}+\frac{134173257}{2^{20}}\cdot x^{2}\cdot y^{4}\cdot z^{2}\\
&\quad -\frac{454875813}{2^{20}}\cdot x^{4}\cdot y^{2}\cdot z^{2}\end{align*}\section{Substitution $q_2=1$}
\subsection{Polynomial $U_h(x,y,z)$}
\begin{align*}U_h(x,y,z)&=12-8\cdot x^{2}-4\cdot x^{4}+24\cdot y^{2}+12\cdot y^{4}+8\cdot z^{2}-4\cdot z^{4}+8\cdot x^{2}\cdot y^{2}\\
&\quad -8\cdot x^{2}\cdot z^{2}-8\cdot y^{2}\cdot z^{2}-32\cdot x^{1}\cdot y^{1}\cdot z^{1}\end{align*}\subsection{Polynomial $Q_h(x,y,z)$}
\begin{align*}Q_h(x,y,z)&=-16\cdot x^{2}-32\cdot x^{4}-16\cdot x^{6}+48\cdot y^{2}+96\cdot y^{4}+48\cdot y^{6}+64\cdot x^{2}\cdot y^{2}\\
&\quad +80\cdot x^{2}\cdot y^{4}-48\cdot x^{2}\cdot z^{2}-48\cdot x^{2}\cdot z^{4}-16\cdot x^{2}\cdot z^{6}+16\cdot x^{4}\cdot y^{2}\\
&\quad -64\cdot x^{4}\cdot z^{2}-32\cdot x^{4}\cdot z^{4}-16\cdot x^{6}\cdot z^{2}+80\cdot y^{2}\cdot z^{2}+16\cdot y^{2}\cdot z^{4}\\
&\quad -16\cdot y^{2}\cdot z^{6}+64\cdot y^{4}\cdot z^{2}-32\cdot y^{4}\cdot z^{4}-16\cdot y^{6}\cdot z^{2}-64\cdot x^{2}\cdot y^{2}\cdot z^{4}\\
&\quad -48\cdot x^{2}\cdot y^{4}\cdot z^{2}-48\cdot x^{4}\cdot y^{2}\cdot z^{2}\end{align*}\subsection{Polynomial $20G(x,y,z)$}
\begin{align*}20G(x,y,z)&=\frac{648739}{2^{18}}+\frac{14823259}{2^{20}}\cdot x^{2}-\frac{387529743}{2^{21}}\cdot x^{4}-\frac{256000221}{2^{20}}\cdot x^{6}\\
&\quad -\frac{134541}{2^{10}}\cdot x^{8}+\frac{44438839}{2^{20}}\cdot y^{2}-\frac{146332863}{2^{21}}\cdot y^{4}\\
&\quad +\frac{44434919}{2^{20}}\cdot y^{6}+\frac{162281}{2^{16}}\cdot y^{8}+\frac{9415967}{2^{19}}\cdot z^{2}-\frac{326439}{2^{12}}\cdot z^{4}\\
&\quad +\frac{222092227}{2^{19}}\cdot z^{6}-\frac{17158373}{2^{17}}\cdot z^{8}+\frac{430091247}{2^{20}}\cdot x^{2}\cdot y^{2}\\
&\quad -\frac{446437383}{2^{20}}\cdot x^{2}\cdot y^{4}+\frac{9649177}{2^{19}}\cdot x^{2}\cdot y^{6}+\frac{134173257}{2^{20}}\cdot x^{2}\cdot z^{2}\\
&\quad -\frac{627763863}{2^{20}}\cdot x^{2}\cdot z^{4}-\frac{324641461}{2^{20}}\cdot x^{2}\cdot z^{6}+\frac{1077571077}{2^{20}}\cdot x^{4}\cdot y^{2}\\
&\quad -\frac{21419991}{2^{18}}\cdot x^{4}\cdot y^{4}-\frac{454875813}{2^{20}}\cdot x^{4}\cdot z^{2}-\frac{1069566183}{2^{21}}\cdot x^{4}\cdot z^{4}\\
&\quad +\frac{223173297}{2^{19}}\cdot x^{6}\cdot y^{2}-\frac{322962261}{2^{20}}\cdot x^{6}\cdot z^{2}-\frac{448012203}{2^{20}}\cdot y^{2}\cdot z^{2}\\
&\quad +\frac{1081992477}{2^{20}}\cdot y^{2}\cdot z^{4}-\frac{258181681}{2^{20}}\cdot y^{2}\cdot z^{6}+\frac{428588667}{2^{20}}\cdot y^{4}\cdot z^{2}\\
&\quad -\frac{383178903}{2^{21}}\cdot y^{4}\cdot z^{4}+\frac{14317079}{2^{20}}\cdot y^{6}\cdot z^{2}+\frac{795393177}{2^{19}}\cdot x^{2}\cdot y^{2}\cdot z^{2}\\
&\quad -\frac{449022393}{2^{20}}\cdot x^{2}\cdot y^{2}\cdot z^{4}+\frac{134241297}{2^{20}}\cdot x^{2}\cdot y^{4}\cdot z^{2}\\
&\quad -\frac{633692043}{2^{20}}\cdot x^{4}\cdot y^{2}\cdot z^{2}\end{align*}\section{Substitution $q_3=1$}
\subsection{Polynomial $U_h(x,y,z)$}
\begin{align*}U_h(x,y,z)&=-4-8\cdot x^{2}-4\cdot x^{4}-8\cdot y^{2}+12\cdot y^{4}+8\cdot z^{2}+12\cdot z^{4}+8\cdot x^{2}\cdot y^{2}\\
&\quad -8\cdot x^{2}\cdot z^{2}+24\cdot y^{2}\cdot z^{2}-32\cdot x^{1}\cdot y^{1}\cdot z^{1}\end{align*}\subsection{Polynomial $Q_h(x,y,z)$}
\begin{align*}Q_h(x,y,z)&=-16\cdot x^{2}-32\cdot x^{4}-16\cdot x^{6}-16\cdot y^{2}-32\cdot y^{4}-16\cdot y^{6}-64\cdot x^{2}\cdot y^{2}\\
&\quad -48\cdot x^{2}\cdot y^{4}-48\cdot x^{2}\cdot z^{2}-48\cdot x^{2}\cdot z^{4}-16\cdot x^{2}\cdot z^{6}-48\cdot x^{4}\cdot y^{2}\\
&\quad -64\cdot x^{4}\cdot z^{2}-32\cdot x^{4}\cdot z^{4}-16\cdot x^{6}\cdot z^{2}+16\cdot y^{2}\cdot z^{2}+80\cdot y^{2}\cdot z^{4}\\
&\quad +48\cdot y^{2}\cdot z^{6}+64\cdot y^{4}\cdot z^{2}+96\cdot y^{4}\cdot z^{4}+48\cdot y^{6}\cdot z^{2}+64\cdot x^{2}\cdot y^{2}\cdot z^{4}\\
&\quad +80\cdot x^{2}\cdot y^{4}\cdot z^{2}+16\cdot x^{4}\cdot y^{2}\cdot z^{2}\end{align*}\subsection{Polynomial $20G(x,y,z)$}
\begin{align*}20G(x,y,z)&=-\frac{17158373}{2^{17}}-\frac{324641461}{2^{20}}\cdot x^{2}-\frac{1069566183}{2^{21}}\cdot x^{4}-\frac{322962261}{2^{20}}\cdot x^{6}\\
&\quad -\frac{134541}{2^{10}}\cdot x^{8}-\frac{258181681}{2^{20}}\cdot y^{2}-\frac{383178903}{2^{21}}\cdot y^{4}\\
&\quad +\frac{14317079}{2^{20}}\cdot y^{6}+\frac{162281}{2^{16}}\cdot y^{8}+\frac{222092227}{2^{19}}\cdot z^{2}\\
&\quad -\frac{326439}{2^{12}}\cdot z^{4}+\frac{9415967}{2^{19}}\cdot z^{6}+\frac{648739}{2^{18}}\cdot z^{8}-\frac{449022393}{2^{20}}\cdot x^{2}\cdot y^{2}\\
&\quad +\frac{134241297}{2^{20}}\cdot x^{2}\cdot y^{4}+\frac{9649177}{2^{19}}\cdot x^{2}\cdot y^{6}-\frac{627763863}{2^{20}}\cdot x^{2}\cdot z^{2}\\
&\quad +\frac{134173257}{2^{20}}\cdot x^{2}\cdot z^{4}+\frac{14823259}{2^{20}}\cdot x^{2}\cdot z^{6}-\frac{633692043}{2^{20}}\cdot x^{4}\cdot y^{2}\\
&\quad -\frac{21419991}{2^{18}}\cdot x^{4}\cdot y^{4}-\frac{454875813}{2^{20}}\cdot x^{4}\cdot z^{2}-\frac{387529743}{2^{21}}\cdot x^{4}\cdot z^{4}\\
&\quad +\frac{223173297}{2^{19}}\cdot x^{6}\cdot y^{2}-\frac{256000221}{2^{20}}\cdot x^{6}\cdot z^{2}+\frac{1081992477}{2^{20}}\cdot y^{2}\cdot z^{2}\\
&\quad -\frac{448012203}{2^{20}}\cdot y^{2}\cdot z^{4}+\frac{44438839}{2^{20}}\cdot y^{2}\cdot z^{6}+\frac{428588667}{2^{20}}\cdot y^{4}\cdot z^{2}\\
&\quad -\frac{146332863}{2^{21}}\cdot y^{4}\cdot z^{4}+\frac{44434919}{2^{20}}\cdot y^{6}\cdot z^{2}+\frac{795393177}{2^{19}}\cdot x^{2}\cdot y^{2}\cdot z^{2}\\
&\quad +\frac{430091247}{2^{20}}\cdot x^{2}\cdot y^{2}\cdot z^{4}-\frac{446437383}{2^{20}}\cdot x^{2}\cdot y^{4}\cdot z^{2}\\
&\quad +\frac{1077571077}{2^{20}}\cdot x^{4}\cdot y^{2}\cdot z^{2}\end{align*}

\end{document}
